\input{preamble}

% OK, start here.
%
\begin{document}

\title{Compléments de cohomologie étale}


\maketitle

\phantomsection
\label{section-phantom}

\tableofcontents


\section{Introduction}
\label{section-introduction}

\noindent
Ce chapitre est le deuxième d'une série de chapitres consacrés à la cohomologie
étale des schémas. Le premier chapitre se trouve dans
Cohomologie étale, section \ref{etale-cohomology-section-introduction}.

\medskip\noindent
La séparation avec le chapitre précédent est, en gros, la suivante : tout ce
qui concerne les « foncteurs image directe à support propre » (cohomologie à support propre
et son adjoint à droite), ainsi que leurs applications, figure dans ce chapitre.






\section{Prolongement des sections}
\label{section-growing}

\noindent
Dans cette section, nous étudions des résultats du type suivant.

\begin{lemma}
\label{lemma-section-support-in-locally-closed-pre}
Soit $X$ un schéma et soit $\mathcal{F}$ un faisceau abélien sur $X_\etale$.
Soit $\varphi : U' \to U$ un morphisme de $X_\etale$. Soit $Z' \subset U'$ un
sous-schéma fermé tel que $Z' \to U' \to U$ soit une immersion fermée
d'image $Z \subset U$. Il existe alors une bijection canonique
$$
\{s \in \mathcal{F}(U) \mid \text{Supp}(s) \subset Z\} =
\{s' \in \mathcal{F}(U') \mid \text{Supp}(s') \subset Z'\}
$$
qui est donnée par restriction si $\varphi^{-1}(Z) = Z'$.
\end{lemma}

\begin{proof}
Considérons le sous-schéma fermé $Z'' = \varphi^{-1}(Z)$ de $U'$.
Alors $Z' \subset Z''$ est fermé, puisque $Z'$ est fermé dans $U'$.
D'autre part, $Z' \to Z''$ est un morphisme étale
(comme morphisme entre schémas étales sur $Z$), et il est donc
ouvert. Ainsi $Z'' = Z' \amalg T$ pour un certain fermé $T$.
Le recouvrement ouvert $U' = (U' \setminus T) \cup (U' \setminus Z')$
montre que
$$
\{s' \in \mathcal{F}(U') \mid \text{Supp}(s') \subset Z'\} =
\{s' \in \mathcal{F}(U' \setminus T) \mid \text{Supp}(s') \subset Z'\}
$$
et le recouvrement étale $\{U' \setminus T \to U, U \setminus Z \to U\}$
montre que
$$
\{s \in \mathcal{F}(U) \mid \text{Supp}(s) \subset Z\} =
\{s' \in \mathcal{F}(U' \setminus T) \mid \text{Supp}(s') \subset Z'\}
$$
Ceci achève la démonstration.
\end{proof}

\begin{lemma}
\label{lemma-section-support-in-locally-closed}
Soit $X$ un schéma et soit $Z \subset X$ un sous-schéma localement fermé.
Soit $\mathcal{F}$ un faisceau abélien sur $X_\etale$. Étant donnés des ouverts
$U, U' \subset X$ contenant $Z$ comme sous-schéma fermé,
il existe une bijection canonique
$$
\{s \in \mathcal{F}(U) \mid \text{Supp}(s) \subset Z\} =
\{s \in \mathcal{F}(U') \mid \text{Supp}(s) \subset Z\}
$$
qui est donnée par restriction si $U' \subset U$.
\end{lemma}

\begin{proof}
Puisque $Z$ est un sous-schéma fermé de $U \cap U'$, il suffit de
démontrer le lemme lorsque $U' \subset U$. Il s'agit alors d'un cas particulier
du lemme \ref{lemma-section-support-in-locally-closed-pre}.
\end{proof}

\noindent
Introduisons une notation quelque peu inhabituelle qui nous sera
utile par la suite. Dans la situation du
lemme \ref{lemma-section-support-in-locally-closed} ci-dessus, posons
$$
H_Z(\mathcal{F}) = \{s \in \mathcal{F}(U) \mid \text{Supp}(s) \subset Z\}
$$
où $U \subset X$ est un sous-schéma ouvert quelconque contenant $Z$ comme
sous-schéma fermé. Si ce manque de précision gêne le lecteur, il peut
choisir $U = X \setminus \partial Z$, où
$\partial Z = \overline{Z}\setminus Z$ est la « frontière » de $Z$ dans $X$.
Cependant, dans plusieurs arguments ci-dessous, la possibilité de choisir
différents ouverts jouera un rôle. Voici quelques propriétés de cette
construction :
\begin{enumerate}
\item
\label{item-inclusion}
Si $Z \subset Z'$ sont des sous-schémas localement fermés de $X$ et si $Z$ est
fermé dans $Z'$, il existe une application injective naturelle
$$
H_Z(\mathcal{F}) \to H_{Z'}(\mathcal{F}).
$$
\item
\label{item-pullback}
Si $f : Y \to X$ est un morphisme de schémas et si $Z \subset X$ est un
sous-schéma localement fermé, il existe une application naturelle
d'image inverse $f^* : H_Z(\mathcal{F}) \to H_{f^{-1}Z}(f^{-1}\mathcal{F})$.
\end{enumerate}
Il sera commode d'étendre notre notation à la situation suivante :
supposons donnés $W \in X_\etale$ et un sous-schéma localement fermé
$Z \subset W$. Nous noterons alors
$$
H_Z(\mathcal{F}) =
\{s \in \mathcal{F}(U) \mid \text{Supp}(s) \subset Z\} =
H_Z(\mathcal{F}|_{W_\etale})
$$
où $U \subset W$ est un sous-schéma ouvert quelconque contenant $Z$
comme sous-schéma fermé, exactement comme ci-dessus\footnote{En fait, le
lemme \ref{lemma-section-support-in-locally-closed-pre}
montre que, si $Z$ est un schéma sur $X$ isomorphe à un sous-schéma
localement fermé d'un objet $W$ de $X_\etale$, alors
le choix de $W$ est sans importance.}.







\section{sections à support propre}
\label{section-compact-support}

\noindent
Une référence pour cette section est \cite[Exposé XVII, section 6]{SGA4}.
Soit $f : X \to Y$ un morphisme de schémas séparé et
localement de type fini. Dans cette section, nous définissons un foncteur
$f_! : \textit{Ab}(X_\etale) \to \textit{Ab}(Y_\etale)$
en définissant $f_!\mathcal{F} \subset f_*\mathcal{F}$ comme
le sous-faisceau des sections dont le support est propre sur $Y$
(au sens approprié).

\medskip\noindent
Avertissement : le foncteur $f_!$ est le faisceau de cohomologie de degré zéro
d'un foncteur $Rf_!$ sur la catégorie dérivée (insérer une référence future),
mais $Rf_!$ n'est pas le foncteur dérivé de $f_!$.

\begin{lemma}
\label{lemma-f-shriek-separated}
Soit $f : X \to Y$ un morphisme de schémas localement de type fini.
Soit $\mathcal{F}$ un faisceau abélien sur $X_\etale$. La règle
$$
Y_\etale \longrightarrow \textit{Ab},\quad
V \longmapsto \{s \in f_*\mathcal{F}(V) = \mathcal{F}(X_V) \mid
\text{Supp}(s) \subset X_V \text{ est propre sur }V\}
$$
définit un sous-faisceau abélien de $f_*\mathcal{F}$.
\end{lemma}

\noindent
Avertissement : ce faisceau n'est pas le « bon » si $f$ n'est pas séparé.

\begin{proof}
Rappelons que le support d'une section est fermé
(Cohomologie étale, lemme \ref{etale-cohomology-lemma-support-section-closed}) ;
les résultats de
Cohomologie des schémas, section \ref{coherent-section-proper-over-base},
s'appliquent donc. Le lemme ci-dessus et le lemme
\ref{coherent-lemma-union-closed-proper-over-base} de Cohomologie des schémas
montrent que notre sous-ensemble de $f_*\mathcal{F}(V)$ est un sous-groupe.
D'après le lemme de Cohomologie des schémas
\ref{coherent-lemma-base-change-closed-proper-over-base}
nous voyons que notre règle définit un sous-préfaisceau.
Enfin, supposons donnés $s \in f_*\mathcal{F}(V)$
et un recouvrement étale $\{V_i \to V\}$ tels que
le support de $s|_{V_i}$ soit propre sur $V_i$.
Le support de $s|_{V_i}$ est l'image réciproque du
support de $s|_V$ (utiliser la caractérisation du support
au moyen des fibres et le lemme
\ref{etale-cohomology-lemma-stalk-pullback} de Cohomologie étale).
Le support de $s$ est donc propre sur $V$ d'après le lemme
\ref{descent-lemma-descending-property-proper-over-base} de Descente.
Notre règle satisfait ainsi la condition de faisceau.
\end{proof}

\begin{lemma}
\label{lemma-separated-etale-shriek}
Soit $j : U \to X$ un morphisme étale séparé. Soit $\mathcal{F}$
un faisceau abélien sur $U_\etale$. L'image de l'application injective
$j_!\mathcal{F} \to j_*\mathcal{F}$
du lemme de Cohomologie étale
\ref{etale-cohomology-lemma-shriek-into-star-separated-etale}
est le sous-faisceau du lemme \ref{lemma-f-shriek-separated}.
\end{lemma}

\noindent
On pourrait aussi repousser ce lemme et le démontrer
à l'aide de la description des fibres des deux faisceaux.

\begin{proof}
La construction de $j_!\mathcal{F} \to j_*\mathcal{F}$ dans la démonstration du
lemme de Cohomologie étale
\ref{etale-cohomology-lemma-shriek-into-star-separated-etale}
s'obtient par la construction d'une application de préfaisceaux
$j_{p!}\mathcal{F} \to j_*\mathcal{F}$, dont l'image est manifestement
contenue dans le sous-faisceau du
lemme \ref{lemma-f-shriek-separated}.
Puisque $j_!\mathcal{F}$ est le faisceau associé à
$j_{p!}\mathcal{F}$, nous en déduisons que l'image de
$j_!\mathcal{F} \to j_*\mathcal{F}$ est contenue dans
ce sous-faisceau. Réciproquement, soit $s \in j_*\mathcal{F}(V)$
de support $Z$ propre sur $V$. Alors $Z \to V$ est
fini, d'image fermée $Z' \subset V$ ; voir le lemme
\ref{more-morphisms-lemma-characterize-finite} de Compléments sur les morphismes.
La restriction de $s$ à $V \setminus Z'$ est nulle, et la section nulle
appartient à l'image de $j_!\mathcal{F} \to j_*\mathcal{F}$.
D'autre part, si $v \in Z'$, nous pouvons trouver
un voisinage étale
$(V', v') \to (V, v)$ tel que l'on ait une décomposition
$U_{V'} = W \amalg U'_1 \amalg \ldots \amalg U'_n$
en sous-schémas ouverts et fermés, où $U'_i \to V'$ est un isomorphisme
et $T_{V'} \subset U'_1 \amalg \ldots \amalg U'_n$ ; voir le lemme
\ref{etale-lemma-etale-etale-local-technical} de Morphismes étales.
En inversant les isomorphismes $U'_i \to V'$,
nous obtenons $n$ morphismes $\varphi'_i : V' \to U$
et des sections $s'_i$ sur $V'$, obtenues par image inverse de $s$.
La section $\sum (\varphi'_i, s'_i)$ de
$j_{p!}\mathcal{F}$ sur $V'$ — voir la formule de $j_{p!}\mathcal{F}(V')$
dans la démonstration du lemme de Cohomologie étale
\ref{etale-cohomology-lemma-shriek-into-star-separated-etale},
a pour image la restriction de $s$ à $V'$ par construction.
Nous concluons que, localement pour la topologie étale, $s$ appartient à l'image
de $j_!\mathcal{F} \to j_*\mathcal{F}$, ce qui achève la démonstration.
\end{proof}

\begin{definition}
\label{definition-f-shriek-separated}
Soit $f : X \to Y$ un morphisme de schémas séparé (!) et
localement de type fini. Soit $\mathcal{F}$ un faisceau abélien sur
$X_\etale$. Le sous-faisceau $f_!\mathcal{F} \subset f_*\mathcal{F}$
construit au lemme \ref{lemma-f-shriek-separated} est appelé
{\it l'image directe à support propre}.
\end{definition}

\noindent
Le lemme \ref{lemma-separated-etale-shriek} montre que ceci ne contredit pas la
définition \ref{etale-cohomology-definition-extension-zero} de Cohomologie étale,
car les deux définitions coïncident lorsqu'elles s'appliquent. Vérifions-le.

\begin{lemma}
\label{lemma-proper-f-shriek}
Soit $f : X \to Y$ un morphisme propre de schémas.
Alors $f_! = f_*$.
\end{lemma}

\begin{proof}
Cela résulte immédiatement de la construction de $f_!$.
\end{proof}

\noindent
Voici une observation très utile.

\begin{remark}[Covariance par rapport aux immersions ouvertes]
\label{remark-covariance-f-shriek-separated}
Soit $f : X \to Y$ un morphisme de schémas séparé et
localement de type fini. Soit $\mathcal{F}$ un faisceau abélien sur $X_\etale$.
Soit $X' \subset X$ un sous-schéma ouvert. Notons $f' : X' \to Y$
la restriction de $f$. Il existe une application injective canonique
$$
f'_!(\mathcal{F}|_{X'}) \longrightarrow f_!\mathcal{F}
$$
En effet, soit $V \in Y_\etale$ et considérons une section
$s' \in f'_*(\mathcal{F}|_{X'})(V) = \mathcal{F}(X' \times_Y V)$
de support $Z'$ propre sur $V$. Alors $Z'$ est aussi fermé dans $X \times_Y V$ ;
voir le lemme de Cohomologie des schémas
\ref{coherent-lemma-functoriality-closed-proper-over-base}.
Il existe donc une unique section
$s \in \mathcal{F}(X \times_Y V) = f_*\mathcal{F}(V)$
dont la restriction à $X' \times_Y V$ est $s'$ et dont la restriction
à $X \times_Y V \setminus Z'$ est nulle ; voir le
lemme \ref{lemma-section-support-in-locally-closed}. Cette construction est
compatible avec les applications de restriction et induit donc l'application voulue
de faisceaux $f'_!(\mathcal{F}|_{X'}) \to f_!\mathcal{F}$, qui est manifestement
injective. Par construction, nous obtenons un diagramme commutatif
$$
\xymatrix{
f'_!(\mathcal{F}|_{X'}) \ar[r] \ar[d] &
f_!\mathcal{F} \ar[d] \\
f'_*(\mathcal{F}|_{X'}) &
f_*\mathcal{F} \ar[l]
}
$$
fonctoriel en $\mathcal{F}$. Il est clair que, pour $X'' \subset X'$ ouvert
et $f'' = f|_{X''} : X'' \to Y$, la composée des applications canoniques
$f''_!\mathcal{F}|_{X''} \to f'_!\mathcal{F}|_{X'} \to f_!\mathcal{F}$
que nous venons de construire est l'application canonique
$f''_!\mathcal{F}|_{X''} \to f_!\mathcal{F}$.
\end{remark}

\begin{lemma}
\label{lemma-compactify-f-shriek-separated}
Soit $Y$ un schéma. Soit $j : X \to \overline{X}$ une immersion
ouverte de schémas sur $Y$, avec $\overline{X}$ propre sur $Y$.
Notons $f : X \to Y$ et $\overline{f} : \overline{X} \to Y$
les morphismes structuraux. Pour $\mathcal{F} \in \textit{Ab}(X_\etale)$,
il existe un isomorphisme canonique (voir la démonstration)
$$
f_!\mathcal{F} \longrightarrow \overline{f}_!j_!\mathcal{F}
$$
Comme $\overline{f}_! = \overline{f}_*$ d'après le
lemme \ref{lemma-proper-f-shriek}, nous obtenons
$\overline{f}_* \circ j_! = f_!$ comme foncteurs de
$\textit{Ab}(X_\etale) \to \textit{Ab}(Y_\etale)$.
\end{lemma}

\begin{proof}
Nous avons $(j_!\mathcal{F})|_X = \mathcal{F}$ ; voir le
lemme \ref{etale-cohomology-lemma-jshriek-open} de Cohomologie étale.
Ainsi, la flèche affichée est l'application injective
$f_!(\mathcal{G}|_X) \to \overline{f}_!\mathcal{G}$
de la remarque \ref{remark-covariance-f-shriek-separated},
pour $\mathcal{G} = j_!\mathcal{F}$. La description explicite
de cette application montre qu'il suffit de prouver ceci : si $V \in Y_\etale$ et
$s \in \overline{f}_!\mathcal{G}(V) = \overline{f}_*\mathcal{G}(V) =
\mathcal{G}(\overline{X}_V)$
est une section, alors le support de $s$ est contenu dans l'ouvert
$X_V \subset \overline{X}_V$. Cela résulte immédiatement de ce que
les fibres de $\mathcal{G}$ sont nulles aux points géométriques
de $\overline{X} \setminus X$.
\end{proof}

\noindent
Nous voulons relier les fibres de $f_!\mathcal{F}$ aux sections à
support propre sur les fibres. Nous avons besoin pour cela d'une définition.

\begin{definition}
\label{definition-compact-support}
Soit $X$ un schéma séparé localement de type fini sur un corps $k$.
Soit $\mathcal{F}$ un faisceau abélien sur $X_\etale$. Nous définissons
$H^0_c(X, \mathcal{F}) \subset H^0(X, \mathcal{F})$ comme l'ensemble
des sections dont le support est propre sur $k$. Les éléments de
$H^0_c(X, \mathcal{F})$ sont appelés {\it sections à support propre}.
\end{definition}

\noindent
Avertissement : cette définition n'est pas la « bonne » si $X$ n'est pas
séparé sur $k$.

\begin{lemma}
\label{lemma-proper-compact-support}
Soit $X$ un schéma propre sur un corps $k$. Alors
$H^0_c(X, \mathcal{F}) = H^0(X, \mathcal{F})$.
\end{lemma}

\begin{proof}
Cela résulte immédiatement de la construction de $H^0_c$.
\end{proof}

\begin{remark}[Immersions ouvertes et sections à support propre]
\label{remark-covariance-compact-support}
Soit $X$ un schéma séparé localement de type fini sur un corps $k$.
Soit $\mathcal{F}$ un faisceau abélien sur $X_\etale$.
Exactement comme dans la remarque \ref{remark-covariance-f-shriek-separated},
pour tout ouvert $X' \subset X$, il existe une application injective
$$
H^0_c(X', \mathcal{F}|_{X'}) \longrightarrow H^0_c(X, \mathcal{F})
$$
et ces applications font de $H^0_c$ un « cofaisceau » sur le site de Zariski de $X$.
\end{remark}

\begin{lemma}
\label{lemma-compactify-compact-support}
Soit $k$ un corps. Soit $j : X \to \overline{X}$ une immersion
ouverte de schémas sur $k$, avec $\overline{X}$ propre sur $k$.
Pour $\mathcal{F} \in \textit{Ab}(X_\etale)$,
il existe un isomorphisme canonique (voir la démonstration)
$$
H^0_c(X, \mathcal{F}) \longrightarrow
H^0_c(\overline{X}, j_!\mathcal{F}) =
H^0(\overline{X}, j_!\mathcal{F})
$$
où l'égalité de droite résulte du
lemme \ref{lemma-proper-compact-support}.
\end{lemma}

\begin{proof}
Nous avons $(j_!\mathcal{F})|_X = \mathcal{F}$ ; voir le
lemme \ref{etale-cohomology-lemma-jshriek-open} de Cohomologie étale.
Ainsi, la flèche affichée est l'application injective
$H^0_c(X, \mathcal{G}|_X) \to H^0_c(\overline{X}, \mathcal{G})$
de la remarque \ref{remark-covariance-compact-support},
pour $\mathcal{G} = j_!\mathcal{F}$. La description explicite
de cette application montre qu'il suffit de prouver ceci : si
$s \in H^0(\overline{X}, \mathcal{G})$ est une section, alors le support de
$s$ est contenu dans l'ouvert $X$. Cela résulte immédiatement de ce que
les fibres de $\mathcal{G}$ sont nulles aux points géométriques
de $\overline{X} \setminus X$.
\end{proof}

\begin{lemma}
\label{lemma-stalk-f-shriek-separated}
Soit $f : X \to Y$ un morphisme de schémas séparé et
localement de type fini. Soit $\mathcal{F}$ un faisceau abélien sur
$X_\etale$. Il existe alors un isomorphisme canonique
$$
(f_!\mathcal{F})_{\overline{y}}
\longrightarrow
H^0_c(X_{\overline{y}}, \mathcal{F}|_{X_{\overline{y}}})
$$
pour tout point géométrique $\overline{y} : \Spec(k) \to Y$.
\end{lemma}

\begin{proof}
Rappelons que $(f_*\mathcal{F})_{\overline{y}} = \colim f_*\mathcal{F}(V)$,
où la limite inductive porte sur les voisinages étales $(V, \overline{v})$
de $\overline{y}$. Si $s \in f_*\mathcal{F}(V) = \mathcal{F}(X_V)$,
l'image inverse de $s$ est une section de $\mathcal{F}$ sur
$(X_V)_{\overline{v}} = X_{\overline{y}}$. Nous obtenons ainsi une application canonique
$$
c_{\overline{y}} :
(f_*\mathcal{F})_{\overline{y}}
\longrightarrow 
H^0(X_{\overline{y}}, \mathcal{F}|_{X_{\overline{y}}})
$$
Nous affirmons que cette application induit une bijection entre les sous-groupes
$(f_!\mathcal{F})_{\overline{y}}$ et
$H^0_c(X_{\overline{y}}, \mathcal{F}|_{X_{\overline{y}}})$.
Cette affirmation implique le lemme, mais elle est légèrement plus précise,
car elle décrit l'identification du lemme comme étant donnée
par l'image inverse des sections de $\mathcal{F}$ sur la fibre géométrique de $f$.

\medskip\noindent
Observons que tout élément
$s \in (f_!\mathcal{F})_{\overline{y}} \subset (f_*\mathcal{F})_{\overline{y}}$
est envoyé par $c_{\overline{y}}$ sur un élément de
$H^0_c(X_{\overline{y}}, \mathcal{F}|_{X_{\overline{y}}}) \subset
H^0(X_{\overline{y}}, \mathcal{F}|_{X_{\overline{y}}})$.
En effet, la formation du support d'une section
commute au changement de base et la propreté est préservée par
changement de base. Ceci construit déjà l'application de l'énoncé du lemme.
Pour montrer que c'est un isomorphisme, nous pouvons travailler localement
pour la topologie de Zariski sur $Y$ ; nous supposons donc $Y$ affine.

\medskip\noindent
Nous utiliserons ci-dessous l'observation suivante :
étant donné un sous-schéma ouvert $X' \subset X$,
si $f' = f|_{X'}$, nous obtenons un diagramme commutatif
$$
\xymatrix{
(f'_!(\mathcal{F}|_{X'}))_{\overline{y}} \ar[r] \ar[d] &
H^0_c(X'_{\overline{y}}, \mathcal{F}|_{X'_{\overline{y}}}) \ar[d] \\
(f_!\mathcal{F})_{\overline{y}} \ar[r] &
H^0_c(X_{\overline{y}}, \mathcal{F}|_{X_{\overline{y}}})
}
$$
où les flèches horizontales sont les applications construites ci-dessus et
les flèches verticales sont celles des
remarques \ref{remark-covariance-f-shriek-separated} et
\ref{remark-covariance-compact-support}.
En effet, étant donnés un voisinage étale $(V, \overline{v})$
de $\overline{y}$ et une section $s \in f_*\mathcal{F}(V) = \mathcal{F}(X_V)$
dont le support $Z$ est contenu dans $X'_V$ et propre sur $V$,
de sorte que $s$ détermine un élément de chacun des groupes
$(f'_!(\mathcal{F}|_{X'}))_{\overline{y}}$ et
$(f_!\mathcal{F})_{\overline{y}}$ qui se correspondent par
la flèche verticale du diagramme, les flèches horizontales envoient alors
ces éléments sur l'image inverse $s|_{X_{\overline{y}}}$ de $s$ dans
$X_{\overline{y}}$, dont le support $Z_{\overline{y}}$ est contenu dans
$X'_{\overline{y}}$ ; cette restriction détermine donc
une paire compatible d'éléments de
$H^0_c(X'_{\overline{y}}, \mathcal{F}|_{X'_{\overline{y}}})$
et de
$H^0_c(X_{\overline{y}}, \mathcal{F}|_{X_{\overline{y}}})$.

\medskip\noindent
Supposons que $s \in (f_!\mathcal{F})_{\overline{y}}$ ait une image nulle dans
$H^0_c(X_{\overline{y}}, \mathcal{F}|_{X_{\overline{y}}})$.
Disons que $s$ correspond à $s \in f_*\mathcal{F}(V) = \mathcal{F}(X_V)$,
de support $Z$ propre sur $V$. Nous pouvons supposer $V$ affine,
et $Z$ est alors quasi-compact. Nous pouvons choisir un ouvert quasi-compact
$X' \subset X$ contenant l'image de $Z$. Alors $Z$ est contenu dans
$X'_V$ et $s$ est donc l'image d'un élément
$s' \in f'_!(\mathcal{F}|_{X'})(V)$, où $f' = f|_{X'}$, comme au
paragraphe précédent. L'image de $s'$ dans
$H^0_c(X'_{\overline{y}}, \mathcal{F}|_{X'_{\overline{y}}})$
est donc nulle. Pour démontrer l'injectivité, nous pouvons remplacer $X$ par
$X'$, c'est-à-dire supposer $X$ quasi-compact. Nous démontrerons
ce cas ci-dessous.

\medskip\noindent
Supposons que
$t \in H^0_c(X_{\overline{y}}, \mathcal{F}|_{X_{\overline{y}}})$.
Alors le support de $t$ est contenu dans un sous-schéma ouvert
quasi-compact $W \subset X_{\overline{y}}$.
Nous pouvons donc trouver un sous-schéma ouvert quasi-compact
$X' \subset X$ tel que $X'_{\overline{y}}$ contienne $W$.
Il est alors clair que $t$ appartient à l'image
de l'application injective
$H^0_c(X'_{\overline{y}}, \mathcal{F}|_{X'_{\overline{y}}}) \to
H^0_c(X_{\overline{y}}, \mathcal{F}|_{X_{\overline{y}}})$.
Pour démontrer la surjectivité, nous pouvons donc remplacer $X$
par $X'$, c'est-à-dire supposer $X$ quasi-compact.
Nous démontrerons ce cas ci-dessous.

\medskip\noindent
Dans ce dernier paragraphe, nous démontrons le lemme lorsque
$X$ est quasi-compact et $Y$ affine. D'après le théorème
\ref{flat-theorem-nagata} de Compléments sur la platitude, il existe une compactification
$j : X \to \overline{X}$ sur $Y$. Posons $\mathcal{G} = j_!\mathcal{F}$,
de sorte que $\mathcal{F} = \mathcal{G}|_X$ d'après le
lemme \ref{etale-cohomology-lemma-jshriek-open} de Cohomologie étale.
La discussion ci-dessus donne un diagramme commutatif
$$
\xymatrix{
(f_!\mathcal{F})_{\overline{y}} \ar[r] \ar[d] &
H^0_c(X_{\overline{y}}, \mathcal{F}|_{X_{\overline{y}}}) \ar[d] \\
(\overline{f}_!\mathcal{G})_{\overline{y}} \ar[r] &
H^0_c(\overline{X}_{\overline{y}}, \mathcal{G}|_{\overline{X}_{\overline{y}}})
}
$$
D'après les lemmes \ref{lemma-compactify-f-shriek-separated} et
\ref{lemma-compactify-compact-support}, les applications verticales
sont des isomorphismes. Nous sommes ramenés au cas du morphisme propre
$\overline{X} \to Y$. Pour un morphisme propre, notre application
est un isomorphisme d'après les
lemmes \ref{lemma-proper-f-shriek} et \ref{lemma-proper-compact-support},
ainsi que le changement de base propre pour les images directes ; voir le
lemme de Cohomologie étale
\ref{etale-cohomology-lemma-proper-pushforward-stalk}.
\end{proof}

\begin{lemma}
\label{lemma-base-change-f-shriek-separated}
Considérons un carré cartésien
$$
\xymatrix{
X' \ar[r]_{g'} \ar[d]_{f'} & X \ar[d]^f \\
Y' \ar[r]^g & Y
}
$$
de schémas, où $f$ est séparé et localement de type fini.
Pour tout faisceau abélien $\mathcal{F}$ sur $X_\etale$, nous avons
$f'_!(g')^{-1}\mathcal{F} = g^{-1}f_!\mathcal{F}$.
\end{lemma}

\begin{proof}
Plus généralement, il existe une application de changement de base
$g^{-1}f_*\mathcal{F} \to f'_*(g')^{-1}\mathcal{F}$ ; voir
Sites, section \ref{sites-section-pullback}.
Nous affirmons que cette application envoie $g^{-1}f_!\mathcal{F}$
dans le sous-faisceau $f'_!(g')^{-1}\mathcal{F}$
et induit l'isomorphisme du lemme.

\medskip\noindent
Choisissons un point géométrique $\overline{y}': \Spec(k) \to Y'$ et notons
$\overline{y} = g \circ \overline{y}'$ son image dans $Y$. Il existe un
diagramme commutatif
$$
\xymatrix{
(f_*\mathcal{F})_{\overline{y}} \ar[r] \ar[d] &
H^0(X_{\overline{y}}, \mathcal{F}|_{X_{\overline{y}}}) \ar[d] \\
(f'_*(g')^{-1}\mathcal{F})_{\overline{y}'} \ar[r] &
H^0(X'_{\overline{y}'}, (g')^{-1}\mathcal{F}|_{X'_{\overline{y}'}})
}
$$
où les applications horizontales sont celles utilisées dans la démonstration du
lemme \ref{lemma-stalk-f-shriek-separated},
et les applications verticales sont celles de changement de base ci-dessus.
Le diagramme commute, car chacune des quatre applications
considérées s'obtient en prenant l'image inverse de sections locales par
un morphisme de schémas, et le diagramme sous-jacent de morphismes
de schémas commute. Puisque le diagramme de l'énoncé du lemme
est cartésien, nous avons $X'_{\overline{y}'} = X_{\overline{y}}$.
Le lemme \ref{lemma-stalk-f-shriek-separated}
et sa démonstration donnent donc un diagramme commutatif
$$
\xymatrix{
(f_*\mathcal{F})_{\overline{y}} \ar[rrr] \ar[ddd] & & &
H^0(X_{\overline{y}}, \mathcal{F}|_{X_{\overline{y}}}) \ar[ddd] \\
& (f_!\mathcal{F})_{\overline{y}} \ar[r] \ar@{..>}[d] \ar[lu] &
H^0_c(X_{\overline{y}}, \mathcal{F}|_{X_{\overline{y}}}) \ar[d] \ar[ru] \\
& (f'_!(g')^{-1}\mathcal{F})_{\overline{y}'} \ar[r] \ar[ld] &
H^0_c(X'_{\overline{y}'}, (g')^{-1}\mathcal{F}|_{X'_{\overline{y}'}}) \ar[rd]\\
(f'_*(g')^{-1}\mathcal{F})_{\overline{y}'} \ar[rrr] & & &
H^0(X'_{\overline{y}'}, (g')^{-1}\mathcal{F}|_{X'_{\overline{y}'}})
}
$$
où les flèches horizontales du carré intérieur sont des isomorphismes
et les deux flèches verticales de droite sont des égalités. En outre, les
flèches sud-est, sud-ouest, nord-est et nord-ouest sont injectives. Il existe donc
une unique flèche pointillée bijective complétant le diagramme. Nous en concluons que
$g^{-1}f_!\mathcal{F} \subset g^{-1}f_*\mathcal{F} \to f'_*(g')^{-1}\mathcal{F}$
est envoyé dans le sous-faisceau
$f'_!(g')^{-1}\mathcal{F} \subset f'_*(g')^{-1}\mathcal{F}$
car il en est ainsi sur les fibres ; voir le théorème
\ref{etale-cohomology-theorem-exactness-stalks} de Cohomologie étale.
Le même théorème implique alors que l'application induite est un isomorphisme,
ce qui achève la démonstration.
\end{proof}

\begin{lemma}
\label{lemma-f-shriek-composition}
Soient $f : X \to Y$ et $g : Y \to Z$ des morphismes composables de schémas,
séparés et localement de type fini. Soit $\mathcal{F}$ un faisceau abélien
sur $X_\etale$. Alors $g_!f_!\mathcal{F} = (g \circ f)_!\mathcal{F}$
comme sous-faisceaux de $(g \circ f)_*\mathcal{F}$.
\end{lemma}

\begin{proof}
Nous recommandons vivement au lecteur de le démontrer lui-même.
Soient $W \in Z_\etale$ et
$s \in (g \circ f)_*\mathcal{F}(W) = \mathcal{F}(X_W)$.
Notons $T \subset X_W$ le support de $s$ ; c'est un sous-ensemble
fermé. La section $s$ appartient à $(g \circ f)_!\mathcal{F}$
si et seulement si $T$ est propre sur $W$. Nous avons
$f_!\mathcal{F} \subset f_*\mathcal{F}$, et donc
$g_!f_!\mathcal{F} \subset g_!f_*\mathcal{F} \subset g_*f_*\mathcal{F}$.
D'autre part, $s$ appartient à $g_!f_!\mathcal{F}$ si et seulement
si (a) $T$ est propre sur $Y_W$ et (b) le support $T'$ de $s$,
considérée comme section de $f_!\mathcal{F}$, est propre sur $W$.
Si (a) est vérifiée, l'image de $T$ dans $Y_W$ est fermée et, puisque
$f_!\mathcal{F} \subset f_*\mathcal{F}$, nous voyons que
$T' \subset Y_W$ est l'image de $T$ (détails omis ; considérer les fibres).

\medskip\noindent
Il reste donc à montrer qu'un sous-ensemble fermé $T \subset X_W$
est propre sur $W$ si et seulement si $T$ est propre sur $Y_W$
et si l'image de $T$ dans $Y_W$ est propre sur $W$. Munissons $T$
de la structure réduite induite de sous-schéma fermé.
Si $T$ est propre sur $W$, alors $T \to Y_W$ est propre d'après le
lemme \ref{morphisms-lemma-image-proper-scheme-closed} de Morphismes,
et l'image de $T$ dans $Y_W$ est propre sur $W$ d'après le
lemme de Cohomologie des schémas
\ref{coherent-lemma-functoriality-closed-proper-over-base}.
Réciproquement, si $T$ est propre sur $Y_W$
et si l'image de $T$ dans $Y_W$ est propre sur $W$,
alors le morphisme $T \to W$ est propre comme composé
de morphismes propres (nous munissons ici l'image fermée de $T$
dans $Y_W$ de sa structure réduite induite, afin de ramener la
question à des morphismes de schémas) ; voir le
lemme \ref{morphisms-lemma-composition-proper} de Morphismes.
\end{proof}

\begin{remark}
\label{remark-f-shriek-base-change-composition}
Les isomorphismes entre foncteurs
construits ci-dessus possèdent les deux propriétés suivantes :
\begin{enumerate}
\item Soient $f : X \to Y$, $g : Y \to Z$ et $h : Z \to T$ des morphismes
composables de schémas, séparés et localement de type fini.
Alors le diagramme
$$
\xymatrix{
(h \circ g \circ f)_! \ar[r] \ar[d] &
(h \circ g)_! \circ f_! \ar[d] \\
h_! \circ (g \circ f)_! \ar[r] &
h_! \circ g_! \circ f_!
}
$$
est commutatif, les flèches étant celles du lemme \ref{lemma-f-shriek-composition}.
\item Supposons donné un diagramme de schémas
$$
\xymatrix{
X' \ar[d]_{f'} \ar[r]_c & X \ar[d]^f \\
Y' \ar[d]_{g'} \ar[r]_b & Y \ar[d]^g \\
Z' \ar[r]^a & Z
}
$$
dont les deux carrés sont cartésiens, et où $f$ et $g$ sont séparés et
localement de type fini. Alors le diagramme
$$
\xymatrix{
a^{-1} \circ (g \circ f)_! \ar[d] \ar[rr] & &
(g' \circ f')_! \circ c^{-1} \ar[d] \\
a^{-1} \circ g_! \circ f_! \ar[r] &
g'_! \circ b^{-1} \circ f_! \ar[r] &
g'_! \circ f'_! \circ c^{-1}
}
$$
est commutatif, les flèches horizontales étant celles du
lemme \ref{lemma-base-change-f-shriek-separated}
et les autres celles du lemme \ref{lemma-f-shriek-composition}.
\end{enumerate}
La partie (1) est vraie parce qu'il existe un diagramme commutatif
analogue pour les images directes. La partie (2) résulte de la compatibilité
très générale des applications de changement de base pour les images directes
(Sites, remarque \ref{sites-remark-compose-base-change})
et du fait que les isomorphismes des
lemmes \ref{lemma-base-change-f-shriek-separated} et
\ref{lemma-f-shriek-composition}
sont construits à l'aide des applications de changement de base correspondantes pour les images directes.
\end{remark}

\begin{lemma}
\label{lemma-colim-f-shriek-separated}
Soit $f : X \to Y$ un morphisme de schémas séparé et
localement de type fini. Soit $X = \bigcup_{i \in I} X_i$ un
recouvrement ouvert tel que, pour tous $i, j \in I$, il existe $k$
avec $X_i \cup X_j \subset X_k$. Notons $f_i : X_i \to Y$
la restriction de $f$. Alors
$$
f_!\mathcal{F} = \colim_{i \in I} f_{i, !}(\mathcal{F}|_{X_i})
$$
fonctoriellement en $\mathcal{F} \in \textit{Ab}(X_\etale)$,
où les morphismes de transition sont ceux construits dans la
remarque \ref{remark-covariance-f-shriek-separated}.
\end{lemma}

\begin{proof}
Il suffit de montrer que l'application canonique de
droite à gauche est bijective lorsqu'on l'évalue sur un objet
quasi-compact $V$ de $Y_\etale$.
La limite inductive du membre de droite est filtrante
et ses morphismes de transition sont injectifs.
Nous pouvons donc utiliser le
lemme \ref{sites-lemma-directed-colimits-sections} de Sites
pour évaluer cette limite. L'énoncé se ramène ainsi à observer
qu'un fermé $Z \subset X_V$ propre sur $V$ est quasi-compact
et est donc contenu dans $X_{i, V}$ pour un certain $i$.
\end{proof}

\begin{lemma}
\label{lemma-f-shriek-separated-direct-sums}
Soit $f : X \to Y$ un morphisme de schémas séparé et
localement de type fini. Alors le foncteur $f_!$ commute aux sommes directes.
\end{lemma}

\begin{proof}
Soit $\mathcal{F} = \bigoplus \mathcal{F}_i$. Pour montrer que l'application
$\bigoplus f_!\mathcal{F}_i \to f_!\mathcal{F}$ est un isomorphisme,
il suffit de montrer que ces faisceaux ont les mêmes sections sur
un objet quasi-compact $V$ de $Y_\etale$. En remplaçant $Y$ par $V$,
il suffit de montrer que
$H^0(Y, f_!\mathcal{F}) \subset H^0(X, \mathcal{F})$
est égal à
$\bigoplus H^0(Y, f_!\mathcal{F}_i)
\subset \bigoplus H^0(X, \mathcal{F}_i)
\subset H^0(X, \bigoplus \mathcal{F}_i)$.
Dans ce cas, en écrivant $X$ comme réunion de ses ouverts quasi-compacts
et en utilisant le lemme \ref{lemma-colim-f-shriek-separated},
nous nous ramenons au cas où $X$ est également quasi-compact.
Alors $H^0(X, \mathcal{F}) = \bigoplus H^0(X, \mathcal{F}_i)$
d'après le théorème \ref{etale-cohomology-theorem-colimit} de Cohomologie étale.
Le lecteur conclura aisément en considérant les supports des sections.
\end{proof}

\begin{lemma}
\label{lemma-lqf-f-shriek-separated-colimits}
Soit $f : X \to Y$ un morphisme de schémas séparé et
localement quasi-fini. Alors
\begin{enumerate}
\item pour $\mathcal{F}$ dans $\textit{Ab}(X_\etale)$ et un point
géométrique $\overline{y} : \Spec(k) \to Y$, nous avons
$$
(f_!\mathcal{F})_{\overline{y}} =
\bigoplus\nolimits_{f(\overline{x}) = \overline{y}} \mathcal{F}_{\overline{x}}
$$
fonctoriellement en $\mathcal{F}$, et
\item le foncteur $f_!$ est exact.
\end{enumerate}
\end{lemma}

\begin{proof}
Le foncteur $f_!$ est exact à gauche par construction. L'exactitude à droite
se vérifie sur les fibres
(Cohomologie étale, théorème \ref{etale-cohomology-theorem-exactness-stalks}).
Il suffit donc de démontrer la partie (1).

\medskip\noindent
Soit $\overline{y} : \Spec(k) \to Y$ un point géométrique.
L'espace topologique sous-jacent au schéma $X_{\overline{y}}$
est discret
(Morphismes, lemme \ref{morphisms-lemma-locally-quasi-finite-fibres}),
et tous les corps résiduels en ses points sont égaux à $k$
(comme extensions finies de $k$). Par conséquent,
$\{\overline{x} : \Spec(k) \to X : f(\overline{x}) = \overline{y}\}$
est égal à l'ensemble des points de $X_{\overline{y}}$.
Le calcul de la fibre résulte donc du
lemme plus général \ref{lemma-stalk-f-shriek-separated}.
\end{proof}









\section{sections à support fini}
\label{section-finite-support}

\noindent
Dans cette section, nous étendons la construction de la
section \ref{section-compact-support} aux morphismes localement
quasi-finis qui ne sont pas nécessairement séparés.

\medskip\noindent
Soit $f : X \to Y$ un morphisme de schémas localement quasi-fini.
Soit $\mathcal{F}$ un faisceau abélien sur $X_\etale$. Pour $V$ dans
$Y_\etale$, notons $X_V = X \times_Y V$ le changement de base. Nous allons
considérer le groupe des sommes formelles finies
\begin{equation}
\label{equation-formal-sum}
s = \sum\nolimits_{i = 1, \ldots, n} (Z_i, s_i)
\end{equation}
où $Z_i \subset X_V$ est un sous-schéma localement fermé tel que le
morphisme $Z_i \to V$ soit fini\footnote{Puisque $f$ est localement quasi-fini,
le morphisme $Z_i \to V$ est fini si et seulement s'il est propre.}
et où $s_i \in H_{Z_i}(\mathcal{F})$. Ici, comme dans la
section \ref{section-growing}, nous posons
$$
H_{Z_i}(\mathcal{F}) =
\{s_i \in \mathcal{F}(U_i) \mid \text{Supp}(s_i) \subset Z_i\}
$$
où $U_i \subset X_V$ est un sous-schéma ouvert contenant $Z_i$ comme
sous-schéma fermé. Nous considérons ces sommes formelles modulo les
relations suivantes :
\begin{enumerate}
\item
\label{item-sum}
$(Z, s) + (Z, s') = (Z, s + s')$,
\item
\label{item-sub}
$(Z, s) = (Z', s)$ si $Z \subset Z'$.
\end{enumerate}
La seconde relation a bien un sens : puisque $Z \to V$ est fini
et $Z' \to V$ séparé, l'inclusion $Z \to Z'$ est fermée et nous
pouvons utiliser l'application décrite en (\ref{item-inclusion}).

\medskip\noindent
Notons $f_{p!}\mathcal{F}(V)$ le quotient du groupe abélien
des sommes formelles (\ref{equation-formal-sum}) par ces relations.
La première relation dit que $f_{p!}\mathcal{F}(V)$ est un quotient
de la somme directe des groupes abéliens $H_Z(\mathcal{F})$,
où $Z \subset X_V$ parcourt les sous-schémas localement fermés finis sur $V$.
La seconde relation dit qu'il s'agit en réalité de la limite inductive
\begin{equation}
\label{equation-colimit-definition}
f_{p!}\mathcal{F}(V) = \colim_Z H_Z(\mathcal{F})
\end{equation}
Cette formule donne une manière abstraite commode de considérer
notre construction.

\medskip\noindent
Observons ensuite que cette construction donne naturellement un préfaisceau
$f_{p!}\mathcal{F}$ de groupes abéliens sur $Y_\etale$.
En effet, un morphisme $V' \to V$ dans $Y_\etale$ fournit le morphisme de changement de base
$X_{V'} \to X_V$. Si $Z \subset X_V$ est un sous-schéma localement fermé
fini sur $V$, alors l'image inverse schématique $Z' \subset X_{V'}$
est finie sur $V'$. En outre, si $U \subset X_V$ est un ouvert tel
que $Z$ soit fermé dans $U$, alors l'image inverse $U' \subset X_{V'}$
est un ouvert tel que $Z'$ soit fermé dans $U'$. L'application de restriction
$\mathcal{F}(U) \to \mathcal{F}(U')$ de $\mathcal{F}$
envoie donc $H_Z(\mathcal{F})$ dans $H_{Z'}(\mathcal{F})$ ; c'est un cas
particulier de la fonctorialité décrite en (\ref{item-pullback}) ci-dessus.
Ces applications sont manifestement compatibles avec les inclusions
$Z_1 \subset Z_2$ de tels sous-schémas localement fermés de $X_V$, et
nous obtenons une application
$$
f_{p!}\mathcal{F}(V) = \colim_Z H_Z(\mathcal{F})
\longrightarrow
\colim_{Z'} H_{Z'}(\mathcal{F}) =
f_{p!}\mathcal{F}(V')
$$
Ces applications font bien de $f_{p!}\mathcal{F}$ un préfaisceau de groupes
abéliens sur $Y_\etale$. Nous omettons les détails.

\medskip\noindent
Enfin, observons que la construction de $f_{p!}\mathcal{F}$
est fonctorielle en $\mathcal{F}$ dans $\textit{Ab}(X_\etale)$.
Ainsi, à tout morphisme localement quasi-fini $f : X \to Y$,
nous avons associé un foncteur
$$
f_{p!} :
\textit{Ab}(X_\etale)
\longrightarrow
\textit{PAb}(Y_\etale)
$$
de la catégorie des faisceaux abéliens sur $X_\etale$ vers celle
des préfaisceaux abéliens sur $Y_\etale$. Avant de définir $f_!$ par
passage au faisceau associé à ce foncteur, vérifions qu'il coïncide avec
la construction de la section \ref{section-compact-support}
et avec celle de la
section \ref{etale-cohomology-section-extension-by-zero} de Cohomologie étale
lorsqu'elles s'appliquent toutes deux.

\begin{lemma}
\label{lemma-finite-support-f-shriek-separated}
Soit $f : X \to Y$ un morphisme de schémas séparé et localement
quasi-fini. Fonctoriellement en $\mathcal{F} \in \textit{Ab}(X_\etale)$,
il existe un isomorphisme canonique (!)
$$
f_{p!}\mathcal{F} \longrightarrow f_!\mathcal{F}
$$
de préfaisceaux abéliens qui identifie le faisceau
$f_!\mathcal{F}$ de la définition \ref{definition-f-shriek-separated}
au préfaisceau $f_{p!}\mathcal{F}$ construit ci-dessus.
\end{lemma}

\begin{proof}
Soit $V$ un objet de $Y_\etale$. Si $Z \subset X_V$ est localement fermé
et fini sur $V$, alors, puisque $f$ est séparé, le
morphisme $Z \to X_V$ est une immersion fermée. De plus, si
$Z_i$, $i = 1, \ldots, n$, sont des sous-schémas fermés de $X_V$ finis
sur $V$, alors $Z_1 \cup \ldots \cup Z_n$ (réunion schématique)
est un sous-schéma fermé fini sur $V$. Dans ce cas, la limite inductive
(\ref{equation-colimit-definition}) qui définit $f_{p!}\mathcal{F}(V)$
est donc filtrante, et $f_{!p}\mathcal{F}(V)$ est simplement égal
à l'ensemble des sections de $\mathcal{F}(X_V)$ dont le support est fini sur $V$.
Tout fermé de $X_V$ propre sur $V$ étant
en fait fini sur $V$ (car $f$ est localement quasi-fini),
nous concluons que cet ensemble est égal à $f_!\mathcal{F}(V)$
par définition.
\end{proof}

\begin{lemma}
\label{lemma-finite-support-stalk}
Soit $f : X \to Y$ un morphisme de schémas localement quasi-fini.
Soit $\overline{y} : \Spec(k) \to Y$ un point géométrique.
Fonctoriellement en $\mathcal{F}$ dans $\textit{Ab}(X_\etale)$, nous avons
$$
(f_{p!}\mathcal{F})_{\overline{y}} =
\bigoplus\nolimits_{f(\overline{x}) = \overline{y}} \mathcal{F}_{\overline{x}}
$$
\end{lemma}

\begin{proof}
Rappelons que la fibre en $\overline{y}$ d'un préfaisceau est définie par la
limite inductive usuelle sur les voisinages étales $(V, \overline{v})$
de $\overline{y}$ ; voir la définition
\ref{etale-cohomology-definition-stalk} de Cohomologie étale. Supposons donc que
$s = \sum_{i = 1, \ldots, n} (Z_i, s_i)$, comme dans
(\ref{equation-formal-sum}), soit un élément de $f_{p!}\mathcal{F}(V)$,
où $(V, \overline{v})$ est un voisinage étale de $\overline{y}$.
Alors, puisque
$$
X_{\overline{y}} = (X_V)_{\overline{v}} \supset Z_{i, \overline{v}}
$$
et puisque $s_i$ est une section de $\mathcal{F}$ sur un voisinage ouvert
de $Z_i$ dans $X_V$, nous pouvons envoyer $s$ sur
$$
\sum\nolimits_{i = 1, \ldots, n} 
\sum\nolimits_{\overline{x} \in Z_{i, \overline{v}}}
\left(\text{classe de }s_i\text{ dans }\mathcal{F}_{\overline{x}}\right)
\quad\in\quad
\bigoplus\nolimits_{f(\overline{x}) = \overline{y}} \mathcal{F}_{\overline{x}}
$$
Nous omettons de vérifier la compatibilité avec les applications de restriction
et le fait que les relations (\ref{item-sum}) $(Z, s) + (Z, s') - (Z, s + s')$
et (\ref{item-sub}) $(Z, s) - (Z', s)$ si $Z \subset Z'$ sont envoyées sur zéro.
Nous obtenons ainsi une application
$$
(f_{p!}\mathcal{F})_{\overline{y}}
\longrightarrow
\bigoplus\nolimits_{f(\overline{x}) = \overline{y}} \mathcal{F}_{\overline{x}}
$$

\medskip\noindent
Montrons que cette flèche est surjective. Il suffit de choisir
$\overline{x}$ tel que $f(\overline{x}) = \overline{y}$ et de montrer qu'un
élément $s$ du facteur direct $\mathcal{F}_{\overline{x}}$ appartient à
l'image. Supposons que $s$ soit représenté par $s \in \mathcal{F}(U)$,
où $(U, \overline{u})$ est un voisinage étale de $\overline{x}$.
Puisque $f$ est localement quasi-fini, le morphisme $U \to Y$
l'est également. D'après le lemme de Compléments sur les morphismes
\ref{more-morphisms-lemma-etale-makes-quasi-finite-finite-multiple-points-var}
nous pouvons trouver un voisinage étale $(V, \overline{v})$ de
$\overline{y}$, un sous-schéma ouvert
$$
W \subset U \times_Y V,
$$
et un point géométrique $\overline{w}$ s'envoyant sur $\overline{u}$ et
$\overline{v}$, tels que $W \to V$ soit fini et que $\overline{w}$ soit l'unique
point géométrique de $W$ s'envoyant sur $\overline{v}$. (Nous omettons le passage
du langage des points géométriques employé ici à celui des points et des
extensions de corps résiduels utilisé dans
l'énoncé du lemme.) Le morphisme $W \to X_V = X \times_Y V$
est étale. Choisissons un voisinage ouvert affine $W' \subset X_V$
de l'image $\overline{w}'$ de $\overline{w}$. Puisque $\overline{w}$
est l'unique point de $W$ au-dessus de $\overline{v}$ et que $W \to V$
est fermé, nous pouvons, après avoir remplacé $V$ par un voisinage ouvert de $\overline{v}$,
supposer que $W \to X_V$ se factorise par $W'$. Alors $W \to W'$ est fini et
étale, et il existe un unique point géométrique $\overline{w}$ de $W$
au-dessus de $\overline{w}'$. Il s'ensuit que $W \to W'$ est une immersion ouverte
au-dessus d'un voisinage ouvert de $\overline{w}'$ dans $W'$ ; voir le
lemme \ref{etale-lemma-finite-etale-one-point} de Morphismes étales.
En rétrécissant $V$ et $W'$, nous pouvons supposer que $W \to W'$ est un isomorphisme.
Ainsi, $s$ peut être regardée comme une section $s'$ de $\mathcal{F}$ sur
le sous-schéma ouvert $W' \subset X_V$, qui est fini sur $V$.
Par définition, $(W', s')$ définit donc un élément de $j_{p!}\mathcal{F}(V)$
dont l'image est $s$, comme souhaité.

\medskip\noindent
Montrons que la flèche est injective. Pour cela, soit
$s = \sum_{i = 1, \ldots, n} (Z_i, s_i)$, comme dans (\ref{equation-formal-sum}),
un élément de $f_{p!}\mathcal{F}(V)$, où $(V, \overline{v})$ est un
voisinage étale de $\overline{y}$. Supposons que l'image de $s$ soit nulle
par l'application construite ci-dessus. Après avoir tout d'abord remplacé
$(V, \overline{v})$ par un voisinage étale de lui-même,
nous pouvons supposer qu'il existe des décompositions
$Z_i = Z_{i, 1} \amalg \ldots \amalg Z_{i, m_i}$ en sous-schémas
ouverts et fermés, telles que chaque $Z_{i, j}$ possède exactement un point géométrique
au-dessus de $\overline{v}$. Dans la décomposition évidente en somme directe
$$
H_{Z_i}(\mathcal{F}) = \bigoplus H_{Z_{i, j}}(\mathcal{F})
$$
supposons que l'élément $s_i$ corresponde à $\sum s_{i, j}$. Les relations
(\ref{item-sum}) et (\ref{item-sub}) permettent de remplacer $s$ par
$\sum_{i = 1, \ldots, n} \sum_{j = 1, \ldots, m_i} (Z_{i, j}, s_{i, j})$.
Autrement dit, nous pouvons supposer que $Z_i$ possède un unique point géométrique
au-dessus de $\overline{v}$. Soient $\overline{x}_1, \ldots, \overline{x}_m$
les points géométriques de $X$ au-dessus de $\overline{y}$ correspondant aux
points géométriques de nos $Z_i$ au-dessus de $\overline{v}$ ; notons que, pour
un même $j \in \{1, \ldots, m\}$, plusieurs indices $i$ peuvent être tels que
$\overline{x}_j$ corresponde à un point de $Z_i$.
D'après le lemme de Compléments sur les morphismes
\ref{more-morphisms-lemma-etale-makes-quasi-finite-finite-multiple-points-var}
appliqué à $X_V \to V$, après avoir remplacé
$(V, \overline{v})$ par un voisinage étale de lui-même,
nous pouvons supposer qu'il existe des sous-schémas ouverts
$$
W_j \subset X \times_Y V,\quad j = 1, \ldots, m
$$
et un point géométrique $\overline{w}_j$ de $W_j$ s'envoyant sur $\overline{x}_j$ et
$\overline{v}$, tels que $W_j \to V$ soit fini et que $\overline{w}_j$ soit l'unique
point géométrique de $W_j$ s'envoyant sur $\overline{v}$.
Après avoir rétréci $V$, nous pouvons supposer que $Z_i \subset W_j$ pour un certain $j$,
et nous disposons de l'application $H_{Z_i}(\mathcal{F}) \to H_{W_j}(\mathcal{F})$.
La relation (\ref{item-sub}) montre donc
que notre élément est équivalent à un élément de la forme
$$
\sum\nolimits_{j = 1, \ldots, m} (W_j, t_j)
$$
pour certains $t_j \in H_{W_j}(\mathcal{F})$. Il est clair que cet élément a
simplement pour image la classe de $t_j$ dans le facteur $\mathcal{F}_{\overline{x}_j}$.
Puisque l'image de $s$ est nulle, celle de $t_j$ dans
$\mathcal{F}_{\overline{x}_j}$ l'est aussi. Il s'ensuit que la restriction de $t_j$
est nulle sur un voisinage ouvert de $\overline{w}_j$ dans $W_j$ ; voir le
lemme \ref{etale-cohomology-lemma-zero-over-image} de Cohomologie étale.
En rétrécissant encore $V$, nous obtenons $t_j = 0$ pour tout $j$, comme souhaité.
\end{proof}

\begin{lemma}
\label{lemma-finite-support-etale-shriek}
Soit $f = j : U \to X$ un morphisme étale de schémas. Notons $j_{p!}$
la construction donnée par la formule de Cohomologie étale
(\ref{etale-cohomology-equation-j-p-shriek})
et $f_{p!}$ la construction précédente. Fonctoriellement en
$\mathcal{F} \in \textit{Ab}(X_\etale)$, il existe une application canonique
$$
j_{p!}\mathcal{F} \longrightarrow f_{p!}\mathcal{F}
$$
de préfaisceaux abéliens qui identifie le faisceau
$j_!\mathcal{F} = (j_{p!}\mathcal{F})^\#$ de la définition
\ref{etale-cohomology-definition-extension-zero} de Cohomologie étale
à $(f_{p!}\mathcal{F})^\#$.
\end{lemma}

\begin{proof}
On pourra lire la démonstration du lemme de Cohomologie étale
\ref{etale-cohomology-lemma-shriek-into-star-separated-etale}
avant celle du présent lemme.
Soit $V$ un objet de $X_\etale$. Rappelons que
$$
j_{p!}\mathcal{F}(V) =
\bigoplus\nolimits_{\varphi : V \to U} \mathcal{F}(V \xrightarrow{\varphi} U)
$$
À $\varphi$ est associé un sous-schéma ouvert
$Z_\varphi \subset U_V = U \times_X V$, à savoir
l'image du graphe de $\varphi$. Grâce à $\varphi$,
nous obtenons un isomorphisme $V \to Z_\varphi$ sur $U$
et pouvons regarder un élément
$$
s_\varphi \in \mathcal{F}(V \xrightarrow{\varphi} U) =
\mathcal{F}(Z_\varphi) = H_{Z_\varphi}(\mathcal{F})
$$
comme une section de $\mathcal{F}$ sur $Z_{\varphi}$. Puisque
$Z_\varphi \subset U_V$ est ouvert, nous avons en fait
$H_{Z_\varphi}(\mathcal{F}) = \mathcal{F}(Z_\varphi)$
et pouvons regarder $s_\varphi$ comme un élément de $H_{Z_\varphi}(\mathcal{F})$.
Cela étant, notre application $j_{p!}\mathcal{F} \to f_{p!}\mathcal{F}$
est définie par la règle
$$
\sum\nolimits_{i = 1, \ldots, n} s_{\varphi_i}
\longmapsto
\sum\nolimits_{i = 1, \ldots, n} (Z_{\varphi_i}, s_{\varphi_i})
$$
où le membre de droite est une somme comme dans (\ref{equation-formal-sum}).
Nous omettons de vérifier la compatibilité avec les applications de restriction
et la fonctorialité en $\mathcal{F}$.

\medskip\noindent
Pour achever la démonstration, affirmons que, pour tout point géométrique
$\overline{y} : \Spec(k) \to Y$, le diagramme suivant est commutatif :
$$
\xymatrix{
(j_{p!}\mathcal{F})_{\overline{y}} \ar[r] \ar[d] &
\bigoplus_{j(\overline{x}) = \overline{y}} \mathcal{F}_{\overline{x}}
\ar@{=}[d] \\
(f_{p!}\mathcal{F})_{\overline{y}} \ar[r] &
\bigoplus_{f(\overline{x}) = \overline{y}} \mathcal{F}_{\overline{x}}
}
$$
où la flèche horizontale supérieure est construite dans la démonstration de la
proposition de Cohomologie étale
\ref{etale-cohomology-proposition-describe-jshriek},
la flèche horizontale inférieure dans celle du
lemme \ref{lemma-finite-support-stalk},
la flèche verticale droite est l'égalité évidente, et
la flèche verticale gauche est l'application induite sur les fibres par celle définie au
paragraphe précédent. L'affirmation résulte immédiatement
de la description explicite de toutes les flèches qui interviennent
ici et dans les références indiquées.
Puisque les flèches horizontales sont des isomorphismes,
la flèche verticale gauche est elle aussi un isomorphisme. Notre application
induit donc un isomorphisme après passage aux faisceaux associés, d'après le
théorème \ref{etale-cohomology-theorem-exactness-stalks} de Cohomologie étale.
\end{proof}

\begin{definition}
\label{definition-f-shriek-lqf}
Soit $f : X \to Y$ un morphisme de schémas localement quasi-fini.
Nous appelons {\it image directe à support propre} le
foncteur
$$
f_! : \textit{Ab}(X_\etale) \longrightarrow \textit{Ab}(Y_\etale)
$$
défini par la formule $f_!\mathcal{F} = (f_{p!}\mathcal{F})^\#$,
c'est-à-dire que $f_!\mathcal{F}$ est le faisceau associé au préfaisceau
$f_{p!}\mathcal{F}$ construit ci-dessus.
\end{definition}

\noindent
D'après le lemme \ref{lemma-finite-support-f-shriek-separated},
ceci ne contredit pas la définition \ref{definition-f-shriek-separated}
(lorsque les deux définitions s'appliquent) et, d'après le
lemme \ref{lemma-finite-support-etale-shriek},
ceci ne contredit pas non plus la
définition \ref{etale-cohomology-definition-extension-zero} de Cohomologie étale
(lorsque les deux définitions s'appliquent).

\begin{lemma}
\label{lemma-lqf-f-shriek-stalk}
Soit $f : X \to Y$ un morphisme de schémas localement quasi-fini. Alors :
\begin{enumerate}
\item pour $\mathcal{F}$ dans $\textit{Ab}(X_\etale)$ et tout point
géométrique $\overline{y} : \Spec(k) \to Y$, on a
$$
(f_!\mathcal{F})_{\overline{y}} =
\bigoplus\nolimits_{f(\overline{x}) = \overline{y}} \mathcal{F}_{\overline{x}}
$$
fonctoriellement en $\mathcal{F}$ ;
\item le foncteur $f_! : \textit{Ab}(X_\etale) \to \textit{Ab}(Y_\etale)$
est exact et commute aux sommes directes.
\end{enumerate}
\end{lemma}

\begin{proof}
La formule pour les fibres découle immédiatement du
lemme \ref{lemma-finite-support-stalk} (et lui est en fait équivalente).
L'exactitude du foncteur en résulte aussitôt,
puisque l'exactitude se vérifie sur les fibres ; voir le
théorème \ref{etale-cohomology-theorem-exactness-stalks} de Cohomologie étale.
\end{proof}

\begin{remark}[Covariance par rapport aux immersions ouvertes]
\label{remark-covariance-lqf-f-shriek}
Soit $f : X \to Y$ un morphisme de schémas localement quasi-fini. Soit
$\mathcal{F}$ un faisceau abélien sur $X_\etale$.
Soit $X' \subset X$ un sous-schéma ouvert, et notons $f' : X' \to Y$
la restriction de $f$.
Il existe une application canonique
$$
f'_!(\mathcal{F}|_{X'}) \longrightarrow f_!\mathcal{F}
$$
Cette application est la faisceautisation d'une application canonique
$$
f'_{p!}(\mathcal{F}|_{X'}) \to f_{p!}\mathcal{F}
$$
construite comme suit. Soit $V \in Y_\etale$, et considérons une section
$s' = \sum_{i = 1, \ldots, n} (Z'_i, s'_i)$ comme dans
(\ref{equation-formal-sum}), définissant un élément de
$f'_{p!}(\mathcal{F}|_{X'})(V)$.
Alors $Z'_i \subset X'_V$ peut aussi être regardé comme un sous-schéma localement fermé
de $X_V$, et l'on a $H_{Z'_i}(\mathcal{F}|_{X'}) = H_{Z'_i}(\mathcal{F})$.
Nous envoyons $s'$ sur exactement la même somme
$s = \sum_{i = 1, \ldots, n} (Z'_i, s'_i)$,
mais regardée maintenant comme un élément de $f_{p!}\mathcal{F}(V)$.
Nous omettons de vérifier que cette construction est compatible avec les
applications de restriction et fonctorielle en $\mathcal{F}$.
Cette construction possède les propriétés suivantes :
\begin{enumerate}
\item Les applications $f'_{p!}\mathcal{F}' \to f_{p!}\mathcal{F}$ et
$f'_!\mathcal{F}' \to f_!\mathcal{F}$ sont compatibles avec
la description des fibres donnée dans les lemmes
\ref{lemma-finite-support-stalk} et \ref{lemma-lqf-f-shriek-stalk}.
\item Si $f$ est séparé, l'application
$f'_{p!}\mathcal{F}' \to f_{p!}\mathcal{F}$ est celle
construite dans la remarque \ref{remark-covariance-f-shriek-separated}
via l'isomorphisme du lemme \ref{lemma-finite-support-f-shriek-separated}.
\item Si $X'' \subset X'$ est un autre ouvert, la composée de
$f''_{p!}(\mathcal{F}|_{X''}) \to f'_{p!}(\mathcal{F}|_{X'}) \to
f_{p!}\mathcal{F}$ est l'application
$f''_{p!}(\mathcal{F}|_{X''}) \to f_{p!}\mathcal{F}$ associée à
l'inclusion $X'' \subset X$. Après passage aux faisceaux associés, il en va
de même de
$f''_!(\mathcal{F}|_{X''}) \to f'_!(\mathcal{F}|_{X'}) \to f_!\mathcal{F}$.
\item L'application $f'_!\mathcal{F}' \to f_!\mathcal{F}$ est injective,
car cela se vérifie sur les fibres.
\end{enumerate}
Tous ces énoncés se démontrent aisément en représentant les éléments
par des sommes finies comme ci-dessus et en examinant leurs images.
\end{remark}

\begin{lemma}
\label{lemma-lqf-colimit-f-shriek}
Soit $f : X \to Y$ un morphisme de schémas localement quasi-fini.
Soit $X = \bigcup_{i \in I} X_i$ un recouvrement ouvert. Il existe alors
un complexe exact
$$
\ldots \to
\bigoplus\nolimits_{i_0, i_1, i_2} f_{i_0i_1i_2, !}
\mathcal{F}|_{X_{i_0i_1i_2}} \to
\bigoplus\nolimits_{i_0, i_1} f_{i_0i_1, !} \mathcal{F}|_{X_{i_0i_1}} \to
\bigoplus\nolimits_{i_0} f_{i_0, !} \mathcal{F}|_{X_{i_0}}
\to f_!\mathcal{F} \to 0
$$
fonctoriel en $\mathcal{F} \in \textit{Ab}(X_\etale)$ ; voir les détails
dans la démonstration.
\end{lemma}

\begin{proof}
Comme d'habitude, posons $X_{i_0 \ldots i_p} = X_{i_0} \cap \ldots \cap X_{i_p}$
et notons $f_{i_0 \ldots i_p}$ la restriction de $f$ à
$X_{i_0 \ldots i_p}$. Les différentielles du complexe sont les applications
construites dans la remarque \ref{remark-covariance-lqf-f-shriek},
avec les règles de signes du complexe de {\v C}ech.
L'exactitude découle aisément de la description des fibres donnée au
lemme \ref{lemma-lqf-f-shriek-stalk}. Nous omettons les détails.
\end{proof}

\begin{remark}[Autre construction]
\label{remark-alternative-lqf-f-shriek}
Le lemme \ref{lemma-lqf-colimit-f-shriek}
donne une autre construction du foncteur $f_!$
pour les morphismes localement quasi-finis $f$.
En effet, étant donné un morphisme de schémas localement quasi-fini $f : X \to Y$,
nous pouvons choisir un recouvrement ouvert $X = \bigcup_{i \in I} X_i$
tel que chaque $f_i : X_i \to Y$ soit séparé ; il suffit par exemple de choisir
un recouvrement ouvert affine de $X$. Nous
pouvons alors définir $f_!\mathcal{F}$ comme le conoyau de l'avant-dernière flèche
du complexe du lemme, c'est-à-dire
$$
f_!\mathcal{F} = \Coker\left(
\bigoplus\nolimits_{i_0, i_1} f_{i_0i_1, !} \mathcal{F}|_{X_{i_0i_1}} \to
\bigoplus\nolimits_{i_0} f_{i_0, !} \mathcal{F}|_{X_{i_0}}
\right)
$$
où l'on peut utiliser les constructions de $f_{i_0, !}$ et de
$f_{i_0i_1, !}$ de la section \ref{section-compact-support},
puisque les morphismes $f_{i_0}$ et $f_{i_0 i_1}$ sont séparés.
On peut alors calculer les fibres de $f_!$ (en se ramenant au cas séparé,
c'est-à-dire au lemme \ref{lemma-lqf-f-shriek-separated-colimits})
et obtenir le résultat du lemme \ref{lemma-lqf-f-shriek-stalk}.
Tous les autres résultats de la présente section s'en
déduisent également.
\end{remark}

\begin{remark}
\label{remark-construct-map-presheaves-downstairs}
Soit $g : Y' \to Y$ un morphisme de schémas.
Pour tout préfaisceau abélien $\mathcal{G}'$ sur $Y'_\etale$, notons
$g_*\mathcal{G}'$ le préfaisceau $V \mapsto \mathcal{G}'(Y' \times_Y V)$.
Si $\alpha : \mathcal{G} \to g_*\mathcal{G}'$ est une application de préfaisceaux abéliens
sur $Y_\etale$, il existe une unique application
$\alpha^\# : \mathcal{G}^\# \to g_*((\mathcal{G}')^\#)$
de faisceaux abéliens sur $Y_\etale$ telle que le diagramme
$$
\xymatrix{
\mathcal{G} \ar[d] \ar[r]_\alpha & g_*\mathcal{G}' \ar[d] \\
\mathcal{G}^\# \ar[r]^-{\alpha^\#} & g_*((\mathcal{G}')^\#)
}
$$
soit commutatif, les flèches verticales provenant des applications canoniques
$\mathcal{G} \to \mathcal{G}^\#$ et $\mathcal{G}' \to (\mathcal{G}')^\#$. Si
$\alpha' : g^{-1}\mathcal{G}^\# \to (\mathcal{G}')^\#$
est l'application adjointe à $\alpha^\#$, alors, pour tout point géométrique
$\overline{y}' : \Spec(k) \to Y'$ d'image
$\overline{y} = g \circ \overline{y}'$ dans $Y$, l'application
$$
\alpha'_{\overline{y}'} :
\mathcal{G}_{\overline{y}} =
(\mathcal{G}^\#)_{\overline{y}} =
(g^{-1}\mathcal{G}^\#)_{\overline{y}'}
\longrightarrow
(\mathcal{G}')^\#_{\overline{y}'} =
\mathcal{G}'_{\overline{y}'}
$$
est définie comme suit : une section $s$ de $\mathcal{G}$ sur un voisinage étale
$(V, \overline{v})$ définit une classe dans la fibre, envoyée sur celle de la section
$\alpha(s)$ de $g_*\mathcal{G}'(V) = \mathcal{G}'(Y' \times_Y V)$
sur le voisinage étale $(Y' \times_Y V, (\overline{y}', \overline{v}))$
dans la fibre de $\mathcal{G}'$ en $\overline{y}'$.
\end{remark}

\begin{lemma}
\label{lemma-lqf-base-change-f-shriek}
Considérons un carré cartésien
$$
\xymatrix{
X' \ar[r]_{g'} \ar[d]_{f'} & X \ar[d]^f \\
Y' \ar[r]^g & Y
}
$$
de schémas, où $f$ est localement quasi-fini. Il existe un isomorphisme
$g^{-1}f_!\mathcal{F} \to f'_!(g')^{-1}\mathcal{F}$, fonctoriel en
$\mathcal{F}$ dans $\textit{Ab}(X_\etale)$, qui est compatible avec
la description des fibres donnée au lemme \ref{lemma-lqf-f-shriek-stalk}
(voir la démonstration pour l'énoncé précis).
\end{lemma}

\begin{proof}
Avec les conventions de la remarque \ref{remark-construct-map-presheaves-downstairs},
nous allons construire explicitement une application
$$
c : f_{p!}\mathcal{F} \longrightarrow g_*f'_{p!}(g')^{-1}\mathcal{F}
$$
de préfaisceaux abéliens sur $Y_\etale$. D'après la discussion de la
remarque \ref{remark-construct-map-presheaves-downstairs},
elle détermine une application canonique
$g^{-1}f_!\mathcal{F} \to f'_!(g')^{-1}\mathcal{F}$. Enfin, nous
montrerons que cette application induit des isomorphismes sur les fibres et conclurons par le
théorème \ref{etale-cohomology-theorem-exactness-stalks} de Cohomologie étale.

\medskip\noindent
Construction de l'application $c$. Soit $V \in Y_\etale$, et considérons une section
$s = \sum_{i = 1, \ldots, n} (Z_i, s_i)$ comme dans
(\ref{equation-formal-sum}), définissant un élément de $f_{p!}\mathcal{F}(V)$.
La valeur de $g_*f'_{p!}(g')^{-1}\mathcal{F}$ en $V$ est
$f'_{p!}(g')^{-1}\mathcal{F}(V')$, où $V' = V \times_Y Y'$.
Notons $Z'_i \subset X'_{V'}$ le changement de base de $Z_i$ à $V'$.
D'après (\ref{item-pullback}), il existe une application de changement de base
$H_{Z_i}(\mathcal{F}) \to H_{Z'_i}((g')^{-1}\mathcal{F})$.
Si $s'_i \in H_{Z'_i}((g')^{-1}\mathcal{F})$ désigne l'image de $s_i$
par changement de base, posons $c(s) = \sum_{i = 1, \ldots, n} (Z'_i, s'_i)$ comme dans
(\ref{equation-formal-sum}), ce qui définit un élément de
$f'_{p!}(g')^{-1}\mathcal{F}(V')$. Nous omettons de vérifier
que cette construction est compatible avec les relations
(\ref{item-sum}) et (\ref{item-sub}), ainsi qu'avec les
applications de restriction. La construction est manifestement
fonctorielle en $\mathcal{F}$.

\medskip\noindent
Soit $\overline{y}' : \Spec(k) \to Y'$ un point géométrique d'image
$\overline{y} = g \circ \overline{y}'$ dans $Y$. Remarquons que
$X'_{\overline{y}'} = X_{\overline{y}}$ par transitivité des
produits fibrés. Ainsi, $g'$ induit une bijection
$\{f'(\overline{x}') = \overline{y}'\} \to \{f(\overline{x}) = \overline{y}\}$
et, si $\overline{x}'$ s'envoie sur $\overline{x}$, alors
$((g')^{-1}\mathcal{F})_{\overline{x}'} = \mathcal{F}_{\overline{x}}$
d'après le lemme \ref{etale-cohomology-lemma-stalk-pullback} de Cohomologie étale.
Affirmons maintenant que le diagramme
$$
\xymatrix{
(g^{-1}f_!\mathcal{F})_{\overline{y}'} \ar@{=}[r] \ar[d] &
(f_!\mathcal{F})_{\overline{y}} \ar[r] \ar[ld] &
\bigoplus\nolimits_{f(\overline{x}) = \overline{y}}
\mathcal{F}_{\overline{x}} \ar[d]
\\
(f'_!(g')^{-1}\mathcal{F})_{\overline{y}'} \ar[rr] & &
\bigoplus\nolimits_{f'(\overline{x}') = \overline{y}'}
(g')^{-1}\mathcal{F}_{\overline{x}'}
}
$$
est commutatif, les flèches horizontales étant celles de la démonstration du
lemme \ref{lemma-finite-support-stalk}, tandis que la flèche verticale droite
est une égalité d'après ce qui précède. La flèche sud-ouest est
décrite dans la remarque \ref{remark-construct-map-presheaves-downstairs}
comme l'application de changement de base, c'est-à-dire qu'elle est
simplement donnée par notre construction $c$ ci-dessus. La description
élémentaire de l'image d'une somme $\sum (Z_i, z_i)$ dans la
fibre en $\overline{x}$, donnée dans la démonstration du
lemme \ref{lemma-finite-support-stalk}, montre immédiatement que le
diagramme commute. Ceci achève la démonstration du lemme.
\end{proof}

\begin{lemma}
\label{lemma-lqf-separated-shriek-composition}
Soient $f' : X \to Y'$ et $g : Y' \to Y$ des morphismes composables de schémas,
où $f'$ et $f = g \circ f'$ sont localement quasi-finis, et $g$ séparé et
localement de type fini. Il existe alors un isomorphisme canonique de foncteurs
$g_! \circ f'_! = f_!$. Cet isomorphisme est compatible avec
\begin{enumerate}
\item[(a)] la covariance par rapport aux immersions ouvertes des
remarques \ref{remark-covariance-f-shriek-separated} et
\ref{remark-covariance-lqf-f-shriek},
\item[(b)] les isomorphismes de changement de base des
lemmes \ref{lemma-lqf-base-change-f-shriek}
et \ref{lemma-base-change-f-shriek-separated} ;
\item[(c)] il coïncide avec l'isomorphisme du lemme \ref{lemma-f-shriek-composition}
via les identifications du lemme \ref{lemma-finite-support-f-shriek-separated}
lorsque $f'$ est séparé.
\end{enumerate}
\end{lemma}

\begin{proof}
Soit $\mathcal{F}$ un faisceau abélien sur $X_\etale$. Avec les conventions de la
remarque \ref{remark-construct-map-presheaves-downstairs}, nous allons
construire explicitement une application
$$
c : f_{p!}\mathcal{F} \longrightarrow g_*f'_{p!}\mathcal{F}
$$
de préfaisceaux abéliens sur $Y_\etale$. D'après la discussion de la
remarque \ref{remark-construct-map-presheaves-downstairs},
elle détermine une application canonique
$c^\# : f_!\mathcal{F} \to g_*f'_!\mathcal{F}$.
Nous montrerons que l'image de $c^\#$ est contenue dans le sous-faisceau
$g_!f'_!\mathcal{F}$, obtenant ainsi une application
$c' : f_!\mathcal{F} \to g_!f'_!\mathcal{F}$. Nous démontrerons ensuite
(a), (b) et (c). Enfin, le point (b)
permettra de montrer que $c'$ est un isomorphisme.

\medskip\noindent
Construction de l'application $c$. Soit $V \in Y_\etale$, et
soit $s = \sum (Z_i, s_i)$ une somme comme dans (\ref{equation-formal-sum}),
définissant un élément de $f_{p!}\mathcal{F}(V)$.
Rappelons que $Z_i \subset X_V = X \times_Y V$
est un sous-schéma localement fermé fini sur $V$.
En posant $V' = Y' \times_Y V$, on obtient $X_{V'} = X \times_{Y'} V' = X_V$.
Ainsi $Z_i \subset X_{V'}$ est localement fermé, et
$Z_i$ est fini sur $V'$ puisque $g$ est séparé
(lemme \ref{morphisms-lemma-finite-permanence} de Morphismes).
Nous pouvons donc poser $c(s) = \sum (Z_i, s_i)$, mais en la regardant maintenant
comme un élément de $f'_{p!}\mathcal{F}(V') = (g_*f'_{p!}\mathcal{F})(V)$.
La construction est manifestement compatible avec les relations
(\ref{item-sum}) et (\ref{item-sub}),
ainsi qu'avec les applications de restriction ; nous obtenons donc l'application $c$.

\medskip\noindent
Remarquons que, dans la discussion précédente, notre section $c(s) = \sum (Z_i, s_i)$ de
$f'_!\mathcal{F}$ sur $V'$ a une restriction nulle sur
$V' \setminus \Im(\coprod Z_i \to V')$. Puisque $\Im(\coprod Z_i \to V')$
est propre sur $V$ (par exemple d'après le lemme de Morphismes
\ref{morphisms-lemma-scheme-theoretic-image-is-proper}),
nous concluons que $c(s)$ définit une section de
$g_!f'_!\mathcal{F} \subset g_*f'_!\mathcal{F}$ sur $V$.
Puisque toute section locale de $f_!\mathcal{F}$ provient localement d'une
section locale de $f_{p!}\mathcal{F}$, nous concluons que l'image
de $c^\#$ est contenue dans $g_!f'_!\mathcal{F}$.
Nous obtenons donc une application induite $c' : f_!\mathcal{F} \to g_!f'_!\mathcal{F}$
par laquelle $c^\#$ se factorise, comme annoncé au premier paragraphe de la démonstration.

\medskip\noindent
Démonstration de (a). Soit $Y'_1 \subset Y'$ un sous-schéma ouvert,
et posons $X_1 = (f')^{-1}(W')$. Nous obtenons un diagramme
$$
\xymatrix{
X_1 \ar[d]_{f'_1} \ar[r]_a \ar@/_2em/[dd]_{f_1} &
X \ar[d]^{f'} \ar@/^2em/[dd]^f \\
Y'_1 \ar[d]_{g_1} \ar[r]_{b'} &
Y' \ar[d]^g \\
Y \ar@{=}[r] &
Y
}
$$
où les flèches horizontales sont des immersions ouvertes. Nous affirmons alors que
le diagramme
$$
\xymatrix{
f_{1, !}\mathcal{F}|_{X_1} \ar[r]_{c'_1} \ar[dd] &
g_{1, !}f'_{1, !}\mathcal{F}|_{X_1} \ar@{=}[d] \\
& g_{1, !}(f'_!\mathcal{F})|_{Y'_1} \ar[d] \\
f_!\mathcal{F} \ar[r]^{c'} &
g_!f'_!\mathcal{F} \ar[r] & g_*f'_!\mathcal{F}
}
$$
commute, la flèche verticale gauche étant celle de la
remarque \ref{remark-covariance-lqf-f-shriek}, et
la flèche verticale droite celle de la remarque \ref{remark-covariance-f-shriek-separated}.
Le signe d'égalité du diagramme vient de ce que $f'_1$
est la restriction de $f'$ à $Y'_1$ et que notre construction
de $f'_!$ est locale sur la base.
Enfin, pour démontrer la commutativité, choisissons un objet $V$ de
$Y_\etale$ et une somme formelle $s_1 = \sum (Z_{1, i}, s_{1, i})$ comme dans
(\ref{equation-formal-sum}), définissant un élément de
$f_{1, p!}\mathcal{F}|_{X_1}(V)$. Cela signifie que
$Z_{1, i} \subset X_1 \times_Y V$ est localement fermé et fini sur $V$,
et que $s_{1, i} \in H_{Z_{1, i}}(\mathcal{F})$.
Suivons cette section
le long des applications qui interviennent : il suffit de montrer que nous
obtenons le même élément de
$g_*f'_!\mathcal{F}(V) = f'_!\mathcal{F}(Y' \times_Y V)$.
En parcourant chacun des deux côtés du diagramme, on voit immédiatement
que l'on obtient l'élément $\sum (Z_{1, i}, s_{1, i})$,
où $Z_{1, i}$ est maintenant regardé comme un sous-schéma localement fermé
de $X \times_{Y'} (Y' \times_Y V) = X \times_Y V$, fini sur
$Y' \times_Y V$.

\medskip\noindent
Démonstration de (b). Soit $b : Y_1 \to Y$ un morphisme de schémas. Formons le
diagramme commutatif
$$
\xymatrix{
X_1 \ar[d]_{f'_1} \ar[r]_a \ar@/_2em/[dd]_{f_1} &
X \ar[d]^{f'} \ar@/^2em/[dd]^f \\
Y'_1 \ar[d]_{g_1} \ar[r]_{b'} &
Y' \ar[d]^g \\
Y_1 \ar[r]^b &
Y
}
$$
dont les carrés sont cartésiens. Nous affirmons que notre construction est compatible
avec les applications de changement de base des lemmes \ref{lemma-lqf-base-change-f-shriek}
et \ref{lemma-base-change-f-shriek-separated}, c'est-à-dire
que le rectangle supérieur du diagramme
$$
\xymatrix{
b^{-1}f_!\mathcal{F} \ar[rr] \ar[d]_{b^{-1}c'} & &
f_{1, !}a^{-1}\mathcal{F} \ar[d]^{c_1'} \\
b^{-1}g_!f'_!\mathcal{F} \ar[r] \ar[d] &
g_{1, !}(b')^{-1}f'_!\mathcal{F} \ar[r] \ar[d] &
g_{1, !}f'_{1, !}a^{-1}\mathcal{F} \ar[d] \\
b^{-1}g_*f'_!\mathcal{F} \ar[r] &
g_{1, *}(b')^{-1}f'_!\mathcal{F} \ar[r] &
g_{1, *}f'_{1, !}a^{-1}\mathcal{F}
}
$$
commute. Cette vérification est élémentaire, et nous
invitons le lecteur à l'omettre. Comme les flèches de la ligne médiane
vers la ligne inférieure sont injectives, il suffit de montrer que
le contour extérieur du diagramme commute.
Pour cela, il suffit de prendre une section locale de
$b^{-1}f_!\mathcal{F}$ et de montrer que les deux chemins donnent la même
section locale de $g_{1, *}f'_{1, !}a^{-1}\mathcal{F}$.
Il suffit même de le vérifier
pour les sections locales qui sont les images inverses par $b$
d'une section $s = \sum (Z_i, s_i)$ de $f_{p!}\mathcal{F}(V)$
comme ci-dessus (puisque de telles images inverses engendrent le faisceau abélien
$b^{-1}f_!\mathcal{F}$). Notons $V_1$, $V'_1$ et $Z_{1, i}$
les changements de base de $V$, $V' = Y' \times_Y V$ et $Z_i$ par $Y_1 \to Y$.
Rappelons que $Z_i$ est un sous-schéma localement fermé de $X_V = X_{V'}$ ;
ainsi $Z_{1, i}$ est un sous-schéma localement fermé
de $(X_1)_{V_1} = (X_1)_{V'_1}$. Alors $b^{-1}c'$ envoie l'image inverse
de $s$ sur l'image inverse de la section locale $c(s) \sum (Z_i, s_i)$, regardée
comme un élément de $f'_{p!}\mathcal{F}(V') = (g_*f'_{p!}\mathcal{F})(V)$.
La composée des deux applications de changement de base inférieures
l'envoie simplement sur $\sum (Z_{i, 1}, s_{1, i})$, regardé comme un
élément de $f'_{1, p!}a^{-1}\mathcal{F}(V'_1) =
g_{1, *}f'_{1, p!}a^{-1}\mathcal{F}(V_1)$.
D'autre part, l'application de changement de base en haut du diagramme
envoie l'image inverse de $s$ sur $\sum (Z_{1, i}, s_{1, i})$, regardé
comme un élément de $f_{1, !}a^{-1}\mathcal{F}(V_1)$.
Enfin, par sa construction même, $c'_1$
envoie bien cet élément sur $\sum (Z_{i, 1}, s_{1, i})$, regardé comme un
élément de $f'_{1, p!}a^{-1}\mathcal{F}(V'_1) =
g_{1, *}f'_{1, p!}a^{-1}\mathcal{F}(V_1)$ ; la commutativité
est ainsi vérifiée.

\medskip\noindent
Démonstration de (c). Elle résulte de la comparaison des définitions
des deux applications ; nous omettons les détails.

\medskip\noindent
Pour achever la démonstration, il suffit de montrer que l'image inverse de
$c'$ par tout point géométrique $\overline{y} : \Spec(k) \to Y$
est un isomorphisme. En effet, prendre l'image inverse par $\overline{y}$
revient à prendre la fibre en $\overline{y}$
(remarque \ref{etale-cohomology-remark-stalk-pullback} de Cohomologie étale),
et nous pouvons donc invoquer le
théorème \ref{etale-cohomology-theorem-exactness-stalks} de Cohomologie étale.
Grâce à la compatibilité (b) que nous venons d'établir, nous 
concluons que nous pouvons supposer que $Y$ est le spectre de $k$,
et il reste à montrer que $c'$ est un isomorphisme.
Il suffit pour cela de montrer que l'application induite
$$
\bigoplus\nolimits_{x \in X} \mathcal{F}_x = H^0(Y, f_!\mathcal{F})
\longrightarrow
H^0(Y, g_!f'_!\mathcal{F}) = H^0_c(Y', f'_!\mathcal{F})
$$
est un isomorphisme. Les égalités résultent des
lemmes \ref{lemma-lqf-f-shriek-stalk} et
\ref{lemma-stalk-f-shriek-separated}.
Rappelons que $X$ est une somme disjointe de
spectres d'anneaux locaux artiniens de corps résiduel $k$ ; voir le
lemme \ref{varieties-lemma-algebraic-scheme-dim-0} de Variétés.
Comme les deux membres commutent aux sommes directes
(nous omettons les détails), nous pouvons supposer que $\mathcal{F}$ est un faisceau gratte-ciel
$x_*A$ concentré en un certain point $x \in X$.
Alors $f'_!\mathcal{F}$ est le faisceau gratte-ciel au point
$y'$, image de $x$ dans $Y$, d'après le lemme \ref{lemma-lqf-f-shriek-stalk}.
Dans ce cas, il est manifeste que notre construction
donne l'application identité $A \to H^0_c(Y', y'_*A) = A$,
comme souhaité.
\end{proof}

\begin{lemma}
\label{lemma-lqf-shriek-composition}
Soient $f : X \to Y$ et $g : Y \to Z$ des morphismes de schémas
localement quasi-finis et composables. Il existe alors un isomorphisme canonique de foncteurs
$$
(g \circ f)_! \longrightarrow g_! \circ f_!
$$
Ces isomorphismes vérifient les propriétés suivantes :
\begin{enumerate}
\item Si $f$ et $g$ sont séparés, l'isomorphisme coïncide avec celui du
lemme \ref{lemma-f-shriek-composition}.
\item Si $g$ est séparé, l'isomorphisme coïncide avec celui du
lemme \ref{lemma-lqf-separated-shriek-composition}.
\item Pour tout point géométrique $\overline{z} : \Spec(k) \to Z$, le diagramme
$$
\xymatrix{
((g \circ f)_!\mathcal{F})_{\overline{z}} \ar[d] \ar[rr] & &
\bigoplus\nolimits_{g(f(\overline{x})) = \overline{z}}
\mathcal{F}_{\overline{x}} \ar@{=}[d] \\
(g_!f_!\mathcal{F})_{\overline{z}} \ar[r] &
\bigoplus\nolimits_{g(\overline{y}) = \overline{z}}
(f_!\mathcal{F})_{\overline{y}} \ar[r] &
\bigoplus\nolimits_{g(f(\overline{x})) = \overline{z}}
\mathcal{F}_{\overline{x}}
}
$$
est commutatif, les flèches horizontales étant données par le
lemme \ref{lemma-lqf-f-shriek-stalk}.
\item Soit $h : Z \to T$ un troisième morphisme de schémas
localement quasi-fini. Alors le diagramme
$$
\xymatrix{
(h \circ g \circ f)_! \ar[r] \ar[d] &
(h \circ g)_! \circ f_! \ar[d] \\
h_! \circ (g \circ f)_! \ar[r] &
h_! \circ g_! \circ f_!
}
$$
est commutatif.
\item Supposons donné un diagramme de schémas
$$
\xymatrix{
X' \ar[d]_{f'} \ar[r]_c & X \ar[d]^f \\
Y' \ar[d]_{g'} \ar[r]_b & Y \ar[d]^g \\
Z' \ar[r]^a & Z
}
$$
dont les deux carrés sont cartésiens et où $f$ et $g$ sont
localement quasi-finis. Alors le diagramme
$$
\xymatrix{
a^{-1} \circ (g \circ f)_! \ar[d] \ar[rr] & &
(g' \circ f')_! \circ c^{-1} \ar[d] \\
a^{-1} \circ g_! \circ f_! \ar[r] &
g'_! \circ b^{-1} \circ f_! \ar[r] &
g'_! \circ f'_! \circ c^{-1}
}
$$
est commutatif, les flèches horizontales étant celles du
lemme \ref{lemma-lqf-base-change-f-shriek}.
\end{enumerate}
\end{lemma}

\begin{proof}
Si $f$ et $g$ sont séparés, il s'agit d'un cas particulier du
lemme \ref{lemma-f-shriek-composition}.
Si $g$ est séparé, il s'agit d'un cas particulier du
lemme \ref{lemma-lqf-separated-shriek-composition},
qui coïncide en outre avec le cas où $f$ et $g$ sont séparés.

\medskip\noindent
Construction dans le cas général. Choisissons un recouvrement ouvert $Y = \bigcup Y_i$
tel que la restriction $g_i : Y_i \to Z$ de $g$ soit séparée.
Posons $X_i = f^{-1}(Y_i)$ et notons $f_i : X_i \to Y_i$ la restriction
de $f$. Notons aussi $h = g \circ f$ et $h_i : X_i \to Z$ la restriction
de $h$. Considérons le diagramme suivant :
$$
\xymatrix{
\bigoplus\nolimits_{i_0, i_1}
h_{i_0i_1, !}\mathcal{F}|_{X_{i_0i_1}} \ar[r] \ar[d] &
\bigoplus\nolimits_{i_0} h_{i_0, !}\mathcal{F}|_{X_{i_0}} \ar[r] \ar[d] &
h_!\mathcal{F} \ar[r] \ar@{..>}[dd] &
0 \\
\bigoplus\nolimits_{i_0, i_1}
g_{i_0i_1, !} f_{i_0i_1, !}\mathcal{F}|_{X_{i_0i_1}} \ar[r] \ar[d] &
\bigoplus\nolimits_{i_0}
g_{i_0, !} f_{i_0, !}\mathcal{F}|_{X_{i_0}} \ar[d] \\
\bigoplus\nolimits_{i_0, i_1}
g_{i_0i_1, !} (f_!\mathcal{F})|_{Y_{i_0i_1}} \ar[r] &
\bigoplus\nolimits_{i_0}
g_{i_0, !} (f_!\mathcal{F})|_{Y_{i_0}} \ar[r] &
g_!f_!\mathcal{F} \ar[r] &
0
}
$$
D'après le lemme \ref{lemma-lqf-colimit-f-shriek}, les lignes supérieure et inférieure
du diagramme sont exactes. D'après le lemme \ref{lemma-lqf-separated-shriek-composition},
le carré supérieur gauche commute. Les flèches verticales du
carré inférieur gauche proviennent des égalités
$(f_!\mathcal{F})|_{Y_{i_0i_1}} = f_{i_0i_1, !}\mathcal{F}|_{X_{i_0i_1}}$ et
$(f_!\mathcal{F})|_{Y_{i_0}} = f_{i_0, !}\mathcal{F}|_{X_{i_0}}$
puisque la construction de $f_!$ est locale sur la base. De plus, ces
égalités sont bien entendu compatibles avec les identifications
$((f_!\mathcal{F})|_{Y_{i_0}})|_{Y_{i_0i_1}} =
(f_!\mathcal{F})|_{Y_{i_0i_1}}$ et
$(f_{i_0, !}\mathcal{F}|_{X_{i_0}})|_{Y_{i_0i_1}} =
f_{i_0i_1, !}\mathcal{F}|_{X_{i_0i_1}}$
qui servent (avec la covariance par rapport aux immersions ouvertes
pour $Y_{i_0i_1} \subset Y_{i_0}$)
à définir les flèches horizontales du carré inférieur gauche.
Ce carré commute donc lui aussi.
Nous en concluons qu'il existe une unique
flèche pointillée comme indiqué dans le diagramme et
que cette flèche est en outre un isomorphisme.

\medskip\noindent
Démonstration des propriétés (1) -- (5). Fixons le recouvrement ouvert $Y = \bigcup Y_i$.
Remarquons que, si $Y \to Z$ est séparé, nous obtenons une flèche pointillée
s'insérant dans le grand diagramme ci-dessus grâce à l'application du
lemme \ref{lemma-lqf-separated-shriek-composition}
(par les propriétés mêmes de ce lemme).
Ceci démontre (2), puis (1) par compatibilité des
applications du lemme \ref{lemma-lqf-separated-shriek-composition}
et du lemme \ref{lemma-f-shriek-composition}.
Ensuite, pour tout schéma $Z'$ sur $Z$, nous obtenons la compatibilité de (5)
pour l'application $(g' \circ f')_! \to g'_! \circ f'_!$
construite au moyen du recouvrement ouvert $Y' = \bigcup b^{-1}(Y_i)$.
Cela résulte clairement de la compatibilité correspondante des applications
construites au lemme \ref{lemma-lqf-separated-shriek-composition}.
En particulier, nous pouvons considérer un point géométrique
$\overline{z} : \Spec(k) \to Z$. Comme les applications
$X_{\overline{z}} \to Y_{\overline{z}} \to \Spec(k)$
sont séparées, le changement de base de
$(g \circ f)_!\mathcal{F} \to g_! f_! \mathcal{F}$
par $\overline{z}$ coïncide avec l'application du
lemme \ref{lemma-f-shriek-composition}.
Le lecteur voit alors immédiatement que l'on obtient la propriété (3).
La propriété (3) garantit bien sûr que notre transformation de foncteurs
$(g \circ f)_! \to g_! \circ f_!$, construite au moyen du recouvrement ouvert
$Y = \bigcup Y_i$, ne dépend pas du choix de ce recouvrement.
Enfin, la propriété (4) se déduit de la propriété (3), déjà démontrée,
en examinant ce qui se passe sur les fibres.
\end{proof}








\section{Pondérations et morphismes trace pour les morphismes localement quasi-finis}
\sectionmark{Pondérations et morphismes trace}
\label{section-weightings}

\noindent
Une référence pour cette section est
\cite[Exposé XVII, Proposition 6.2.5]{SGA4}.

\medskip\noindent
Soit $f : X \to Y$ un morphisme de schémas localement quasi-fini.
Soit $w : X \to \mathbf{Z}$ une pondération de $f$ ; voir la
définition \ref{more-morphisms-definition-weighting} de Compléments sur les morphismes.
Soit $\mathcal{F}$ un faisceau abélien sur $Y_\etale$.
Dans cette section, nous montrerons qu'il existe un morphisme
$$
\text{Tr}_{f, w, \mathcal{F}} :
f_!f^{-1}\mathcal{F}
\longrightarrow
\mathcal{F}
$$
de faisceaux abéliens sur $Y_\etale$, caractérisé par la propriété suivante :
sur les fibres en un point géométrique $\overline{y}$ de $Y$, on obtient l'application
$$
\bigoplus\nolimits_{f(\overline{x}) = \overline{y}} w(\overline{x}) :
(f_!f^{-1}\mathcal{F})_{\overline{y}} =
\bigoplus\nolimits_{f(\overline{x}) = \overline{y}}
\mathcal{F}_{\overline{y}}
\longrightarrow
\mathcal{F}_{\overline{y}}
$$
Comme indiqué, la flèche est donnée par multiplication par l'entier
$w(\overline{x})$ sur le facteur correspondant à $\overline{x}$.
L'égalité située à gauche de la flèche résulte du
lemme \ref{lemma-lqf-f-shriek-stalk}, combiné au
lemme \ref{etale-cohomology-lemma-stalk-pullback} de Cohomologie étale.

\medskip\noindent
Si le morphisme $f : X \to Y$ est plat, localement quasi-fini et localement de
présentation finie, il existe une pondération canonique,
et l'on obtient un morphisme trace canonique dont la formation est
compatible au changement de base ; voir
l'exemple \ref{example-trace-for-flat-quasi-finite}.
Si $Y$ est un schéma localement noethérien unibranche et si $f : X \to Y$
est localement quasi-fini, on peut également définir une pondération
(naturelle) de $f$ et l'on dispose aussi de morphismes trace dans ce cas ; voir
l'exemple \ref{example-trace-for-quasi-finite-over-normal}.

\begin{lemma}
\label{lemma-f-shriek-projection}
Soit $f : X \to Y$ un morphisme de schémas localement quasi-fini.
Soit $\Lambda$ un anneau.
Soit $\mathcal{F}$ un faisceau de $\Lambda$-modules sur $X_\etale$
et soit $\mathcal{G}$ un faisceau de $\Lambda$-modules sur $Y_\etale$.
Il existe un isomorphisme canonique
$$
can :
f_!\mathcal{F} \otimes_\Lambda \mathcal{G}
\longrightarrow
f_!(\mathcal{F} \otimes_\Lambda f^{-1}\mathcal{G})
$$
de faisceaux de $\Lambda$-modules sur $Y_\etale$.
\end{lemma}

\begin{proof}
Rappelons que $f_!\mathcal{F} = (f_{p!}\mathcal{F})^\#$
d'après la définition \ref{definition-f-shriek-lqf}, où
$f_{p!}\mathcal{F}$ est le préfaisceau
construit à la section \ref{section-finite-support}.
Pour construire la flèche, il suffit donc de construire une application
$$
f_{p!}\mathcal{F} \otimes_{p, \Lambda} \mathcal{G}
\longrightarrow
f_{p!}(\mathcal{F} \otimes_\Lambda f^{-1}\mathcal{G})
$$
de préfaisceaux sur $Y_\etale$. Ici, le symbole $\otimes_{p, \Lambda}$
désigne le produit tensoriel de préfaisceaux ; voir la
section \ref{sites-modules-section-tensor-product} de Modules sur les sites.
Soit $V$ un objet de $Y_\etale$. Rappelons que
$$
f_{p!}\mathcal{F}(V) = \colim_Z H_Z(\mathcal{F})
\quad\text{et}\quad
f_{p!}(\mathcal{F} \otimes_\Lambda f^{-1}\mathcal{G})(V) =
\colim_Z H_Z(\mathcal{F} \otimes_\Lambda f^{-1}\mathcal{G})
$$
Voir la section \ref{section-finite-support}. Notre application est définie
sur les tenseurs purs par la règle
$$
(Z, s) \otimes t \longmapsto (Z, s \otimes f^{-1}t)
$$
(voir ci-dessous pour les notations), puis prolongée par linéarité à tout
$(f_{p!}\mathcal{F} \otimes_{p, \Lambda} \mathcal{G})(V) =
f_{p!}\mathcal{F}(V) \otimes_\Lambda \mathcal{G}(V)$.
Les notations employées ici sont les suivantes :
\begin{enumerate}
\item $Z \subset X_V$ est un sous-schéma localement fermé fini sur $V$ ;
\item $s \in H_Z(\mathcal{F})$, ce qui signifie que $s \in \mathcal{F}(U)$
avec $\text{Supp}(s) \subset Z$ pour un ouvert $U \subset X_V$ tel que
$Z \subset U$ soit fermé ;
\item $t \in \mathcal{G}(V)$, d'image $f^{-1}t \in f^{-1}\mathcal{G}(U)$.
\end{enumerate}
Puisque le support de $s \in \mathcal{F}(U)$ est contenu dans $Z$, il est clair
que le support de $s \otimes f^{-1}t$ est lui aussi contenu dans $Z$.
Le couple $(Z, s \otimes f^{-1}t)$ a donc bien un sens.
Il est immédiat que la construction commute aux applications de transition
de la limite inductive $\colim_Z H_Z(\mathcal{F})$ et qu'elle est
compatible avec les applications de restriction. Enfin, il est tout aussi clair
que la construction est compatible avec les identifications des
fibres de $f_!$ données au lemme \ref{lemma-lqf-f-shriek-stalk}.
Autrement dit, l'application $can$ que nous avons construite sur les fibres en un point géométrique
$\overline{y}$ s'insère dans le diagramme commutatif
$$
\xymatrix{
(f_!\mathcal{F} \otimes_\Lambda \mathcal{G})_{\overline{y}}
\ar[r]_-{can_{\overline{y}}} \ar[d] &
f_!(\mathcal{F} \otimes_\Lambda f^{-1}\mathcal{G})_{\overline{y}} \ar[d] \\
(\bigoplus \mathcal{F}_{\overline{x}})
\otimes_\Lambda \mathcal{G}_{\overline{y}} \ar[r] &
\bigoplus
(\mathcal{F}_{\overline{x}} \otimes_\Lambda \mathcal{G}_{\overline{y}})
}
$$
où les sommes directes portent sur les points géométriques $\overline{x}$
au-dessus de $\overline{y}$, les flèches verticales sont les identifications
du lemme \ref{lemma-lqf-f-shriek-stalk}, et la flèche horizontale inférieure
est l'isomorphisme évident.
Nous concluons que $can$ est un isomorphisme, comme souhaité.
\end{proof}

\begin{lemma}
\label{lemma-trace-map-exists}
Soit $f : X \to Y$ un morphisme de schémas localement quasi-fini.
Soit $w : X \to \mathbf{Z}$ une pondération de $f$. Pour tout faisceau abélien
$\mathcal{F}$ sur $Y$, il existe un unique morphisme trace
$\text{Tr}_{f, w, \mathcal{F}} : f_!f^{-1}\mathcal{F} \to \mathcal{F}$
ayant le comportement prescrit sur les fibres.
\end{lemma}

\begin{proof}
D'après le lemme \ref{lemma-f-shriek-projection}, nous avons une identification
$f_!f^{-1}\mathcal{F} = f_!\underline{\mathbf{Z}} \otimes \mathcal{F}$
compatible avec la description des fibres de ces faisceaux aux
points géométriques. Il suffit donc de construire le morphisme
$$
\text{Tr}_{f, w, \underline{\mathbf{Z}}} :
f_!\underline{\mathbf{Z}}
\longrightarrow
\underline{\mathbf{Z}}
$$
ayant le comportement prescrit sur les fibres. D'après la
définition \ref{definition-f-shriek-lqf}, on a
$f_!\underline{\mathbf{Z}} = (f_{p!}\underline{\mathbf{Z}})^\#$,
où $f_{p!}\underline{\mathbf{Z}}$ est le préfaisceau
construit à la section \ref{section-finite-support}.
Il suffit donc de construire une application
$$
f_{p!}\underline{\mathbf{Z}} \longrightarrow \underline{\mathbf{Z}}
$$
de préfaisceaux sur $Y_\etale$. Soit $V$ un objet de $Y_\etale$.
Rappelons, d'après la section \ref{section-finite-support}, que
$$
f_{p!}\underline{\mathbf{Z}}(V) = \colim_Z H_Z(\underline{\mathbf{Z}})
$$
Ici, la limite inductive est prise sur l'ensemble ordonné des sous-schémas
localement fermés $Z \subset X_V$ qui sont finis sur $V$. Pour chaque tel
$Z$, nous allons définir une application
$$
H_Z(\underline{\mathbf{Z}}) \longrightarrow \underline{\mathbf{Z}}(V)
$$
compatible avec les applications qui définissent la limite inductive.

\medskip\noindent
Soit $Z \subset X_V$ localement fermé et fini sur $V$.
Choisissons un ouvert $U \subset X_V$ contenant $Z$ comme partie fermée.
Un élément $s$ de $H_Z(\underline{\mathbf{Z}})$ est une section
$s \in \underline{\mathbf{Z}}(U)$ dont le support est contenu dans $Z$.
Soit $U_n \subset U$ la partie ouverte et fermée où la valeur
de $s$ est $n \in \mathbf{Z}$. La condition sur le support donne
$Z \cap U_n = U_n$ pour $n \not = 0$. Ainsi, pour $n \not = 0$, l'ouvert
$U_n$ est aussi fermé dans $Z$ (comme complémentaire de tous les autres),
et $U_n \to V$ est donc fini puisque $Z$ est fini sur $V$.
Par définition même d'une pondération, la
fonction $\int_{U_n \to V} w|_{U_n}$ est alors localement constante sur $V$,
et nous pouvons la regarder comme un élément de $\underline{\mathbf{Z}}(V)$.
Notre construction envoie $(Z, s)$ sur l'élément
$$
\sum\nolimits_{n \in \mathbf{Z},\ n \not = 0}
n \left(\int_{U_n \to V} w|_{U_n}\right)
\quad \in \quad \underline{\mathbf{Z}}(V)
$$
La somme est localement finie sur $V$, donc bien définie ; nous omettons les détails
(dans toute la discussion, le lecteur peut commencer par choisir des ouverts affines et
s'assurer que tous les schémas qui interviennent dans l'argument sont quasi-compacts, afin que
la somme soit finie). Nous omettons de vérifier que cette construction
est compatible avec les applications de la limite inductive et avec les applications de
restriction qui définissent $f_{p!}\underline{\mathbf{Z}}$.

\medskip\noindent
Soit $\overline{y}$ un point géométrique de $Y$ au-dessus du point $y \in Y$.
En prenant les fibres en $\overline{y}$, la construction précédente détermine une application
$$
(f_!\underline{\mathbf{Z}})_{\overline{y}}
= \bigoplus\nolimits_{f(\overline{x}) = \overline{y}} \mathbf{Z}
\longrightarrow
\mathbf{Z} = \underline{\mathbf{Z}}_{\overline{y}}
$$
Pour achever la démonstration, montrons que cette application est donnée par multiplication par
$w(\overline{x})$ sur le facteur correspondant à $\overline{x}$.
Choisissons en effet $\overline{x}$ au-dessus de $\overline{y}$.
Il existe un voisinage étale $(V, \overline{v}) \to (Y, \overline{y})$
tel que $X_V$ contienne un ouvert $U$ fini sur $V$,
tel que $\overline{x}$ soit le seul à appartenir à $U$ parmi les points géométriques
de $X$ relevant $\overline{y}$.
Ceci résulte du lemme de Compléments sur les morphismes
\ref{more-morphisms-lemma-etale-makes-quasi-finite-finite-multiple-points-var};
nous omettons certains détails. Alors $(U, 1)$ définit une section de
$f_!\underline{\mathbf{Z}}$ sur $V$, dont l'image vaut $1$ dans le facteur
correspondant à $\overline{x}$ et zéro dans les autres facteurs
(voir la démonstration du lemme \ref{lemma-finite-support-stalk}), et
notre construction précédente envoie $(U, 1)$ sur $\int_{U \to V} w|_U$,
qui est constante de valeur $w(\overline{x})$ dans un voisinage
de $\overline{v}$, comme souhaité.
\end{proof}

\begin{lemma}
\label{lemma-properties-trace-map}
Soit $f : X \to Y$ un morphisme de schémas localement quasi-fini.
Soit $w : X \to \mathbf{Z}$ une pondération de $f$. Les morphismes trace
construits ci-dessus vérifient les propriétés suivantes :
\begin{enumerate}
\item $\text{Tr}_{f, w, \mathcal{F}}$ est fonctoriel en $\mathcal{F}$ ;
\item $\text{Tr}_{f, w, \mathcal{F}}$ est compatible à tout changement de base ;
\item étant donnés un anneau $\Lambda$ et $K$ dans $D(Y_\etale, \Lambda)$,
on obtient $\text{Tr}_{f, w, K} : f_!f^{-1}K \to K$, fonctoriel en $K$
et compatible à tout changement de base.
\end{enumerate}
\end{lemma}

\begin{proof}
Le point (1) résulte soit de la construction du morphisme trace
dans la démonstration du lemme \ref{lemma-trace-map-exists},
soit, plus simplement, de ce que la caractérisation de l'application
impose cette fonctorialité sur toutes les fibres.
Soit
$$
\xymatrix{
X' \ar[r]_{g'} \ar[d]_{f'} & X \ar[d]^f \\
Y' \ar[r]^g & Y
}
$$
un diagramme cartésien de schémas. Alors la fonction
$w' = w \circ g' : X' \to \mathbf{Z}$ est une pondération de $f'$ d'après le
lemme \ref{more-morphisms-lemma-weighting-base-change} de Compléments sur les morphismes.
Le point (2) signifie que le diagramme
$$
\xymatrix{
g^{-1}f_!f^{-1}\mathcal{F}
\ar[rr]_-{g^{-1}\text{Tr}_{f, w, \mathcal{F}}} \ar@{=}[d] & &
g^{-1}\mathcal{F} \ar@{=}[d] \\
f'_!(f')^{-1}g^{-1}\mathcal{F}
\ar[rr]^-{\text{Tr}_{f', w', g^{-1}\mathcal{F}}} & &
g^{-1}\mathcal{F}
}
$$
est commutatif, l'égalité verticale gauche étant donnée par
$$
g^{-1}f_!f^{-1}\mathcal{F} =
f'_!(g')^{-1}f^{-1}\mathcal{F} =
f'_!(f')^{-1}g^{-1}\mathcal{F}
$$
où le premier signe d'égalité est donné par le lemme \ref{lemma-lqf-base-change-f-shriek}
(changement de base pour l'image directe à support propre). La commutativité de ce diagramme
résulte de la caractérisation de l'action de nos morphismes trace
sur les fibres et du fait que la flèche de changement de base du
lemme \ref{lemma-lqf-base-change-f-shriek} respecte la description des fibres.

\medskip\noindent
Les points (1) et (2) étant acquis, le point (3) résulte de ce que les foncteurs
$f^{-1} : D(Y_\etale, \Lambda) \to D(X_\etale, \Lambda)$ et
$f_! : D(X_\etale, \Lambda) \to D(Y_\etale, \Lambda)$
s'obtiennent en appliquant $f^{-1}$ et $f_!$ à des complexes quelconques
de modules représentant les objets considérés.
\end{proof}

\begin{lemma}
\label{lemma-trace-composition}
Soient $f : X \to Y$ et $g : Y \to Z$ des morphismes localement quasi-finis.
Soit $w_f : X \to \mathbf{Z}$ une pondération de $f$, et soit
$w_g : Y \to \mathbf{Z}$ une pondération de $g$. Pour
$K \in D(Z_\etale, \Lambda)$, la composée
$$
(g \circ f)_!(g \circ f)^{-1}K =
g_! f_! f^{-1} g^{-1}K
\xrightarrow{g_! \text{Tr}_{f, w_f, g^{-1}K}}
g_!g^{-1}K
\xrightarrow{\text{Tr}_{g, w_g, K}}
K
$$
est égale à $\text{Tr}_{g \circ f, w_{g \circ f}, K}$, où
$w_{g \circ f}(x) = w_f(x) w_g(f(x))$.
\end{lemma}

\begin{proof}
On a $(g \circ f)_! = g_! \circ f_!$ d'après le
lemme \ref{lemma-lqf-shriek-composition}.
Le lemme \ref{more-morphisms-lemma-weighting-composition} de Compléments sur les morphismes
montre que $w_{g \circ f}$ est une pondération de $g \circ f$ ;
l'énoncé a donc un sens. L'égalité se vérifie sur les fibres.
Nous omettons les détails.
\end{proof}

\begin{example}[Morphisme trace dans le cas plat quasi-fini]
\label{example-trace-for-flat-quasi-finite}
Soit $f : X \to Y$ un morphisme de schémas plat,
localement quasi-fini et localement de présentation finie.
On obtient une pondération positive canonique $w : X \to \mathbf{Z}$
en posant
$$
w(x) = \text{longueur}_{\mathcal{O}_{X, x}}
(\mathcal{O}_{X, x}/\mathfrak m_{f(x)} \mathcal{O}_{X, x})
[\kappa(x) : \kappa(f(x))]_i
$$
Voir le lemme de Compléments sur les morphismes
\ref{more-morphisms-lemma-weighting-flat-quasi-finite}.
Ainsi, d'après les lemmes \ref{lemma-trace-map-exists} et
\ref{lemma-properties-trace-map}, on obtient pour $f$ des morphismes trace
$$
\text{Tr}_{f, K} : f_!f^{-1}K \longrightarrow K
$$
fonctoriels en $K$ dans $D(Y_\etale, \Lambda)$ et compatibles
à tout changement de base. Notons que tout changement de base
$f' : X' \to Y'$ de $f$ possède les mêmes propriétés et que la restriction de $w$
est la pondération canonique de $f'$.
\end{example}

\begin{remark}
\label{remark-trace-etale-counit}
Soit $j : U \to X$ un morphisme étale de schémas. Alors le morphisme trace
$\text{Tr} : j_!j^{-1}K \to K$ de
l'exemple \ref{example-trace-for-flat-quasi-finite} est égal à la
flèche d'adjonction associée à l'adjonction entre $j_!$ et $j^{-1}$.
Nous avons déjà employé le terme « trace » pour cette flèche d'adjonction dans la
section \ref{etale-cohomology-section-trace-method} de Cohomologie étale.
\end{remark}

\begin{example}[Morphisme trace dans le cas quasi-fini sur un schéma normal]
\label{example-trace-for-quasi-finite-over-normal}
Soit $Y$ un schéma géométriquement unibranche et localement noethérien ;
par exemple, $Y$ peut être une variété normale. Soit $f : X \to Y$ un morphisme de schémas
localement quasi-fini. Il existe alors une
pondération positive $w : X \to \mathbf{Z}$ de $f$,
définie approximativement en associant à $x$ le
« degré séparable générique »
de $\mathcal{O}_{X, x}^{sh}$ sur $\mathcal{O}_{Y, f(x)}^{sh}$.
Voir le lemme de Compléments sur les morphismes
\ref{more-morphisms-lemma-weighting-quasi-finite-Noetherian}.
Ainsi, d'après les lemmes \ref{lemma-trace-map-exists} et
\ref{lemma-properties-trace-map}, on obtient pour $f$ et $w$ des morphismes trace
$$
\text{Tr}_{f, w, K} : f_!f^{-1}K \longrightarrow K
$$
fonctoriels en $K$ dans $D(Y_\etale, \Lambda)$ et compatibles
à tout changement de base. Toutefois, dans ce cas, pour un changement de base
$f' : X' \to Y'$ de $f$, la restriction de $w$ à $X'$ ne possède pas en général
d'interprétation « naturelle » en termes du
morphisme $f'$.
\end{example}
















\section{Image inverse extraordinaire pour les morphismes localement quasi-finis}
\sectionmark{Image inverse extraordinaire}
\label{section-duality-locally-quasi-finite}

\noindent
Pour tout morphisme de schémas localement quasi-fini $f : X \to Y$, le
foncteur $f_! : \textit{Ab}(X_\etale) \to \textit{Ab}(Y_\etale)$ commute
aux sommes directes et est exact ; voir le lemme \ref{lemma-lqf-f-shriek-stalk}.
Ceci suggère qu'il admet un adjoint à droite, que nous noterons $f^!$.

\medskip\noindent
Avertissement : ce foncteur est la version non dérivée !

\begin{lemma}
\label{lemma-lqf-f-upper-shriek}
Soit $f : X \to Y$ un morphisme de schémas localement quasi-fini.
\begin{enumerate}
\item Le foncteur $f_! : \textit{Ab}(X_\etale) \to \textit{Ab}(Y_\etale)$
admet un adjoint à droite $f^! : \textit{Ab}(Y_\etale) \to \textit{Ab}(X_\etale)$.
\item On a
$f^!(\overline{y}_*A) = \prod_{f(\overline{x}) = \overline{y}} \overline{x}_*A$.
\item Si $\Lambda$ est un anneau, le foncteur
$f_! : \textit{Mod}(X_\etale, \Lambda) \to \textit{Mod}(Y_\etale, \Lambda)$
admet un adjoint à droite
$f^! : \textit{Mod}(Y_\etale, \Lambda) \to \textit{Mod}(X_\etale, \Lambda)$
qui coïncide avec $f^!$ sur les faisceaux abéliens sous-jacents.
\end{enumerate}
\end{lemma}

\begin{proof}
Démonstration de (1). Soit $E \subset \Ob(\textit{Ab}(Y_\etale))$ la classe
formée des produits de faisceaux gratte-ciel. Nous affirmons ce qui suit :
\begin{enumerate}
\item[(a)] tout $\mathcal{G}$ de $\textit{Ab}(Y_\etale)$ est un sous-faisceau
d'un élément de $E$ ;
\item[(b)] pour tout $\mathcal{G} \in E$, il existe un objet
$\mathcal{H}$ de $\textit{Ab}(X_\etale)$ tel que
$\Hom(f_!\mathcal{F}, \mathcal{G}) = \Hom(\mathcal{F}, \mathcal{H})$
fonctoriellement en $\mathcal{F}$.
\end{enumerate}
Une fois cette affirmation vérifiée, le dual du
lemme \ref{homology-lemma-partially-defined-adjoint} d'Homologie
fournit le foncteur adjoint $f^!$.

\medskip\noindent
Le point (a) est vrai, car on peut envoyer $\mathcal{G}$ dans le faisceau
$\prod \overline{y}_*\mathcal{G}_{\overline{y}}$, où le
produit porte sur tous les points géométriques de $Y$. C'est une injection d'après le
théorème \ref{etale-cohomology-theorem-exactness-stalks} de Cohomologie étale.
(Il s'agit de la première étape de la résolution de Godement lorsque
l'on travaille avec des faisceaux abéliens sur des espaces topologiques.)

\medskip\noindent
Le point (b) et le point (2) du lemme s'établissent comme suit.
Supposons que $\mathcal{G} = \prod \overline{y}_*A_{\overline{y}}$
pour certains groupes abéliens $A_{\overline{y}}$. Alors
$$
\Hom(f_!\mathcal{F}, \mathcal{G}) =
\prod \Hom(f_!\mathcal{F}, \overline{y}_*A_{\overline{y}})
$$
Il suffit donc de trouver des faisceaux abéliens $\mathcal{H}_{\overline{y}}$
sur $X_\etale$ qui représentent les foncteurs
$\mathcal{F} \mapsto \Hom(f_!\mathcal{F}, \overline{y}_*A_{\overline{y}})$
et de prendre $\mathcal{H} = \prod \mathcal{H}_{\overline{y}}$.
Nous sommes ainsi ramenés au cas $\mathcal{H} = \overline{y}_*A$
pour un point géométrique fixé $\overline{y} : \Spec(k) \to Y$
et un groupe abélien fixé $A$. Nous affirmons que, dans ce cas,
$\mathcal{H} = \prod_{f(\overline{x}) = \overline{y}} \overline{x}_*A$ convient.
Ceci achèvera la démonstration des points (1) et (2) du lemme.
En effet, on a
$$
\Hom(f_!\mathcal{F}, \overline{y}_*A) =
\Hom_{\textit{Ab}}((f_!\mathcal{F})_{\overline{y}}, A) =
\Hom_{\textit{Ab}}(\bigoplus\nolimits_{f(\overline{x}) = \overline{y}}
\mathcal{F}_{\overline{x}}, A)
$$
d'une part, d'après la description des fibres donnée au
lemme \ref{lemma-lqf-f-shriek-stalk},
et, d'autre part, on a
$$
\Hom(\mathcal{F}, \mathcal{H}) =
\prod\nolimits_{f(\overline{x}) = \overline{y}}
\Hom(\mathcal{F}, \overline{x}_*A) =
\prod\nolimits_{f(\overline{x}) = \overline{y}}
\Hom_{\textit{Ab}}(\mathcal{F}_{\overline{x}}, A)
$$
Nous laissons au lecteur le soin d'identifier ces deux expressions comme foncteurs en $\mathcal{F}$.

\medskip\noindent
Démonstration du point (3). Remarquons qu'un objet de $\textit{Mod}(X_\etale, \Lambda)$
n'est autre qu'un objet $\mathcal{F}$ de $\textit{Ab}(X_\etale)$
muni d'une application $\Lambda \to \text{End}(\mathcal{F})$. Les
foncteurs $f_!$ et $f^!$ de (1) définissent les foncteurs $f_!$ et $f^!$ comme en (3).
Un calcul immédiat montre qu'ils sont adjoints.
\end{proof}

\begin{lemma}
\label{lemma-etale-upper-shriek}
Soit $j : U \to X$ un morphisme étale. Alors $j^! = j^{-1}$.
\end{lemma}

\begin{proof}
Cela résulte du fait que le foncteur $j_!$ défini à la
section \ref{section-finite-support}
coïncide avec $j_!$ tel qu'il est défini dans Cohomologie étale, à la section
\ref{etale-cohomology-section-extension-by-zero} ; voir le
lemme \ref{lemma-finite-support-etale-shriek}. Enfin, dans la
section \ref{etale-cohomology-section-extension-by-zero} de Cohomologie étale,
le foncteur $j_!$ est défini
comme l'adjoint à gauche de $j^{-1}$ ; on conclut donc par
unicité des foncteurs adjoints.
\end{proof}

\begin{lemma}
\label{lemma-upper-shriek-restriction}
Soient $f : X \to Y$ et $g : Y \to Z$ des morphismes séparés et localement
quasi-finis. Il existe un isomorphisme canonique $(g \circ f)^! \to f^! \circ g^!$.
Étant donné un troisième morphisme localement quasi-fini $h : Z \to T$,
le diagramme
$$
\xymatrix{
(h \circ g \circ f)^! \ar[r] \ar[d] &
f^! \circ (h \circ g)^! \ar[d] \\
(g \circ f)^! \circ h^! \ar[r] & f^! \circ g^! \circ h^!
}
$$
est commutatif.
\end{lemma}

\begin{proof}
Par unicité des foncteurs adjoints, ceci se traduit immédiatement
par l'énoncé (dual) correspondant pour les foncteurs $f_!$.
Voir le lemme \ref{lemma-lqf-shriek-composition}.
\end{proof}

\begin{lemma}
\label{lemma-upper-shriek-restriction-etale}
Soient $j : U \to X$ et $j' : V \to U$ des morphismes étales.
L'isomorphisme $(j \circ j')^{-1} = (j')^{-1} \circ j^{-1}$
et l'isomorphisme $(j \circ j')^! = (j')^! \circ j^!$ du
lemme \ref{lemma-upper-shriek-restriction}
coïncident via l'isomorphisme du lemme \ref{lemma-etale-upper-shriek}.
\end{lemma}

\begin{proof}
La démonstration est omise.
\end{proof}

\begin{lemma}
\label{lemma-lqf-base-change-upper-shriek}
Considérons un carré cartésien
$$
\xymatrix{
X' \ar[r]_{g'} \ar[d]_{f'} & X \ar[d]^f \\
Y' \ar[r]^g & Y
}
$$
de schémas, où $f$ est localement quasi-fini. Pour tout faisceau abélien $\mathcal{F}$
sur $Y'_\etale$, on a $(g')_*(f')^!\mathcal{F} = f^!g_*\mathcal{F}$.
\end{lemma}

\begin{proof}
Par unicité des foncteurs adjoints, ceci résulte de
l'énoncé (dual) correspondant pour les foncteurs $f_!$.
Voir le lemme \ref{lemma-lqf-base-change-f-shriek}.
\end{proof}

\begin{remark}
\label{remark-pointed-sets}
Les résultats de cette section se généralisent aux faisceaux d'ensembles pointés.
En effet, pour un site $\mathcal{C}$, notons $\Sh^*(\mathcal{C})$ la catégorie des
faisceaux d'ensembles pointés. Les constructions de la présente section et de la précédente
s'appliquent, mutatis mutandis, aux faisceaux d'ensembles pointés. Ainsi, étant donné un morphisme
de schémas localement quasi-fini $f : X \to Y$, on obtient
un couple de foncteurs adjoints
$$
f_! : \Sh^*(X_\etale) \longrightarrow \Sh^*(Y_\etale)
\quad\text{et}\quad
f^! : \Sh^*(Y_\etale) \longrightarrow \Sh^*(X_\etale)
$$
tel que, pour tout point géométrique $\overline{y}$ de $Y$, il existe des
isomorphismes
$$
(f_!\mathcal{F})_{\overline{y}} =
\coprod\nolimits_{f(\overline{x}) = \overline{y}}
\mathcal{F}_{\overline{x}}
$$
(le coproduit étant pris dans la catégorie des ensembles pointés) fonctoriels en
$\mathcal{F} \in \Sh^*(X_\etale)$ et des isomorphismes
$$
f^!(\overline{y}_*S) =
\prod\nolimits_{f(\overline{x}) = \overline{y}}
\overline{x}_*S
$$
fonctoriels par rapport à l'ensemble pointé $S$. Si
$F : \textit{Ab}(X_\etale) \to \Sh^*(X_\etale)$ et
$F : \textit{Ab}(Y_\etale) \to \Sh^*(Y_\etale)$
désignent les foncteurs d'oubli, la compatibilité des constructions
garantit l'existence d'applications canoniques
$$
f_!F(\mathcal{F}) \longrightarrow F(f_!\mathcal{F})
$$
fonctorielles en $\mathcal{F} \in \textit{Ab}(X_\etale)$ et
$$
F(f^!\mathcal{G}) \longrightarrow f^!F(\mathcal{G})
$$
fonctorielles en $\mathcal{G} \in \textit{Ab}(Y_\etale)$,
qui induisent les applications évidentes sur les fibres, resp.\ les faisceaux gratte-ciel.
En fait, la transformation $F \circ f^! \to f^! \circ F$ est un isomorphisme
(car $f^!$ commute aux produits).
\end{remark}







\section{Image inverse extraordinaire dérivée pour les morphismes localement quasi-finis}
\sectionmark{Image inverse extraordinaire dérivée}
\label{section-derived-duality-locally-quasi-finite}

\noindent
On peut passer aux foncteurs dérivés dans la
section \ref{section-duality-locally-quasi-finite},
et obtenir l'énoncé suivant.

\begin{lemma}
\label{lemma-lqf-shriek-derived}
Soit $f : X \to Y$ un morphisme de schémas localement quasi-fini.
Soit $\Lambda$ un anneau. Les foncteurs $f_!$ et $f^!$ de la
définition \ref{definition-f-shriek-lqf} et du
lemme \ref{lemma-lqf-f-upper-shriek}
induisent des foncteurs adjoints $f_! : D(X_\etale, \Lambda) \to D(Y_\etale, \Lambda)$
et $Rf^! : D(Y_\etale, \Lambda) \to D(X_\etale, \Lambda)$
sur les catégories dérivées.
\end{lemma}

\noindent
Dans le cas séparé, le foncteur $f_!$ est défini à la
section \ref{section-compact-support}.

\begin{proof}
Ceci résulte immédiatement du
lemme \ref{derived-lemma-derived-adjoint-functors} de Catégories dérivées,
du fait que $f_!$ est exact (lemme \ref{lemma-lqf-f-shriek-stalk})
et donc que $Lf_! = f_!$,
et du fait que l'on dispose d'assez de complexes K-injectifs de $\Lambda$-modules
sur $Y_\etale$ pour que $Rf^!$ soit défini.
\end{proof}

\begin{remark}
\label{remark-lqf-shriek-derived-torsion}
Soit $f : X \to Y$ un morphisme de schémas localement quasi-fini. Soit
$\Lambda$ un anneau. Le foncteur
$f_! : D(X_\etale, \Lambda) \to D(Y_\etale, \Lambda)$ du
lemme \ref{lemma-lqf-shriek-derived}
envoie les complexes dont les faisceaux de cohomologie sont de torsion
sur des complexes dont les faisceaux de cohomologie sont de torsion.
Ceci résulte immédiatement de la description des fibres de $f_!$ ;
voir le lemme \ref{lemma-lqf-f-shriek-stalk}.
\end{remark}

\begin{lemma}
\label{lemma-mayer-vietoris-shriek}
Soit $X$ un schéma. Supposons $X = U \cup V$, avec $U$ et $V$ ouverts.
Soit $\Lambda$ un anneau. Soit $K \in D(X_\etale, \Lambda)$. Il existe
un triangle distingué
$$
j_{U \cap V!}K|_{U \cap V} \to
j_{U!}K|_U \oplus j_{V!}K|_V \to K \to 
j_{U \cap V!}K|_{U \cap V}[1]
$$
dans $D(X_\etale, \Lambda)$, avec les notations évidentes.
\end{lemma}

\begin{proof}
Comme les foncteurs de restriction et d'image directe à support propre employés
sont exacts, il suffit de montrer que, pour tout faisceau abélien $\mathcal{F}$ sur
$X_\etale$, la suite
$$
0 \to j_{U \cap V!}\mathcal{F}|_{U \cap V} \to
j_{U!}\mathcal{F}|_U \oplus j_{V!}\mathcal{F}|_V \to \mathcal{F} \to 0
$$
est exacte. Il suffit de le vérifier sur les fibres.
\end{proof}

\begin{lemma}
\label{lemma-triangle-associated-to-open}
Soit $X$ un schéma. Soit $Z \subset X$ un sous-schéma fermé, et soit
$U \subset X$ son complémentaire. Notons $i : Z \to X$ et $j : U \to X$
les morphismes d'inclusion. Soit $\Lambda$ un anneau.
Soit $K \in D(X_\etale, \Lambda)$. Il existe un triangle distingué
$$
j_!j^{-1}K \to K \to i_*i^{-1}K \to j_!j^{-1}K[1]
$$
dans $D(X_\etale, \Lambda)$.
\end{lemma}

\begin{proof}
Conséquence immédiate du lemme de Cohomologie étale
\ref{etale-cohomology-lemma-ses-associated-to-open}
et du fait que les foncteurs $j_!$, $j^{-1}$, $i_*$, $i^{-1}$
sont exacts ; leurs versions dérivées se calculent donc en
appliquant ces foncteurs à tout complexe de faisceaux représentant $K$.
\end{proof}





\section{Préliminaires à l'image directe dérivée à support propre par compactification}
\sectionmark{Préliminaires à l'image directe dérivée}
\label{section-prelim}

\noindent
Dans cette section, nous démontrons quelques lemmes sur l'existence
de certains isomorphismes naturels de foncteurs qui résultent immédiatement
du changement de base propre.

\begin{lemma}
\label{lemma-shriek-proper-and-open}
Considérons un diagramme commutatif de schémas
$$
\xymatrix{
X' \ar[r]_{g'} \ar[d]_{f'} & X \ar[d]^f \\
Y' \ar[r]^g & Y
}
$$
où $f$ et $f'$ sont propres, et $g$ et $g'$ séparés et localement quasi-finis.
Soit $\Lambda$ un anneau. Fonctoriellement en $K \in D(X'_\etale, \Lambda)$,
il existe un morphisme canonique
$$
g_!Rf'_*K \longrightarrow Rf_*(g'_!K)
$$
dans $D(Y_\etale, \Lambda)$. Cette application est un isomorphisme si
(a) $K$ est borné inférieurement et ses faisceaux de cohomologie sont de torsion, ou si
(b) $\Lambda$ est un anneau de torsion.
\end{lemma}

\begin{proof}
Représentons $K$ par un complexe K-injectif $\mathcal{J}^\bullet$ de
faisceaux de $\Lambda$-modules sur $X'_\etale$. Choisissons un quasi-isomorphisme
$g'_!\mathcal{J}^\bullet \to \mathcal{I}^\bullet$ vers un
complexe K-injectif $\mathcal{I}^\bullet$ de
faisceaux de $\Lambda$-modules sur $X_\etale$.
Nous pouvons alors considérer l'application
$$
g_!f'_*\mathcal{J}^\bullet =
g_!f'_!\mathcal{J}^\bullet =
f_!g'_!\mathcal{J}^\bullet =
f_*g'_!\mathcal{J}^\bullet \to
f_*\mathcal{I}^\bullet
$$
où la première et la troisième égalité résultent du
lemme \ref{lemma-proper-f-shriek},
et la deuxième du
lemme \ref{lemma-f-shriek-composition}, qui affirme que
$g_! \circ f'_!$ et $f_! \circ g'_!$ sont tous deux égaux à
$(g \circ f')_! = (f \circ g')_!$ comme sous-faisceaux de
$(g \circ f')_* = (f \circ g')_*$.

\medskip\noindent
Supposons que $\Lambda$ soit de torsion, c'est-à-dire que nous soyons dans le cas (b).
Avec les notations précédentes, il suffit de montrer que
$f_*g'_!\mathcal{J}^\bullet \to f_*\mathcal{I}^\bullet$
est un isomorphisme. La question est locale sur $Y$.
Nous pouvons donc supposer que la dimension
des fibres de $f$ est bornée ; voir le lemme de Morphismes
\ref{morphisms-lemma-morphism-finite-type-bounded-dimension}.
Alors $Rf_*$ est de dimension cohomologique finie ; voir le
lemme de Cohomologie étale
\ref{etale-cohomology-lemma-cohomological-dimension-proper}.
Ainsi, d'après le lemme \ref{derived-lemma-unbounded-right-derived} de Catégories dérivées,
il suffit de montrer que $R^qf_*(g'_!\mathcal{J}) = 0$ pour $q > 0$
et pour tout faisceau injectif de $\Lambda$-modules $\mathcal{J}$
sur $X'_\etale$.

\medskip\noindent
La fibre de $R^qf_*(g'_!\mathcal{J})$ en un point géométrique
$\overline{y}$ est égale à
$H^q(X_{\overline{y}}, (g'_!\mathcal{J})|_{X_{\overline{y}}})$ d'après le
lemme de Cohomologie étale
\ref{etale-cohomology-lemma-proper-base-change-stalk}.
Puisque la formation de $g'_!$ commute au changement de base
(lemme \ref{lemma-base-change-f-shriek-separated}),
ce groupe est égal à
$$
H^q(X_{\overline{y}}, g'_{\overline{y}, !}(\mathcal{J}|_{X'_{\overline{y}}}))
$$
où $g'_{\overline{y}} : X'_{\overline{y}} \to X_{\overline{y}}$
est le morphisme induit entre les fibres géométriques.
Puisque $Y' \to Y$ est localement quasi-fini,
$X'_{\overline{y}}$ est la somme disjointe des fibres
$X'_{\overline{y}'}$ aux points géométriques $\overline{y}'$ de $Y'$
au-dessus de $\overline{y}$. Notons
$g'_{\overline{y}'} : X'_{\overline{y}'} \to X_{\overline{y}}$
la restriction de $g'_{\overline{y}}$ à $X'_{\overline{y}'}$.
Le groupe de cohomologie précédent est donc égal à
$$
H^q(X_{\overline{y}},
\bigoplus\nolimits_{\overline{y}'/\overline{y}}
g'_{\overline{y}', !}(\mathcal{J}|_{X'_{\overline{y}'}}))
$$
par exemple d'après le lemme \ref{lemma-colim-f-shriek-separated} (mais cela
résulte aussi immédiatement de la définition de $g'_{\overline{y}, !}$
à la section \ref{section-compact-support}).
Comme la cohomologie étale sur $X_{\overline{y}}$
commute aux sommes directes
(théorème \ref{etale-cohomology-theorem-colimit} de Cohomologie étale),
il suffit de montrer que
$$
H^q(X_{\overline{y}}, g'_{\overline{y}', !}(\mathcal{J}|_{X'_{\overline{y}'}}))
$$
est nul. Remarquons que
$g_{\overline{y}'} : X'_{\overline{y}'} \to X_{\overline{y}}$
est un morphisme entre schémas propres sur $\overline{y}$ et est donc
lui-même propre. Comme il est aussi localement quasi-fini, on en conclut que
$g_{\overline{y}'}$ est fini. Par conséquent,
$g'_{\overline{y}', !} = g'_{\overline{y}', *} = Rg'_{\overline{y}', *}$.
La suite spectrale de Leray nous ramène à montrer que
$$
H^q(X'_{\overline{y}'}, \mathcal{J}|_{X'_{\overline{y}'}})
$$
est nul. Comme $\Lambda$ est de torsion, cela résulte du changement de base propre
(lemme de Cohomologie étale
\ref{etale-cohomology-lemma-proper-base-change-stalk})
puisque les images directes supérieures de $\mathcal{J}$ par $f'$ sont nulles.

\medskip\noindent
Démonstration dans le cas (a). Nous allons le déduire du cas (b) par des arguments standards.
Nous montrerons que l'application induite $g_! R^pf'_* K \to R^pf_*(g'_!K)$
est un isomorphisme pour tout $p \in \mathbf{Z}$. Fixons un entier
$p_0 \in \mathbf{Z}$. Soit $a$ un entier tel que $H^j(K) = 0$
pour $j < a$. Nous allons démontrer que $g_! R^pf'_* K \to R^pf_*(g'_!K)$
est un isomorphisme pour $p \leq p_0$, par récurrence descendante sur $a$. Si
$a > p_0$, les deux membres de
l'application sont nuls pour $p \leq p_0$ par annulation triviale ; voir le
lemme \ref{derived-lemma-negative-vanishing} de Catégories dérivées
(et utiliser que $g_!$ et $g'_!$ sont des foncteurs exacts).
Supposons $a \leq p_0$. Considérons le triangle distingué
$$
H^a(K)[-a] \to K \to \tau_{\geq a + 1}K
$$
Par récurrence, le résultat vaut pour $\tau_{\geq a + 1}K$.
Au paragraphe suivant, nous le démontrerons pour $H^a(K)[-a]$.
Le lemme des cinq, appliqué au morphisme entre les
suites exactes longues de faisceaux de cohomologie
associées au morphisme de triangles distingués
$$
\xymatrix{
g_! Rf'_*(H^a(K)[-a]) \ar[d] \ar[r] &
g_! Rf'_* K \ar[r] \ar[d] &
g_! Rf'_* \tau_{\geq a + 1} K \ar[d] \\
Rf_*(g'_!(H^a(K)[-a])) \ar[r] &
Rf_*(g'_!K) \ar[r] &
Rf_*(g'_!\tau_{|geq a + 1}K)
}
$$
donne alors le résultat pour $K$. Nous omettons certains détails.

\medskip\noindent
Soit $\mathcal{F}$ un faisceau abélien de torsion sur $X'_\etale$.
Pour achever la démonstration, montrons que
$g_! Rf'_*\mathcal{F} \to R^pf_*(g'_!\mathcal{F})$
est un isomorphisme pour tout $p$.
On peut écrire $\mathcal{F} = \bigcup \mathcal{F}[n]$, où
$\mathcal{F}[n] = \Ker(n : \mathcal{F} \to \mathcal{F})$.
Nous avons l'isomorphisme pour $\mathcal{F}[n]$ d'après le cas (b).
Puisque les foncteurs $g_!$, $g'_!$, $R^pf_*$, $R^pf'_*$ commutent
aux limites inductives filtrantes (cela résulte du
lemme \ref{lemma-lqf-f-shriek-separated-colimits} et du
lemme de Cohomologie étale
\ref{etale-cohomology-lemma-relative-colimit-general})
la démonstration est achevée.
\end{proof}

\begin{lemma}
\label{lemma-shriek-proper-and-open-compose}
Considérons un diagramme commutatif de schémas
$$
\xymatrix{
X' \ar[r]_k \ar[d]_{f'} & X \ar[d]^f \\
Y' \ar[r]_l \ar[d]_{g'} & Y \ar[d]^g \\
Z' \ar[r]^m & Z
}
$$
où $f$, $f'$, $g$ et $g'$ sont propres, et
$k$, $l$ et $m$ séparés et localement quasi-finis.
Alors les isomorphismes du lemme \ref{lemma-shriek-proper-and-open}
associés aux deux carrés se composent pour donner l'isomorphisme
associé au rectangle extérieur (voir la démonstration pour un énoncé précis).
\end{lemma}

\begin{proof}
L'énoncé signifie que, si l'on écrit
$R(g \circ f)_* = Rg_* \circ Rf_*$ et
$R(g' \circ f')_* = Rg'_* \circ Rf'_*$,
alors l'isomorphisme
$m_! \circ Rg'_* \circ Rf'_* \to Rg_* \circ Rf_* \circ k_!$
associé au rectangle extérieur est égal à la composée
$$
m_! \circ Rg'_* \circ Rf'_* \to
Rg_* \circ l_! \circ Rf'_* \to
Rg_* \circ Rf_* \circ k_!
$$
des deux morphismes associés aux carrés du diagramme. Pour le montrer, choisissons
un complexe K-injectif $\mathcal{J}^\bullet$ de $\Lambda$-modules
sur $X'_\etale$ et un quasi-isomorphisme
$k_!\mathcal{J}^\bullet \to \mathcal{I}^\bullet$
vers un complexe K-injectif $\mathcal{I}^\bullet$ de $\Lambda$-modules
sur $X_\etale$. La démonstration du lemme \ref{lemma-shriek-proper-and-open}
montre que le morphisme canonique
$$
a : l_!f'_*\mathcal{J}^\bullet \to f_*\mathcal{I}^\bullet
$$
est un quasi-isomorphisme, lequel fournit la
deuxième flèche après application de $Rg_*$. D'après le
lemme \ref{sites-cohomology-lemma-K-injective-flat} de Cohomologie sur les sites,
le complexe $f_*\mathcal{I}^\bullet$, resp.\ $f'_*\mathcal{J}^\bullet$,
est un complexe K-injectif de $\Lambda$-modules sur 
$Y_\etale$, resp.\ $Y'_\etale$.
(Ce recours est un peu abusif et pourrait être évité.)
En particulier, le même raisonnement montre que le morphisme canonique
$$
b : m_!g'_*f'_*\mathcal{J}^\bullet \to g_*f_*\mathcal{I}^\bullet
$$
est un quasi-isomorphisme, et ce quasi-isomorphisme représente
la première flèche. Enfin, la démonstration du lemme \ref{lemma-shriek-proper-and-open}
montre que $g_*l_!f'_!\mathcal{J}^\bullet$ représente
$Rg_*(l_!f'_*\mathcal{J}^\bullet)$, car $f'_*\mathcal{J}^\bullet$
est K-injectif. Ainsi $Rg_*(a) = g_*(a)$, et la composée
$g_*(a) \circ b$ est la flèche du lemme \ref{lemma-shriek-proper-and-open}
associée au rectangle.
\end{proof}

\begin{lemma}
\label{lemma-shriek-proper-and-open-compose-horizontal}
Considérons un diagramme commutatif de schémas
$$
\xymatrix{
X'' \ar[r]_{g'} \ar[d]_{f''} &
X' \ar[r]_g \ar[d]_{f'} &
X \ar[d]^f \\
Y'' \ar[r]^{h'} &
Y' \ar[r]^h &
Y
}
$$
où $f$, $f'$ et $f''$ sont propres, et
$g$, $g'$, $h$ et $h'$ séparés et localement quasi-finis.
Alors les isomorphismes du lemme \ref{lemma-shriek-proper-and-open}
associés aux deux carrés se composent pour donner l'isomorphisme
associé au rectangle extérieur (voir la démonstration pour un énoncé précis).
\end{lemma}

\begin{proof}
L'énoncé signifie que, si l'on écrit
$(h \circ h')_! = h_! \circ h'_!$ et
$(g \circ g')_! = g_! \circ g'_!$
au moyen des égalités du lemme \ref{lemma-f-shriek-composition},
alors l'isomorphisme
$h_! \circ h'_! \circ Rf''_* \to Rf_* \circ g_! \circ g'_!$
associé au rectangle extérieur est égal à la composée
$$
h_! \circ h'_! \circ Rf''_* \to
h_! \circ Rf'_* \circ g'_! \to
Rf_* \circ g_! \circ g'_!
$$
des deux morphismes associés aux carrés du diagramme. Pour le montrer, choisissons
un complexe K-injectif $\mathcal{I}^\bullet$ de $\Lambda$-modules
sur $X''_\etale$ et un quasi-isomorphisme
$g'_!\mathcal{I}^\bullet \to \mathcal{J}^\bullet$
vers un complexe K-injectif $\mathcal{J}^\bullet$ de $\Lambda$-modules
sur $X'_\etale$. Choisissons ensuite un quasi-isomorphisme
$g_!\mathcal{J}^\bullet \to \mathcal{K}^\bullet$
vers un complexe K-injectif $\mathcal{K}^\bullet$ de $\Lambda$-modules
sur $X_\etale$.
La démonstration du lemme \ref{lemma-shriek-proper-and-open} montre que les morphismes
canoniques
$$
h'_!f''_*\mathcal{I}^\bullet \to f'_*\mathcal{J}^\bullet
\quad\text{et}\quad
h_!f'_*\mathcal{J}^\bullet \to f_*\mathcal{K}^\bullet
$$
sont des quasi-isomorphismes, qui définissent respectivement la
première et la deuxième flèche ci-dessus. Puisque $g_!$ est un foncteur exact
(lemme \ref{lemma-lqf-f-shriek-separated-colimits}),
on voit que $g_!g'_!\mathcal{I}^\bullet \to \mathcal{K}^\bullet$
est un quasi-isomorphisme ; par conséquent, le morphisme canonique
$$
h_!h'_!f''_*\mathcal{I}^\bullet \to f_*\mathcal{K}^\bullet
$$
est un quasi-isomorphisme et représente le morphisme associé au rectangle
extérieur dans la catégorie dérivée. Il est clair que ce morphisme est la
composée des deux autres, ce qui achève la démonstration.
\end{proof}

\begin{remark}
\label{remark-going-around}
Considérons un diagramme commutatif
$$
\xymatrix{
X'' \ar[r]_{k'} \ar[d]_{f''} & X' \ar[r]_k \ar[d]_{f'} & X \ar[d]^f \\
Y'' \ar[r]^{l'} \ar[d]_{g''} & Y' \ar[r]^l \ar[d]_{g'} & Y \ar[d]^g \\
Z'' \ar[r]^{m'} & Z' \ar[r]^m & Z
}
$$
de schémas dont les flèches verticales sont propres et les flèches horizontales
séparées et localement quasi-finies.
Étiquetons les carrés du diagramme $A$, $B$, $C$, $D$
comme suit :
$$
\begin{matrix}
A & B \\
C & D
\end{matrix}
$$
Les morphismes du lemme \ref{lemma-shriek-proper-and-open}
associés aux carrés sont alors (où l'on utilise $Rf_* = f_*$, etc.)
$$
\begin{matrix}
\gamma_A : l'_! \circ f''_* \to f'_* \circ k'_! &
\gamma_B : l_! \circ f'_* \to f_* \circ k_! \\
\gamma_C : m'_! \circ g''_* \to g'_* \circ l'_! &
\gamma_D : m_! \circ g'_* \to g_* \circ l_!
\end{matrix}
$$
Pour les rectangles $2 \times 1$ et $1 \times 2$, nous disposons de quatre autres
morphismes
$$
\begin{matrix}
\gamma_{A + B} :
(l \circ l')_! \circ f''_* \to f_* \circ (k \circ k')_*  \\
\gamma_{C + D} :
(m \circ m')_! \circ g''_* \to g_* \circ (l \circ l')_! \\
\gamma_{A + C} :
m'_! \circ (g'' \circ f'')_* \to (g' \circ f')_* \circ k'_! \\
\gamma_{B + D} :
m_! \circ (g' \circ f')_* \to (g \circ f)_* \circ k_!
\end{matrix}
$$
D'après le lemme \ref{lemma-shriek-proper-and-open-compose-horizontal}, on a
$$
\gamma_{A + B} = \gamma_B \circ \gamma_A, \quad
\gamma_{C + D} = \gamma_D \circ \gamma_C
$$
et, d'après le lemme \ref{lemma-shriek-proper-and-open-compose}, on a
$$
\gamma_{A + C} = \gamma_A \circ \gamma_C, \quad
\gamma_{B + D} = \gamma_B \circ \gamma_D
$$
Il serait ici plus exact d'écrire
$\gamma_{A + B} = (\gamma_B \star \text{id}_{k'_!}) \circ
(\text{id}_{l_!} \star \gamma_A)$ avec les notations de la
section \ref{categories-section-formal-cat-cat} de Catégories
et de même pour les autres.
Cela étant, nous obtenons a priori deux transformations
$$
m_! \circ m'_! \circ g''_* \circ f''_*
\longrightarrow
g_* \circ f_* \circ k_! \circ k'_!
$$
à savoir
$$
\gamma_B \circ \gamma_D \circ \gamma_A \circ \gamma_C =
\gamma_{B + D} \circ \gamma_{A + C}
$$
et
$$
\gamma_B \circ \gamma_A \circ \gamma_D \circ \gamma_C =
\gamma_{A + B} \circ \gamma_{C + D}
$$
Le but de cette remarque est de signaler que ces transformations
sont égales. Pour le voir, il suffit en effet de montrer que
$$
\xymatrix{
m_! \circ g'_* \circ l'_! \circ f''_* \ar[r]_{\gamma_D} \ar[d]_{\gamma_A} &
g_* \circ l_! \circ l'_! \circ f''_* \ar[d]^{\gamma_A} \\
m_! \circ g'_* \circ f'_* \circ k'_! \ar[r]^{\gamma_D} &
g_* \circ l_! \circ f'_* \circ k'_!
}
$$
commute. Cela résulte du fait que les carrés $A$ et $D$ ne se rencontrent qu'en
un point ; plus précisément, cela résulte du
lemme \ref{categories-lemma-properties-2-cat-cats} de Catégories,
ou, plus simplement, de la discussion qui précède la
définition \ref{categories-definition-horizontal-composition} de Catégories.
\end{remark}

\begin{lemma}
\label{lemma-shriek-proper-and-open-base-change}
Soit $b : Y_1 \to Y$ un morphisme de schémas. Considérons un diagramme commutatif
de schémas
$$
\vcenter{
\xymatrix{
X' \ar[r]_{g'} \ar[d]_{f'} & X \ar[d]^f \\
Y' \ar[r]^g & Y
}
}
\quad\text{et soit}\quad
\vcenter{
\xymatrix{
X'_1 \ar[r]_{g'_1} \ar[d]_{f'_1} & X_1 \ar[d]^{f_1} \\
Y'_1 \ar[r]^{g_1} & Y_1
}
}
$$
son changement de base par $b$. Supposons $f$ et $f'$ propres, et
$g$ et $g'$ séparés et localement quasi-finis.
Pour un anneau $\Lambda$ et $K$ dans $D(X'_\etale, \Lambda)$,
le diagramme suivant est commutatif :
$$
\xymatrix{
b^{-1}g_!Rf'_*K \ar[d] \ar[r] &
g_{1, !}(b')^{-1}Rf'_*K \ar[r] &
g_{1, !}Rf'_{1, *}(a')^{-1}K \ar[d] \\
b^{-1}Rf_*g'_!K \ar[r] &
Rf_{1, *}a^{-1}g'_!K \ar[r] &
Rf_{1, *}g'_{1, !}(a')^{-1}K
}
$$
dans $D(Y_{1, \etale}, \Lambda)$, où $a : X_1 \to X$, $a' : X'_1 \to X'$,
$b' : Y'_1 \to Y'$ sont les projections, les flèches verticales sont celles
du lemme \ref{lemma-shriek-proper-and-open},
et les flèches horizontales sont la flèche de changement de base
(de la section
\ref{etale-cohomology-section-base-change-preliminaries} de Cohomologie étale)
et la flèche de changement de base du lemme \ref{lemma-base-change-f-shriek-separated}.
\end{lemma}

\begin{proof}
Représentons $K$ par un complexe K-injectif $\mathcal{J}^\bullet$ de
faisceaux de $\Lambda$-modules sur $X'_\etale$. Choisissons un quasi-isomorphisme
$g'_!\mathcal{J}^\bullet \to \mathcal{I}^\bullet$ vers un
complexe K-injectif $\mathcal{I}^\bullet$ de
faisceaux de $\Lambda$-modules sur $X_\etale$.
La démonstration du lemme \ref{lemma-shriek-proper-and-open}
construit $g_!Rf'_*K \to Rf_*g'_!K$ sous la forme
$$
g_!f'_*\mathcal{J}^\bullet =
g_!f'_!\mathcal{J}^\bullet =
f_!g'_!\mathcal{J}^\bullet =
f_*g'_!\mathcal{J}^\bullet \to
f_*\mathcal{I}^\bullet
$$
Choisissons un quasi-isomorphisme
$(a')^{-1}\mathcal{J}^\bullet \to \mathcal{J}_1^\bullet$
vers un complexe K-injectif $\mathcal{J}_1^\bullet$ de
faisceaux de $\Lambda$-modules sur $X'_{1, \etale}$.
On peut alors choisir un diagramme de complexes
$$
\xymatrix{
g'_{1, !}\mathcal{J}_1^\bullet \ar[rr] & &
\mathcal{I}_1^\bullet \\
g'_{1, !}(a')^{-1}\mathcal{J}^\bullet \ar[u] \ar@{=}[r] &
a^{-1}g'_!\mathcal{J}^\bullet \ar[r] &
a^{-1}\mathcal{I}^\bullet \ar[u]
}
$$
commutatif à homotopie près, où toutes les flèches sont des quasi-isomorphismes ;
l'égalité résulte du lemme \ref{lemma-proper-f-shriek},
et $\mathcal{I}_1^\bullet$ est un complexe K-injectif de faisceaux de
$\Lambda$-modules sur $X_{1, \etale}$. Le morphisme
$g_{1, !}Rf'_{1, *}(a')^{-1}K \to Rf_{1, *}g'_{1, !}(a')^{-1}K$ est donné par
$$
g_{1, !}f'_{1, *}\mathcal{J}_1^\bullet =
g_{1, !}f'_{1, !}\mathcal{J}_1^\bullet =
f_{1, !}g'_{1, !}\mathcal{J}_1^\bullet =
f_{1, *}g'_{1, !}\mathcal{J}_1^\bullet \to
f_{1, *}\mathcal{I}_1^\bullet
$$
Les identifications correspondant aux $3$ signes d'égalité dans les deux flèches
sont compatibles avec les morphismes de changement de base, c'est-à-dire que le diagramme
$$
\xymatrix{
b^{-1}g_!f'_*\mathcal{J}^\bullet \ar@{=}[d] \ar[r] &
g_{1, !}(b')^{-1}f'_*\mathcal{J}^\bullet \ar[r] &
g_{1, !}f'_{1, *}(a')^{-1}\mathcal{J}^\bullet \ar@{=}[d] \\
b^{-1}f_*g'_!\mathcal{J}^\bullet \ar[r] &
f_{1, *}a^{-1}g'_!\mathcal{J}^\bullet \ar[r] &
f_{1, *}g'_{1, !}(a')^{-1}\mathcal{J}^\bullet
}
$$
de complexes de faisceaux abéliens commute. Pour le montrer, il suffit de
vérifier la commutativité du diagramme après avoir remplacé $g_!, g_{1, !}, g'_!, g'_{1, !}$
par $g_*, g_{1, *}, g'_*, g'_{1, *}$ (car les foncteurs d'image directe à support propre
sont définis comme des sous-foncteurs des foncteurs $*$ et les
morphismes de changement de base sont définis de manière compatible avec cette inclusion ; voir la
démonstration du lemme \ref{lemma-base-change-f-shriek-separated}).
Pour ce nouveau diagramme, la commutativité résulte de la compatibilité
des morphismes de changement de base avec la composition horizontale et verticale des diagrammes ; voir les
remarques \ref{sites-remark-compose-base-change} et
\ref{sites-remark-compose-base-change-horizontal}
de Sites, de sorte que chacun des deux parcours du diagramme donne le morphisme de
changement de base de $f \circ g' = g \circ f'$ par $b$.
Puisque, bien sûr,
$$
\xymatrix{
g_{1, !}f'_{1, *}(a')^{-1}\mathcal{J}^\bullet \ar@{=}[d] \ar[r] &
g_{1, !}f'_{1, *}\mathcal{J}_1^\bullet \ar@{=}[d] \\
f_{1, *}g'_{1, !}(a')^{-1}\mathcal{J}^\bullet \ar[r] &
f_{1, *}g'_{1, !}\mathcal{J}_1^\bullet 
}
$$
commute, il suffit, pour achever la démonstration, de montrer que
$$
\xymatrix{
b^{-1}f_*g'_!\mathcal{J}^\bullet \ar[r] \ar[d] &
f_{1, *}a^{-1}g'_!\mathcal{J}^\bullet \ar[r] \ar[d] &
f_{1, *}g'_{1, !}(a')^{-1}\mathcal{J}^\bullet \ar[r] &
f_{1, *}g'_{1, !}\mathcal{J}_1^\bullet \ar[d] \\
b^{-1}f_*\mathcal{I}^\bullet \ar[r] &
f_{1, *}a^{-1}\mathcal{I}^\bullet \ar[rr] & &
f_{1, *}\mathcal{I}_1^\bullet
}
$$
commute dans la catégorie dérivée, ce qui résulte du choix des morphismes
effectué plus haut.
\end{proof}

\begin{lemma}
\label{lemma-shriek-lqf-and-proper}
Considérons un diagramme commutatif de schémas
$$
\xymatrix{
X \ar[r]_f \ar[rd]_g & Y \ar[d]^h \\
& Z
}
$$
où $f$ et $g$ sont localement quasi-finis et $h$ propre. Soit $\Lambda$ un anneau.
Fonctoriellement en $K \in D(X_\etale, \Lambda)$, il existe un morphisme canonique
$$
g_!K \longrightarrow Rh_*(f_!K)
$$
dans $D(Z_\etale, \Lambda)$. Ce morphisme est un isomorphisme si
(a) $K$ est borné inférieurement et ses faisceaux de cohomologie sont de torsion, ou si
(b) $\Lambda$ est un anneau de torsion.
\end{lemma}

\begin{proof}
C'est un cas particulier du lemme \ref{lemma-shriek-proper-and-open}
si $f$ et $g$ sont séparés. Nous recommandons au lecteur de ne pas lire la démonstration
dans le cas général,
car nous utiliserons principalement le cas où $f$ et $g$ sont séparés.

\medskip\noindent
Représentons $K$ par un complexe $\mathcal{K}^\bullet$ de faisceaux de
$\Lambda$-modules sur $X_\etale$. Choisissons un quasi-isomorphisme
$f_!\mathcal{K}^\bullet \to \mathcal{I}^\bullet$ vers un complexe K-injectif
$\mathcal{I}^\bullet$ de faisceaux de $\Lambda$-modules sur $Y_\etale$.
Considérons le morphisme
$$
g_!\mathcal{K}^\bullet =
h_!f_!\mathcal{K}^\bullet =
h_*f_!\mathcal{K}^\bullet
\longrightarrow
h_*\mathcal{I}^\bullet
$$
où les égalités résultent des
lemmes \ref{lemma-lqf-separated-shriek-composition}
et \ref{lemma-proper-f-shriek}. Ce morphisme de complexes détermine le morphisme
$g_!K \to Rh_*(f_!K)$ de l'énoncé du lemme.

\medskip\noindent
Supposons que $\Lambda$ soit un anneau de torsion, c'est-à-dire que nous sommes dans le cas (b). Pour vérifier que le morphisme
est un isomorphisme, nous pouvons travailler localement sur $Z$.
Nous pouvons donc supposer que la dimension des fibres de $h$ est bornée ;
voir le lemme de Morphismes
\ref{morphisms-lemma-morphism-finite-type-bounded-dimension}.
Alors $Rh_*$ est de dimension cohomologique finie ; voir le
lemme de Cohomologie étale
\ref{etale-cohomology-lemma-cohomological-dimension-proper}.
Ainsi, d'après le lemme \ref{derived-lemma-unbounded-right-derived} de Catégories dérivées,
si nous montrons que $R^qh_*(f_!\mathcal{F}) = 0$ pour $q > 0$
et pour tout faisceau $\mathcal{F}$ de $\Lambda$-modules sur $X_\etale$, alors
$h_*f_!\mathcal{K}^\bullet \to h_*\mathcal{I}^\bullet$ est un quasi-isomorphisme.

\medskip\noindent
Remarquons que $\mathcal{G} = f_!\mathcal{F}$ est un faisceau de $\Lambda$-modules
sur $Y$ dont les fibres ne sont non nulles qu'aux points $y \in Y$ tels que
$\kappa(y)/\kappa(h(y))$ soit une extension finie. Cela résulte de la
description des fibres de $f_!\mathcal{F}$ au
lemme \ref{lemma-lqf-f-shriek-stalk}
et du fait que $f$ et $g$ sont tous deux localement quasi-finis.
Ainsi, d'après le théorème de changement de base propre (lemme de Cohomologie étale
\ref{etale-cohomology-lemma-proper-base-change-stalk})
il suffit de montrer que $H^q(Y_{\overline{z}}, \mathcal{H}) = 0$,
où $\mathcal{H}$ est un faisceau sur le schéma propre $Y_{\overline{z}}$
sur $\kappa(\overline{z})$, dont le support est contenu dans l'ensemble
des points fermés. L'annulation voulue résulte donc du lemme de Cohomologie étale
\ref{etale-cohomology-lemma-supported-in-closed-points}.

\medskip\noindent
Le cas (a) se déduit du cas (b) par exactement le même argument que celui employé
dans la démonstration du lemme \ref{lemma-shriek-proper-and-open}
(en utilisant le lemme \ref{lemma-lqf-f-shriek-stalk} à la place du
lemme \ref{lemma-lqf-f-shriek-separated-colimits}).
\end{proof}







\section{Image directe dérivée à support propre au moyen de compactifications}
\sectionmark{Image directe dérivée par compactification}
\label{section-derived-lower-shriek-compactification}

\noindent
Soit $f : X \to Y$ un morphisme de schémas séparé et de type fini,
où $Y$ est quasi-compact et quasi-séparé. Choisissons une compactification
$j : X \to \overline{X}$ sur $Y$ ; voir le
théorème \ref{flat-theorem-nagata} de Compléments sur la platitude.
Soit $\Lambda$ un anneau. Notons $D^+_{tors}(X_\etale, \Lambda)$
la sous-catégorie triangulée strictement pleine et saturée de
$D(X_\etale, \Lambda)$ formée des objets $K$ bornés inférieurement
et dont les faisceaux de cohomologie sont de torsion. Nous considérerons le
foncteur
$$
Rf_! = R\overline{f}_* \circ j_! :
D^+_{tors}(X_\etale, \Lambda)
\longrightarrow
D^+_{tors}(Y_\etale, \Lambda)
$$
où $\overline{f} : \overline{X} \to Y$ est le morphisme structural.
Cette définition a un sens : le foncteur $j_!$ envoie $D^+_{tors}(X_\etale, \Lambda)$ dans
$D^+_{tors}(\overline{X}_\etale, \Lambda)$ d'après la
remarque \ref{remark-lqf-shriek-derived-torsion}, et $R\overline{f}_*$ envoie
$D^+_{tors}(\overline{X}_\etale, \Lambda)$ dans $D^+_{tors}(Y_\etale, \Lambda)$
d'après le lemme \ref{etale-cohomology-lemma-torsion-direct-image} de Cohomologie étale.
Si $\Lambda$ est un anneau de torsion, nous définissons alors
$$
Rf_! = R\overline{f}_* \circ j_! :
D(X_\etale, \Lambda)
\longrightarrow
D(Y_\etale, \Lambda)
$$
Voici le lemme indispensable.

\begin{lemma}
\label{lemma-shriek-well-defined}
Soit $f : X \to Y$ un morphisme séparé de type fini entre schémas quasi-compacts
et quasi-séparés. Les foncteurs $Rf_!$ construits ci-dessus sont,
à isomorphisme canonique près, indépendants du choix de la
compactification.
\end{lemma}

\begin{proof}
Nous le démontrerons pour le foncteur
$Rf_! : D(X_\etale, \Lambda) \to D(Y_\etale, \Lambda)$
quand $\Lambda$ est un anneau de torsion ; le cas du foncteur
$Rf_! : D^+_{tors}(X_\etale, \Lambda) \to D^+_{tors}(Y_\etale, \Lambda)$
se démontre exactement de la même manière.

\medskip\noindent
Considérons la catégorie des compactifications de $X$ sur $Y$, qui est
cofiltrante d'après le théorème \ref{flat-theorem-nagata} et les
lemmes \ref{flat-lemma-compactifications-cofiltered} et
\ref{flat-lemma-compactifyable} de Compléments sur la platitude.
À toute compactification
$$
j : X \to \overline{X},\quad \overline{f} : \overline{X} \to Y
$$
la construction précédente associe le foncteur $R\overline{f}_* \circ j_! :
D(X_\etale, \Lambda) \to D(Y_\etale, \Lambda)$.
Précisons un peu. Étant donné un complexe $\mathcal{K}^\bullet$
de faisceaux de $\Lambda$-modules sur $X_\etale$, choisissons un quasi-isomorphisme
$j_!\mathcal{K}^\bullet \to \mathcal{I}^\bullet$ vers un complexe K-injectif
de faisceaux de $\Lambda$-modules sur $\overline{X}_\etale$.
Notre foncteur envoie alors $\mathcal{K}^\bullet$ sur
$\overline{f}_*\mathcal{I}^\bullet$.

\medskip\noindent
Supposons donné un morphisme $g : \overline{X}_1 \to \overline{X}_2$
entre les compactifications $j_i : X \to \overline{X}_i$ sur $Y$.
On obtient alors un isomorphisme
$$
R\overline{f}_{2, *} \circ j_{2, !} =
R\overline{f}_{2, *} \circ Rg_* \circ j_{1, !} =
R\overline{f}_{1, *} \circ j_{1, !}
$$
en utilisant le lemme \ref{lemma-shriek-lqf-and-proper} pour la première égalité.

\medskip\noindent
Pour achever la démonstration, puisque la catégorie des compactifications de $X$ sur $Y$
est cofiltrante, il suffit de montrer que la composition des morphismes de
compactifications de $X$ sur $Y$ correspond à la composition
des isomorphismes de foncteurs\footnote{En effet, si $\alpha, \beta : F \to G$
sont des morphismes de foncteurs et $\gamma : G \to H$ un isomorphisme
de foncteurs tel que $\gamma \circ \alpha = \gamma \circ \beta$, alors
on conclut que $\alpha = \beta$.}. Pour cela, supposons que
$j_3 : X \to \overline{X}_3$
soit une troisième compactification et que $h : \overline{X}_2 \to \overline{X}_3$
soit un morphisme de compactifications. Il faut alors montrer que la
composée
$$
R\overline{f}_{3, *} \circ j_{3, !} =
R\overline{f}_{3, *} \circ Rh_* \circ j_{2, !} =
R\overline{f}_{2, *} \circ j_{2, !} =
R\overline{f}_{2, *} \circ Rg_* \circ j_{1, !} =
R\overline{f}_{1, *} \circ j_{1, !}
$$
est égale à l'isomorphisme de foncteurs construit en utilisant seulement
$j_3$, $g \circ h$ et $j_1$. Un calcul montre qu'il suffit de
prouver que la composée des morphismes
$$
j_{3, !} \to Rh_* \circ j_{2, !} \to Rh_* \circ Rg_* \circ j_{1, !}
$$
du lemme \ref{lemma-shriek-lqf-and-proper} coïncide avec le morphisme correspondant
$j_{3, !} \to R(h \circ g)_* \circ j_{1, !}$
via l'identification $R(h \circ g)_* = Rh_* \circ Rg_*$.
Comme le morphisme du lemme \ref{lemma-shriek-lqf-and-proper}
est un cas particulier de celui du lemme \ref{lemma-shriek-proper-and-open}
(puisque $j_1$ et $j_2$ sont séparés), cela résulte immédiatement du
lemme \ref{lemma-shriek-proper-and-open-compose}.
\end{proof}

\begin{lemma}
\label{lemma-shriek-composition}
Soient $f : X \to Y$ et $g : Y \to Z$ des morphismes séparés de type fini
entre schémas quasi-compacts et quasi-séparés. Il existe alors un
isomorphisme canonique $Rg_! \circ Rf_! \to R(g \circ f)_!$.
\end{lemma}

\begin{proof}
Choisissons une compactification $i : Y \to \overline{Y}$ de $Y$ sur $Z$.
Choisissons une compactification $X \to \overline{X}$ de $X$ sur
$\overline{Y}$. On applique deux fois le théorème \ref{flat-theorem-nagata}
et le lemme \ref{flat-lemma-compactifyable} de Compléments sur la platitude.
Soit $U$ l'image inverse de $Y$ dans $\overline{X}$ ;
on obtient ainsi le diagramme commutatif
$$
\xymatrix{
X \ar[r]_j \ar[d]_f &
U \ar[dl]^{f'} \ar[r]_{j'} &
\overline{X} \ar[dl]^{\overline{f}} \\
Y \ar[r]_i \ar[d]_g &
\overline{Y} \ar[dl]^{\overline{g}} \\
Z
}
$$
On a alors
\begin{align*}
R(g \circ f)_!
& =
R(\overline{g} \circ \overline{f})_* \circ (j' \circ j)_! \\
& =
R\overline{g}_* \circ R\overline{f}_* \circ j'_! \circ j_! \\
& =
R\overline{g}_* \circ i_! \circ Rf'_* \circ j_! \\
& =
Rg_! \circ Rf_!
\end{align*}
La première égalité est la définition de $R(g \circ f)_!$.
La deuxième utilise les identifications
$R(\overline{g} \circ \overline{f})_* = R\overline{g}_* \circ R\overline{f}_*$
et $(j' \circ j)_! = j'_! \circ j_!$ du lemme \ref{lemma-f-shriek-composition}.
L'identification $i_! \circ Rf'_* \to R\overline{f}_* \circ j_!$
employée dans la troisième égalité est celle du lemme \ref{lemma-shriek-proper-and-open}.
La quatrième et dernière égalité est la définition de $Rg_!$ et $Rf_!$.
Pour achever la démonstration, montrons que
cet isomorphisme est indépendant des choix effectués.

\medskip\noindent
Supposons donnés deux diagrammes
$$
\vcenter{
\xymatrix{
X \ar[r]_{j_1} \ar[d] &
U_1 \ar[dl]^{f_1} \ar[r]_{j'_1} &
\overline{X}_1 \ar[dl]^{\overline{f}_1} \\
Y \ar[r]_{i_1} \ar[d] &
\overline{Y}_1 \ar[dl]^{\overline{g}_1} \\
Z
}
}
\quad\text{et}\quad
\vcenter{
\xymatrix{
X \ar[r]_{j_2} \ar[d] &
U_2 \ar[dl]^{f_2} \ar[r]_{j'_2} &
\overline{X}_2 \ar[dl]^{\overline{f}_2} \\
Y \ar[r]_{i_2} \ar[d] &
\overline{Y}_2 \ar[dl]^{\overline{g}_2} \\
Z
}
}
$$
Nous pouvons d'abord choisir une compactification $i : Y \to \overline{Y}$
de $Y$ sur $Z$ qui domine à la fois $\overline{Y}_1$ et $\overline{Y}_2$ ;
voir le lemme \ref{flat-lemma-compactifications-cofiltered} de Compléments sur la platitude.
D'après le lemme \ref{flat-lemma-right-multiplicative-system} de Compléments sur la platitude et les
lemmes \ref{categories-lemma-morphisms-right-localization} et
\ref{categories-lemma-equality-morphisms-right-localization}
de Catégories, nous pouvons choisir une compactification $X \to \overline{X}$ de
$X$ sur $\overline{Y}$, avec des morphismes $\overline{X} \to \overline{X}_1$
et $\overline{X} \to \overline{X}_2$, telle que la composée
$\overline{X} \to \overline{Y} \to \overline{Y}_1$ soit égale à
la composée $\overline{X} \to \overline{X}_1 \to \overline{Y}_1$
et que la composée
$\overline{X} \to \overline{Y} \to \overline{Y}_2$ soit égale à
la composée $\overline{X} \to \overline{X}_2 \to \overline{Y}_2$.
Il suffit donc de comparer les morphismes
déterminés par nos diagrammes lorsque l'on dispose d'un diagramme commutatif
de la forme suivante :
$$
\xymatrix{
X \ar[rr]_{j_1} \ar@{=}[d] & &
U_1 \ar[d]^{h'} \ar[ddll] \ar[rr]_{j'_1} & &
\overline{X}_1 \ar[d]^h \ar[ddll] \\
X \ar'[r][rr]^-{j_2} \ar[d] & &
U_2 \ar'[dl][ddll] \ar'[r][rr]^-{j'_2} & &
\overline{X}_2 \ar[ddll] \\
Y \ar[rr]^{i_1} \ar@{=}[d] & & \overline{Y}_1 \ar[d]^k \\
Y \ar[rr]^{i_2} \ar[d] & & \overline{Y}_2 \ar[dll] \\
Z
}
$$
Chacun des carrés
$$
\xymatrix{
X \ar[r]_{j_1} \ar[d]_{\text{id}} \ar@{}[dr]|A &
U_1 \ar[d]^{h'} \\
X \ar[r]^{j_2} &
U_2
}
\quad
\xymatrix{
U_2 \ar[r]_{j_2'} \ar[d]_{f_2} \ar@{}[dr]|B &
\overline{X}_2 \ar[d]^{\overline{f}_2} \\
Y \ar[r]^{i_2} &
\overline{Y}_2
}
\quad
\xymatrix{
U_1 \ar[r]_{j_1'} \ar[d]_{f_1} \ar@{}[dr]|C &
\overline{X}_1 \ar[d]^{\overline{f}_1} \\
Y \ar[r]^{i_1} &
\overline{Y}_1
}
\quad
\xymatrix{
Y \ar[r]_{i_1} \ar[d]_{\text{id}} \ar@{}[dr]|D &
\overline{Y}_1 \ar[d]^k \\
Y \ar[r]^{i_2} &
\overline{Y}_2
}
\quad
\xymatrix{
X \ar[r]_{j_1' \circ j_1} \ar[d]_{\text{id}} \ar@{}[dr]|E &
\overline{X}_1 \ar[d]^h \\
X \ar[r]^{j_2} &
\overline{X}_2
}
$$
donne lieu à un isomorphisme comme suit :
\begin{align*}
\gamma_A & :
j_{2, !} \to Rh'_* \circ j_{1, !}  \\
\gamma_B & :
i_{2, !} \circ Rf_{2, *} \to R\overline{f}_{2, *} \circ j'_{2, !} \\
\gamma_C & :
i_{1, !} \circ Rf_{1, *} \to R\overline{f}_{1, *} \circ j'_{1, !} \\
\gamma_D & :
i_{2, !} \to Rk_* \circ i_{1, !} \\
\gamma_E & :
j_{2, !} \to Rh_* \circ (j'_1 \circ j_1)_!
\end{align*}
en appliquant la flèche du lemme \ref{lemma-shriek-proper-and-open}
(qui est la même que celle du lemme \ref{lemma-shriek-lqf-and-proper}
lorsque la flèche verticale gauche est l'identité). Posons
\begin{align*}
F_1 & = Rf_{1, *} \circ j_{1, !} \\
F_2 & = Rf_{2, *} \circ j_{2, !} \\
G_1 & = R\overline{g}_{1, *} \circ i_{1, !} \\
G_2 & = R\overline{g}_{2, *} \circ i_{2, !} \\
C_1 & = R(\overline{g}_1 \circ \overline{f}_1)_* \circ (j'_1 \circ j_1)_! \\
C_2 & = R(\overline{g}_2 \circ \overline{f}_2)_* \circ (j'_2 \circ j_2)_!
\end{align*}
La construction donnée dans le premier paragraphe de la démonstration
et dans le lemme \ref{lemma-shriek-well-defined} utilise
\begin{enumerate}
\item $\gamma_C$ pour la flèche $G_1 \circ F_1 \to C_1$,
\item $\gamma_B$ pour la flèche $G_2 \circ F_2 \to C_2 $,
\item $\gamma_A$ pour la flèche $F_2 \to F_1$,
\item $\gamma_D$ pour la flèche $G_2 \to G_1$, et
\item $\gamma_E$ pour la flèche $C_2 \to C_1$.
\end{enumerate}
Il faut donc montrer que le diagramme
$$
\xymatrix{
C_2 \ar[rr]_{\gamma_E} & &
C_1 \\
G_2 \circ F_2 \ar[rr]^{\gamma_D \circ \gamma_A} \ar[u]^{\gamma_B} & &
G_1 \circ F_1 \ar[u]_{\gamma_C}
}
$$
est commutatif. Nous utiliserons les
lemmes \ref{lemma-shriek-proper-and-open-compose} et
\ref{lemma-shriek-proper-and-open-compose-horizontal}
et reprendrons, avec le même léger abus, les notations de la
remarque \ref{remark-going-around} (en omettant notamment, dans la notation,
les produits $\star$ par des transformations
identiques).
On peut écrire $\gamma_E = \gamma_F \circ \gamma_A$, où
$$
\xymatrix{
U_1 \ar[r]_{j'_1} \ar[d]_{h'} \ar@{}[rd]|F &
\overline{X}_1 \ar[d]^h \\
U_2 \ar[r]^{j'_2} &
\overline{X}_2
}
$$
On a donc
$$
\gamma_E \circ \gamma_B = \gamma_F \circ \gamma_A  \circ \gamma_B
= \gamma_F \circ \gamma_B \circ \gamma_A
$$
la dernière égalité venant de ce que les deux carrés $A$ et $B$ ne
se rencontrent qu'en un point (comme dans le dernier argument de la
remarque \ref{remark-going-around}). Il suffit donc de prouver que
$\gamma_C \circ \gamma_D = \gamma_F \circ \gamma_B$.
Comme chacune de ces deux flèches est égale à celle associée au carré
$$
\xymatrix{
U_1 \ar[r] \ar[d] & \overline{X}_1 \ar[d] \\
Y \ar[r] & \overline{Y}_2
}
$$
on conclut.
\end{proof}

\begin{lemma}
\label{lemma-pseudo-functor}
Soient $f : X \to Y$, $g : Y \to Z$, $h : Z \to T$ des morphismes séparés de
type fini entre schémas quasi-compacts et quasi-séparés. Alors
le diagramme
$$
\xymatrix{
Rh_! \circ Rg_! \circ Rf_! \ar[r]_{\gamma_C} \ar[d]^{\gamma_A} &
R(h \circ g)_! \circ Rf_! \ar[d]_{\gamma_{A + B}} \\
Rh_! \circ R(g \circ f)_! \ar[r]^{\gamma_{B + C}} &
R(h \circ g \circ f)_!
}
$$
d'isomorphismes du lemme \ref{lemma-shriek-composition} est commutatif
(pour la signification des $\gamma$, voir la démonstration).
\end{lemma}

\begin{proof}
Pour cela, choisissons une compactification $\overline{Z}$
de $Z$ sur $T$, puis une compactification $\overline{Y}$ de $Y$ sur
$\overline{Z}$, et enfin une compactification $\overline{X}$ de
$X$ sur $\overline{Y}$. On utilise le
théorème \ref{flat-theorem-nagata} et le
lemme \ref{flat-lemma-compactifyable} de Compléments sur la platitude.
Soit $W \subset \overline{Y}$ l'image inverse de $Z$ par
$\overline{Y} \to \overline{Z}$, et soient $U \subset V \subset \overline{X}$
les images inverses respectives des deux sous-schémas $Y \subset W$ par $\overline{X} \to \overline{Y}$.
On obtient le diagramme suivant :
$$
\xymatrix{
X \ar[d]_f \ar[r] & U \ar[r] \ar[d] \ar@{}[dr]|A &
V \ar[d] \ar[r] \ar@{}[rd]|B & \overline{X} \ar[d] \\
Y \ar[d]_g \ar[r] & Y \ar[r] \ar[d] & W \ar[r] \ar[d] \ar@{}[rd]|C &
\overline{Y} \ar[d] \\
Z \ar[d]_h \ar[r] & Z \ar[d] \ar[r] & Z \ar[d] \ar[r] & \overline{Z} \ar[d] \\
T \ar[r] & T \ar[r] & T \ar[r] & T
}
$$
Sans multiplier les notations, mais en raisonnant exactement
comme dans la démonstration du lemme \ref{lemma-shriek-composition},
on voit que les flèches du premier diagramme affiché utilisent les
flèches du lemme \ref{lemma-shriek-proper-and-open} pour les rectangles
$A + B$, $B + C$, $A$ et $C$, comme l'indique le diagramme de
l'énoncé du lemme. D'après les
lemmes \ref{lemma-shriek-proper-and-open-compose} et
\ref{lemma-shriek-proper-and-open-compose-horizontal}
on a $\gamma_{A + B} = \gamma_B \circ \gamma_A$ et
$\gamma_{B + C} = \gamma_B \circ \gamma_C$ ; on en conclut
que l'égalité voulue est vraie pourvu que
$\gamma_A \circ \gamma_C = \gamma_C \circ \gamma_A$.
C'est bien le cas, car les deux carrés $A$ et $C$ ne
se rencontrent qu'en un point (comme dans le dernier argument de la
remarque \ref{remark-going-around}).
\end{proof}

\begin{lemma}
\label{lemma-base-change-shriek}
Considérons un carré cartésien
$$
\xymatrix{
X' \ar[r]_{g'} \ar[d]_{f'} & X \ar[d]^f \\
Y' \ar[r]^g & Y
}
$$
de schémas quasi-compacts et quasi-séparés, où $f$ est séparé et
de type fini. Il existe alors un isomorphisme canonique
$$
g^{-1} \circ Rf_! \to Rf'_! \circ (g')^{-1}
$$
De plus, ces isomorphismes sont compatibles avec ceux
du lemme \ref{lemma-shriek-composition}.
\end{lemma}

\begin{proof}
Choisissons une compactification $j : X \to \overline{X}$ sur $Y$
et notons $\overline{f} : \overline{X} \to Y$ le morphisme structural.
Notons $j' : X' \to \overline{X}'$ et $\overline{f}' : \overline{X}' \to Y'$
les morphismes déduits de $j$ et de $\overline{f}$ par changement de base.
Puisque $Rf_! = R\overline{f}_* \circ j_!$ et $Rf'_! =
R\overline{f}'_* \circ j'_!$, l'isomorphisme se construit au moyen de
$$
g^{-1} \circ R\overline{f}_* \circ j_! \to
R\overline{f}'_* \circ (\overline{g}')^{-1} \circ j_! \to
R\overline{f}'_* \circ j'_! \circ (g')^{-1} 
$$
où la première flèche est l'isomorphisme fourni par le
théorème de changement de base propre (lemme de Cohomologie étale
\ref{etale-cohomology-lemma-proper-base-change} dans le cas borné inférieurement
à cohomologie de torsion, et le lemme de Cohomologie étale
\ref{etale-cohomology-lemma-proper-base-change-mod-n} dans le
cas où $\Lambda$ est un anneau de torsion), tandis que la deuxième est l'isomorphisme du
lemme \ref{lemma-base-change-f-shriek-separated}.

\medskip\noindent
Pour achever la démonstration, il faut établir deux points : d'abord que
l'isomorphisme de foncteurs ainsi obtenu ne dépend pas du
choix de la compactification ; ensuite que, si l'on empile verticalement
deux diagrammes de changement de base comme dans le lemme,
les isomorphismes de changement de base sont compatibles avec ceux
du lemme \ref{lemma-shriek-composition}.
Un argument direct, que nous omettons, montre que ces deux points découlent
de la compatibilité au changement de base des isomorphismes
\begin{enumerate}
\item $Rg_* \circ Rf_* = R(g \circ f)_*$ pour $f : X \to Y$ et $g : Y \to Z$
propres,
\item $g_! \circ f_! = (g \circ f)_!$ pour $f : X \to Y$ et $g : Y \to Z$
séparés et quasi-finis, et
\item $g_! \circ Rf'_* = Rf_* \circ g'_!$ pour $f : X \to Y$ et
$f' : X' \to Y'$ propres, et $g : Y' \to Y$ et $g' : X' \to X$
séparés et quasi-finis, avec $f \circ g' = g \circ f'$,
\end{enumerate}
sont compatibles au changement de base. Cela vaut pour (1) d'après la
remarque \ref{sites-cohomology-remark-compose-base-change} de Cohomologie sur les sites,
pour (2) d'après la remarque \ref{remark-f-shriek-base-change-composition}, et pour
(3) d'après le lemme \ref{lemma-shriek-proper-and-open-base-change}.
\end{proof}

\begin{remark}
\label{remark-shriek-to-star}
Soit $f : X \to Y$ un morphisme de schémas séparé et de type fini,
où $Y$ est quasi-compact et quasi-séparé. Nous construirons plus loin
un morphisme
$$
Rf_!K \longrightarrow Rf_*K
$$
fonctoriel en $K$ dans $D^+_{tors}(X_\etale, \Lambda)$, ou dans
$D(X_\etale, \Lambda)$ si $\Lambda$ est un anneau de torsion. Dans les deux cas, cette transformation
de foncteurs est compatible avec
\begin{enumerate}
\item l'isomorphisme $Rg_! \circ Rf_! \to R(g \circ f)_!$ du
lemme \ref{lemma-shriek-composition} et l'isomorphisme
$Rg_* \circ Rf_* \to R(g \circ f)_*$ du
lemme de Cohomologie sur les sites
\ref{sites-cohomology-lemma-derived-pushforward-composition} ;
\item l'isomorphisme $g^{-1} \circ Rf_! \to Rf'_! \circ (g')^{-1}$
du lemme \ref{lemma-base-change-shriek}
et la flèche de changement de base de la
remarque \ref{sites-cohomology-remark-base-change} de Cohomologie sur les sites.
\end{enumerate}
En effet, choisissons une compactification $j : X \to \overline{X}$ sur $Y$
et notons $\overline{f} : \overline{X} \to Y$ le morphisme structural.
Puisque $Rf_! = R\overline{f}_* \circ j_!$ et
$Rf_* = R\overline{f}_* \circ Rj_*$, il suffit de construire une
transformation de foncteurs $j_! \to Rj_*$. Nous utilisons pour cela la
transformation canonique $j_! \to j_*$ du
lemme de Cohomologie étale
\ref{etale-cohomology-lemma-shriek-into-star-separated-etale}.
Nous omettons la démonstration de l'indépendance de la transformation obtenue
par rapport au choix de la compactification, ainsi que celle des
compatibilités (1) et (2).
\end{remark}













\section{Propriétés de l'image directe dérivée à support propre}
\label{section-derived-lower-shriek-properties}

\noindent
Voici quelques propriétés de l'image directe dérivée à support propre.

\begin{lemma}
\label{lemma-derived-lower-shriek-commute-direct-sums}
Soit $f : X \to Y$ un morphisme séparé de type fini entre schémas quasi-compacts
et quasi-séparés. Soit $\Lambda$ un anneau.
\begin{enumerate}
\item Considérons une famille d'objets $K_i \in D^+_{tors}(X_\etale, \Lambda)$, $i \in I$.
Supposons donné $a \in \mathbf{Z}$ tel que
$H^n(K_i) = 0$ pour $n < a$ et $i \in I$. Alors $Rf_!(\bigoplus_i K_i) =
\bigoplus_i Rf_!K_i$.
\item Si $\Lambda$ est un anneau de torsion, le foncteur
$Rf_! : D(X_\etale, \Lambda) \to D(Y_\etale, \Lambda)$
commute aux sommes directes.
\end{enumerate}
\end{lemma}

\begin{proof}
Par construction, il suffit de le démontrer lorsque $f$ est une immersion ouverte
et lorsque $f$ est un morphisme propre. Pour toute immersion ouverte $j : U \to X$
de schémas, le foncteur $j_! : D(U_\etale) \to D(X_\etale)$ est
adjoint à gauche de l'image inverse $j^{-1} : D(X_\etale) \to D(U_\etale)$
et commute donc aux sommes directes ; voir le lemme de Cohomologie sur les sites
\ref{sites-cohomology-lemma-adjoint-lower-shriek-restrict}.
Dans le cas propre, on a $Rf_! = Rf_*$ et le résultat découle du
lemme de Cohomologie étale
\ref{etale-cohomology-lemma-direct-sum-bounded-below-pushforward}
dans le cas borné inférieurement, et du
lemme \ref{etale-cohomology-lemma-proper-mod-n-direct-sums} de Cohomologie étale
lorsque $\Lambda$ est un anneau de torsion.
\end{proof}

\begin{lemma}
\label{lemma-derived-lower-shriek-bounded}
Soit $f : X \to Y$ un morphisme séparé de type fini entre schémas quasi-compacts
et quasi-séparés. Soit $\Lambda$ un anneau. Les foncteurs $Rf_!$
construits à la section \ref{section-derived-lower-shriek-compactification}
sont bornés au sens suivant : il existe un entier $N$ tel que,
pour $E \in D^+_{tors}(X_\etale, \Lambda)$ ou $E \in D(X_\etale, \Lambda)$
si $\Lambda$ est un anneau de torsion, on ait :
\begin{enumerate}
\item $H^i(Rf_!(\tau_{\leq a}E) \to H^i(Rf_!(E))$ est un isomorphisme
pour $i \leq a$ ;
\item $H^i(Rf_!(E)) \to H^i(Rf_!(\tau_{\geq b - N}E))$ est un isomorphisme
pour $i \geq b$ ;
\item si $H^i(E) = 0$ pour $i \not \in [a, b]$ avec
$-\infty \leq a \leq b \leq \infty$, alors $H^i(Rf_!(E)) = 0$
pour $i \not \in [a, b + N]$.
\end{enumerate}
\end{lemma}

\begin{proof}
Supposons que $\Lambda$ soit un anneau de torsion et considérons le foncteur
$Rf_! : D(X_\etale, \Lambda) \to D(Y_\etale, \Lambda)$.
Par construction, il suffit de le démontrer lorsque $f$ est une immersion ouverte
et lorsque $f$ est un morphisme propre. Pour toute immersion ouverte $j : U \to X$
de schémas, le foncteur $j_! : D(U_\etale) \to D(X_\etale)$ est exact ;
l'énoncé vaut donc avec $N = 0$ dans ce cas.
Si $f$ est propre, alors $Rf_! = Rf_*$, c'est-à-dire qu'il s'agit d'un foncteur
dérivé à droite. La borne à gauche résulte donc du
lemme \ref{derived-lemma-negative-vanishing} de Catégories dérivées.
De plus, dans ce cas, $f_* : \textit{Mod}(X_\etale, \Lambda)
\to \textit{Mod}(Y_\etale, \Lambda)$ est de dimension cohomologique finie d'après le
lemme \ref{morphisms-lemma-morphism-finite-type-bounded-dimension} de Morphismes
et le lemme de Cohomologie étale
\ref{etale-cohomology-lemma-cohomological-dimension-proper}.
On conclut donc par le
lemme \ref{derived-lemma-unbounded-right-derived} de Catégories dérivées.

\medskip\noindent
Supposons maintenant $\Lambda$ quelconque et considérons le foncteur
$Rf_! : D^+_{tors}(X_\etale, \Lambda) \to D^+_{tors}(Y_\etale, \Lambda)$.
On se ramène encore immédiatement au cas où $f$ est propre et
$Rf_! = Rf_*$. Le point (1) est encore immédiat. Pour démontrer le point (3),
on peut raisonner par récurrence sur $b - a$, en utilisant les triangles distingués
de troncation et le lemme de Cohomologie étale
\ref{etale-cohomology-lemma-cohomological-dimension-proper}.
Le point (2) résulte du point (3). Nous omettons les détails.
\end{proof}

\begin{lemma}
\label{lemma-Rf-shriek-for-quasi-finite}
Soit $f : X \to Y$ un morphisme quasi-fini séparé entre schémas quasi-compacts et
quasi-séparés. Alors les foncteurs $Rf_!$ construits à la
section \ref{section-derived-lower-shriek-compactification}
coïncident avec la restriction du foncteur
$f_! : D(X_\etale, \Lambda) \to D(Y_\etale, \Lambda)$
construit à la section \ref{section-derived-duality-locally-quasi-finite}
à leurs domaines de définition communs.
\end{lemma}

\begin{proof}
Par le théorème principal de Zariski (lemme de Compléments sur les morphismes
\ref{more-morphisms-lemma-quasi-finite-separated-pass-through-finite})
on peut trouver une immersion ouverte $j : X \to \overline{X}$ et un morphisme fini
$\overline{f} : \overline{X} \to Y$ tels que $f = \overline{f} \circ j$.
Par construction, on a $Rf_! = R\overline{f}_* \circ j_!$.
Comme $\overline{f}$ est fini, on a $R\overline{f}_* = \overline{f}_*$
d'après la proposition de Cohomologie étale
\ref{etale-cohomology-proposition-finite-higher-direct-image-zero}.
Le lemme en résulte, car $\overline{f}_* \circ j_! = f_!$,
par exemple d'après le lemme \ref{lemma-compactify-f-shriek-separated}.
\end{proof}

\begin{lemma}
\label{lemma-relative-mayer-vietoris}
Soit $f : X \to Y$ un morphisme séparé de type fini entre schémas quasi-compacts
et quasi-séparés. Soient $U$ et $V$ des ouverts quasi-compacts de $X$
tels que $X = U \cup V$. Notons $a : U \to Y$, $b : V \to Y$ et
$c : U \cap V \to Y$ les restrictions de $f$. Soit $\Lambda$ un anneau.
Pour $K$ dans $D^+_{tors}(X_\etale, \Lambda)$, ou $K \in D(X_\etale, \Lambda)$
si $\Lambda$ est un anneau de torsion, on a un triangle distingué
$$
Rc_!(K|_{U \cap V}) \to
Ra_!(K|_U) \oplus Rb_!(K|_V) \to
Rf_!K \to 
Rc_!(K|_{U \cap V})[1]
$$
dans $D(Y_\etale, \Lambda)$.
\end{lemma}

\begin{proof}
Cela résulte du lemme \ref{lemma-mayer-vietoris-shriek},
de l'égalité $Rf_! \circ Rj_{U!} = Ra_!$ donnée par le
lemme \ref{lemma-shriek-composition}, et de l'égalité
$Rj_{U!} = j_{U!}$ donnée par le lemme \ref{lemma-Rf-shriek-for-quasi-finite}.
\end{proof}

\begin{lemma}
\label{lemma-relative-triangle-associated-to-open}
Soit $f : X \to Y$ un morphisme séparé de type fini entre schémas quasi-compacts
et quasi-séparés. Soit $U$ un ouvert quasi-compact de $X$ dont
le complémentaire est $Z \subset X$. Notons $g : U \to Y$ et $h : Z \to Y$
les restrictions de $f$. Soit $\Lambda$ un anneau.
Pour $K$ dans $D^+_{tors}(X_\etale, \Lambda)$, ou $K \in D(X_\etale, \Lambda)$
si $\Lambda$ est un anneau de torsion, on a un triangle distingué
$$
Rg_!(K|_U) \to
Rf_!K \to
Rh_!(K|_Z) \to
Rg_!(K|_U)[1]
$$
dans $D(Y_\etale, \Lambda)$.
\end{lemma}

\begin{proof}
Cela résulte du
lemme \ref{lemma-triangle-associated-to-open},
des égalités $Rf_! \circ Rj_! = Rg_!$ et $Rf_! \circ Ri_!$ données par le
lemme \ref{lemma-shriek-composition}, ainsi que des égalités
$Rj_! = j_!$ et $Ri_! = i_! = i_*$ données par le
lemme \ref{lemma-Rf-shriek-for-quasi-finite}.
\end{proof}

\begin{lemma}
\label{lemma-shriek-and-thickening}
Soit $f' : X' \to Y$ un morphisme séparé de type fini entre schémas quasi-compacts
et quasi-séparés. Soit $i : X \to X'$ un épaississement et
notons $f = f' \circ i$. Soit $\Lambda$ un anneau.
Pour $K'$ dans $D^+_{tors}(X'_\etale, \Lambda)$, ou $K' \in D(X'_\etale, \Lambda)$
si $\Lambda$ est un anneau de torsion, on a $Rf_!i^{-1}K' = Rf'_!K'$.
\end{lemma}

\begin{proof}
Cela résulte de ce que $i^{-1}$ et $i_* = i_!$ sont des équivalences de catégories
quasi-inverses l'une de l'autre par invariance topologique du petit topos étale
(théorème de Cohomologie étale
\ref{etale-cohomology-theorem-topological-invariance}),
ainsi que du
lemme \ref{lemma-shriek-composition}.
\end{proof}

\begin{lemma}
\label{lemma-projection-formula-Rf-lower-shriek}
Soit $f : X \to Y$ un morphisme séparé de type fini entre schémas quasi-compacts
et quasi-séparés. Soit $\Lambda$ un anneau de torsion. Soient
$E \in D(X_\etale, \Lambda)$ et $K \in D(Y_\etale, \Lambda)$. Alors
$$
Rf_!E \otimes_\Lambda^\mathbf{L} K =
Rf_!(E \otimes_\Lambda^\mathbf{L} f^{-1}K)
$$
dans $D(Y_\etale, \Lambda)$.
\end{lemma}

\begin{proof}
Choisissons $j : X \to \overline{X}$ et $\overline{f} : \overline{X} \to Y$
comme dans la construction de $Rf_!$. On a
$j_!E \otimes_\Lambda^\mathbf{L} \overline{f}^{-1}K =
j_!(E \otimes_\Lambda^\mathbf{L} f^{-1}K)$ d'après
le lemme \ref{sites-cohomology-lemma-j-shriek-and-tensor} de Cohomologie sur les sites.
On obtient alors le résultat en appliquant le
lemme de Cohomologie étale
\ref{etale-cohomology-lemma-projection-formula-proper-mod-n}
et en utilisant $f^{-1} = j^{-1} \circ \overline{f}^{-1}$ et
$Rf_! = R\overline{f}_*j_!$.
\end{proof}

\begin{remark}
\label{remark-Rf-lower-shriek-change-of-rings}
Soit $\Lambda_1 \to \Lambda_2$ un homomorphisme d'anneaux de torsion.
Soit $f : X \to Y$ un morphisme séparé de type fini entre schémas quasi-compacts
et quasi-séparés. Le diagramme
$$
\xymatrix{
D(X_\etale, \Lambda_2) \ar[r]_{res} \ar[d]_{Rf_!} &
D(X_\etale, \Lambda_1) \ar[d]^{Rf_!} \\
D(Y_\etale, \Lambda_2) \ar[r]^{res} &
D(Y_\etale, \Lambda_1)
}
$$
est commutatif, où $res$ est le foncteur de restriction des scalaires qui considère un
$\Lambda_2$-module comme un $\Lambda_1$-module au moyen du morphisme d'anneaux donné.
Écrivons $Rf_! = R\overline{f}_* \circ j_!$ pour une factorisation
$f = \overline{f} \circ j$ comme dans la
section \ref{section-derived-lower-shriek-compactification}. On voit que
le résultat se vérifie directement pour $j_!$ et résulte, pour $R\overline{f}_*$,
du lemme de Cohomologie sur les sites
\ref{sites-cohomology-lemma-modules-abelian-unbounded}.
D'autre part, le diagramme
$$
\xymatrix{
D(X_\etale, \Lambda_1)
\ar[r]_{- \otimes_{\Lambda_1}^\mathbf{L} \Lambda_2} \ar[d]_{Rf_!} &
D(X_\etale, \Lambda_2) \ar[d]^{Rf_!} \\
D(Y_\etale, \Lambda_1)
\ar[r]^{- \otimes_{\Lambda_1}^\mathbf{L} \Lambda_2} &
D(Y_\etale, \Lambda_2)
}
$$
est également commutatif d'après le
lemme \ref{lemma-projection-formula-Rf-lower-shriek}.
\end{remark}

\begin{remark}
\label{remark-projection-formula-for-internal-hom}
Soit $f : X \to Y$ un morphisme séparé de type fini entre schémas quasi-compacts
et quasi-séparés. Soit $\Lambda$ un anneau de torsion.
Soient $K$ et $L$ des objets de $D(X_\etale, \Lambda)$. Nous affirmons
qu'il existe un morphisme canonique
$$
\alpha :
Rf_*R\SheafHom_\Lambda(K, L)
\longrightarrow
R\SheafHom_\Lambda(Rf_!K, Rf_!L)
$$
fonctoriel en $K$ et $L$. Choisissons en effet $j : X \to \overline{X}$ et
$\overline{f} : \overline{X} \to Y$ comme dans la construction de $Rf_!$.
Définissons d'abord un morphisme
$$
\beta :
Rj_*R\SheafHom_\Lambda(K, L)
\longrightarrow
R\SheafHom_\Lambda(j_!K, j_!L)
$$
Par la construction du Hom interne dans la catégorie dérivée, cela revient
à définir un morphisme
$$
\beta' :
Rj_*R\SheafHom_\Lambda(K, L) \otimes_\Lambda^\mathbf{L} j_!K
\longrightarrow
j_!L
$$
Voir Cohomologie sur les sites, section \ref{sites-cohomology-section-internal-hom}.
La source de $\beta'$ est égale à
$$
j_!\left(R\SheafHom_\Lambda(K, L) \otimes_\Lambda^\mathbf{L} K\right)
$$
par le lemme \ref{sites-cohomology-lemma-j-shriek-and-tensor} de Cohomologie sur les sites.
On peut donc poser $\beta' = j_!\beta''$, où
$\beta'' : R\SheafHom_\Lambda(K, L) \otimes_\Lambda^\mathbf{L} K \to L$
correspond à l'identité de $R\SheafHom_\Lambda(K, L)$
par la propriété universelle du Hom interne mentionnée ci-dessus. D'après la
remarque de Cohomologie sur les sites
\ref{sites-cohomology-remark-projection-formula-for-internal-hom}
on a un morphisme canonique
$$
\gamma :
R\overline{f}_*R\SheafHom_\Lambda(j_!K, j_!L)
\longrightarrow
R\SheafHom_\Lambda(R\overline{f}_*j_!K, R\overline{f}_*j_!L)
$$
Comme $Rf_! = R\overline{f}_*j_!$ et $Rf_* = R\overline{f}_* Rj_*$
(par Leray), on obtient le morphisme recherché
$\alpha = \gamma \circ R\overline{f}_*\beta$.
\end{remark}









\section{Image inverse extraordinaire dérivée}
\label{section-derived-upper-shriek}

\noindent
Nous obtenons $Rf^!$ au moyen d'un théorème de représentabilité de Brown.

\begin{lemma}
\label{lemma-upper-shriek-derived}
Soit $f : X \to Y$ un morphisme séparé de type fini entre schémas quasi-compacts
et quasi-séparés. Soit $\Lambda$ un anneau de torsion.
Le foncteur
$Rf_! : D(X_\etale, \Lambda) \to D(Y_\etale, \Lambda)$
admet un adjoint à droite
$Rf^! : D(Y_\etale, \Lambda) \to D(X_\etale, \Lambda)$.
\end{lemma}

\begin{proof}
Cela résulte de la
proposition \ref{injectives-proposition-brown} d'Injectifs
et du lemme \ref{lemma-derived-lower-shriek-commute-direct-sums} ci-dessus.
\end{proof}

\begin{lemma}
\label{lemma-Rf-upper-shriek-for-quasi-finite}
Soit $f : X \to Y$ un morphisme quasi-fini séparé entre schémas quasi-compacts
et quasi-séparés. Soit $\Lambda$ un anneau de torsion.
Le foncteur $Rf^! : D(Y_\etale, \Lambda) \to D(X_\etale, \Lambda)$
du lemme \ref{lemma-upper-shriek-derived} s'identifie au foncteur
$Rf^!$ du lemme \ref{lemma-lqf-shriek-derived}.
\end{lemma}

\begin{proof}
Cela résulte de l'unicité des adjoints, puisque $Rf_! = f_!$ d'après le
lemme \ref{lemma-Rf-shriek-for-quasi-finite}.
\end{proof}

\begin{lemma}
\label{lemma-Rf-upper-shriek-for-etale}
Soit $j : U \to X$ un morphisme étale séparé entre schémas quasi-compacts
et quasi-séparés. Soit $\Lambda$ un anneau de torsion.
Le foncteur $Rj^! : D(X_\etale, \Lambda) \to D(U_\etale, \Lambda)$
s'identifie à $j^{-1}$.
\end{lemma}

\begin{proof}
En effet, $Rj^!$ et $j^{-1}$ sont tous deux adjoints à droite du foncteur
$Rj_! = j_!$. Voir par exemple les lemmes
\ref{lemma-Rf-upper-shriek-for-quasi-finite} et
\ref{lemma-etale-upper-shriek}.
\end{proof}

\begin{lemma}
\label{lemma-twisted-inverse-image-bounded-below}
Soit $f : X \to Y$ un morphisme séparé de type fini entre
schémas quasi-compacts et quasi-séparés. Soit $\Lambda$ un anneau de
torsion. Le foncteur $Rf^!$ envoie $D^+(Y_\etale, \Lambda)$ dans
$D^+(X_\etale, \Lambda)$. Plus précisément, il existe un entier
$N \geq 0$ tel que, si $K \in D(Y_\etale, \Lambda)$ vérifie $H^i(K) = 0$
pour $i < a$, alors $H^i(Rf^!K) = 0$ pour $i < a - N$.
\end{lemma}

\begin{proof}
Soit $N$ l'entier fourni par le lemme \ref{lemma-derived-lower-shriek-bounded}.
Par construction, pour $K \in D(Y_\etale, \Lambda)$ et
$L \in \in D(X_\etale, \Lambda)$, on a
$\Hom_X(L, Rf^!K) = \Hom_Y(Rf_!L, K)$. Supposons $H^i(K) = 0$
pour $i < a$. Prenons alors $L = \tau_{\leq a - N - 1}Rf^!K$. D'après le
lemme \ref{lemma-derived-lower-shriek-bounded},
le complexe $Rf_!L$ a ses faisceaux de cohomologie nuls aux
degrés $\leq a - 1$. Ainsi, $\Hom_Y(Rf_!L, K) = 0$ d'après le
lemme \ref{derived-lemma-negative-exts} de Catégories dérivées.
Le morphisme canonique $\tau_{\leq a - N - 1}Rf^!K \to Rf^!K$
est donc nul, ce qui entraîne $H^i(Rf^!K) = 0$ pour $i \leq a - N - 1$.
\end{proof}

\noindent
Soit $f : X \to Y$ un morphisme séparé de type fini entre schémas quasi-séparés
et quasi-compacts. Soit $\Lambda$ un anneau de torsion.
Pour tous $K \in D(Y_\etale, \Lambda)$ et $L \in D(X_\etale, \Lambda)$,
on obtient un morphisme canonique
\begin{equation}
\label{equation-sheafy-trace}
Rf_*R\SheafHom_\Lambda(L, Rf^!K) \longrightarrow R\SheafHom_\Lambda(Rf_!L, K)
\end{equation}
Ce morphisme est construit comme la composée
$$
Rf_*R\SheafHom_\Lambda(L, Rf^!K) \to
R\SheafHom_\Lambda(Rf_!L, Rf_!Rf^!K) \to
R\SheafHom_\Lambda(Rf_!L, K)
$$
où la première flèche est celle de la
remarque \ref{remark-projection-formula-for-internal-hom}
et la seconde est le morphisme d'adjonction $Rf_!Rf^!K \to K$.

\begin{lemma}
\label{lemma-iso-on-RSheafHom}
Soit $f : X \to Y$ un morphisme séparé de type fini entre schémas quasi-compacts et
quasi-séparés. Soit $\Lambda$ un anneau de torsion.
Pour tous $K \in D(Y_\etale, \Lambda)$ et $L \in D(X_\etale, \Lambda)$,
le morphisme (\ref{equation-sheafy-trace})
$$
Rf_*R\SheafHom_\Lambda(L, Rf^!K) \longrightarrow R\SheafHom_\Lambda(Rf_!L, K)
$$
est un isomorphisme.
\end{lemma}

\begin{proof}
Pour démontrer le lemme, il suffit de montrer que, pour tout $M \in D(Y_\etale, \Lambda)$,
le morphisme (\ref{equation-sheafy-trace}) induit une bijection
$$
\Hom_Y(M, Rf_*R\SheafHom_\Lambda(L, Rf^!K))
\longrightarrow
\Hom_Y(M, R\SheafHom_\Lambda(Rf_!L, K))
$$
Pour le voir, utilisons la suite d'égalités suivante :
\begin{align*}
\Hom_Y(M, Rf_*R\SheafHom_\Lambda(L, Rf^!K))
& =
\Hom_X(f^{-1}M, R\SheafHom_\Lambda(L, Rf^!K)) \\
& =
\Hom_X(f^{-1}M \otimes_\Lambda^\mathbf{L} L, Rf^!K) \\
& =
\Hom_Y(Rf_!(f^{-1}M \otimes_\Lambda^\mathbf{L} L), K) \\
& =
\Hom_Y(M \otimes_\Lambda^\mathbf{L} Rf_!L, K) \\
& =
\Hom_Y(M, R\SheafHom_\Lambda(Rf_!L, K))
\end{align*}
La première égalité résulte du
lemme \ref{sites-cohomology-lemma-adjoint} de Cohomologie sur les sites.
La deuxième résulte du
lemme \ref{sites-cohomology-lemma-internal-hom} de Cohomologie sur les sites.
La troisième résulte de la construction de $Rf^!$.
La quatrième résulte du lemme \ref{lemma-projection-formula-Rf-lower-shriek}
(c'est l'étape essentielle).
La cinquième résulte du
lemme \ref{sites-cohomology-lemma-internal-hom} de Cohomologie sur les sites.
\end{proof}

\begin{lemma}
\label{lemma-iso-global-hom}
Soit $f : X \to Y$ un morphisme séparé de type fini entre schémas
quasi-séparés et quasi-compacts. Soit $\Lambda$ un anneau de torsion.
Pour tous $K \in D(Y_\etale, \Lambda)$ et $L \in D(X_\etale, \Lambda)$, le morphisme
(\ref{equation-sheafy-trace}) induit un isomorphisme
$$
R\Hom_X(L, Rf^!K) \longrightarrow R\Hom_Y(Rf_!L, K)
$$
entre les Hom dérivés globaux.
\end{lemma}

\begin{proof}
Par la construction donnée dans la
section \ref{sites-cohomology-section-global-RHom} de Cohomologie sur les sites,
on a
$$
R\Hom_X(L, Rf^!K) =
R\Gamma(X, R\SheafHom_\Lambda(L, Rf^!K)) =
R\Gamma(Y, Rf_*R\SheafHom_\Lambda(L, Rf^!K))
$$
(la seconde égalité résultant de Leray) et
$$
R\Hom_Y(Rf_!L, K) = R\Gamma(Y, R\SheafHom_\Lambda(Rf_!L, K))
$$
Le lemme résulte donc du lemme \ref{lemma-iso-on-RSheafHom}.
\end{proof}

\begin{lemma}
\label{lemma-cartesian-square-Rf-upper-shriek}
Considérons un carré cartésien
$$
\xymatrix{
X' \ar[r]_{g'} \ar[d]_{f'} & X \ar[d]^f \\
Y' \ar[r]^g & Y
}
$$
de schémas quasi-compacts et quasi-séparés, où $f$ est séparé et
de type fini. Alors $Rf^! \circ Rg_* = Rg'_* \circ R(f')^!$.
\end{lemma}

\begin{proof}
Par unicité des foncteurs adjoints, cela résulte du changement de base
pour l'image directe dérivée à support propre : on a
$g^{-1} \circ Rf_! = Rf'_! \circ (g')^{-1}$ d'après
le lemme \ref{lemma-base-change-shriek}.
\end{proof}

\begin{remark}
\label{remark-Rf-upper-shriek-change-of-rings}
Soit $\Lambda_1 \to \Lambda_2$ un homomorphisme d'anneaux de torsion.
Soit $f : X \to Y$ un morphisme séparé de type fini entre schémas quasi-compacts
et quasi-séparés. Le diagramme
$$
\xymatrix{
D(X_\etale, \Lambda_2) \ar[r]_{res} &
D(X_\etale, \Lambda_1) \\
D(Y_\etale, \Lambda_2) \ar[r]^{res} \ar[u]^{Rf^!} &
D(Y_\etale, \Lambda_1) \ar[u]_{Rf^!}
}
$$
est commutatif, où $res$ est le foncteur de restriction des scalaires qui considère un
$\Lambda_2$-module comme un $\Lambda_1$-module au moyen du morphisme d'anneaux donné.
Cela résulte de l'unicité des adjoints, du second diagramme commutatif
de la remarque \ref{remark-Rf-lower-shriek-change-of-rings}, et du fait que
$$
\Hom_{\Lambda_2}(K_1 \otimes_{\Lambda_1}^\mathbf{L} \Lambda_2, K_2) =
\Hom_{\Lambda_1}(K_1, res(K_2))
$$
Cette égalité, tant pour les objets sur $X_\etale$ que pour ceux sur
$Y_\etale$, est un cas très particulier du
lemme \ref{sites-cohomology-lemma-adjoint} de Cohomologie sur les sites.
\end{remark}











\section{Cohomologie à support propre}
\label{section-compactly-supported-cohomology}

\noindent
Soit $k$ un corps. Soit $\Lambda$ un anneau. Soit $X$ un schéma séparé
de type fini sur $k$, de morphisme structural $f : X \to \Spec(k)$.
À la section \ref{section-derived-lower-shriek-compactification},
nous avons défini le foncteur
$Rf_! : D^+_{tors}(X_\etale, \Lambda) \to D^+_{tors}(\Spec(k), \Lambda)$
et le foncteur $Rf_! : D(X_\etale, \Lambda) \to D(\Spec(k), \Lambda)$
si $\Lambda$ est un anneau de torsion. En composant avec le foncteur des sections globales
sur $\Spec(k)$, on obtient ce que nous appellerons la cohomologie à support propre.

\begin{definition}
\label{definition-cohomology-compact-support}
Soit $X$ un schéma séparé de type fini sur un corps $k$.
Soit $\Lambda$ un anneau. Soit $K$ un objet de
$D^+_{tors}(X_\etale, \Lambda)$
ou de $D(X_\etale, \Lambda)$ lorsque $\Lambda$ est un anneau de torsion.
La {\it cohomologie de $K$ à support propre}, ou la
{\it cohomologie à support propre de $K$}, est
$$
R\Gamma_c(X, K) = R\Gamma(\Spec(k), Rf_!K)
$$
où $f : X \to \Spec(k)$ est le morphisme structural. Nous
écrirons $H^i_c(X, K) = H^i(R\Gamma_c(X, K))$.
\end{definition}

\noindent
Nous vérifierons que cette définition est compatible avec la
définition \ref{definition-compact-support} grâce au
lemme \ref{lemma-compact-support-h0}.
L'intérêt de cette définition réside dans le résultat suivant.

\begin{lemma}
\label{lemma-stalk-R-f-shriek}
Soit $f : X \to Y$ un morphisme séparé de type fini entre schémas,
où $Y$ est quasi-compact et quasi-séparé. Soit $K$ un objet de
$D^+_{tors}(X_\etale, \Lambda)$ ou de $D(X_\etale, \Lambda)$ lorsque
$\Lambda$ est un anneau de torsion. Il existe alors un isomorphisme canonique
$$
(Rf_!K)_{\overline{y}}
\longrightarrow
R\Gamma_c(X_{\overline{y}}, K|_{X_{\overline{y}}})
$$
dans $D(\Lambda)$ pour tout point géométrique $\overline{y} : \Spec(k) \to Y$.
\end{lemma}

\begin{proof}
C'est une conséquence immédiate du lemme \ref{lemma-base-change-shriek} et des
définitions.
\end{proof}

\begin{lemma}
\label{lemma-compact-support-h0}
Soit $X$ un schéma séparé de type fini sur un corps $k$.
Si $\mathcal{F}$ est un faisceau abélien de torsion, alors le groupe abélien
$H^0_c(X, \mathcal{F})$ défini dans la définition \ref{definition-compact-support}
coïncide avec le groupe abélien $H^0_c(X, \mathcal{F})$ défini dans la
définition \ref{definition-cohomology-compact-support}.
\end{lemma}

\begin{proof}
Choisissons une compactification $j : X \to \overline{X}$ sur $k$.
Dans les deux cas, le groupe est défini comme $H^0(\overline{X}, j_!\mathcal{F})$.
Pour la première version, cela résulte du
lemme \ref{lemma-compactify-compact-support} ;
pour la seconde, cela résulte de la construction.
\end{proof}

\begin{lemma}
\label{lemma-finiteness-curves}
Soit $k$ un corps algébriquement clos. Soit $X$ un schéma séparé
de type fini sur $k$, de dimension $\leq 1$.
Soit $\Lambda$ un anneau noethérien.
Soit $\mathcal{F}$ un faisceau constructible de $\Lambda$-modules
sur $X$, supposé de torsion. Alors $H^q_c(X, \mathcal{F})$ est un
$\Lambda$-module de type fini.
\end{lemma}

\begin{proof}
Cela résulte du théorème de Cohomologie étale
\ref{etale-cohomology-theorem-vanishing-affine-curves-coefficients}.
Choisissons en effet une compactification $j : X \to \overline{X}$.
Après avoir remplacé $\overline{X}$ par l'adhérence schématique
de $X$, on voit que l'on peut supposer $\dim(\overline{X}) \leq 1$.
Alors $H^q_c(X, \mathcal{F}) = H^q(\overline{X}, j_!\mathcal{F})$
et le théorème s'applique.
\end{proof}

\begin{remark}[Covariance de la cohomologie à support propre]
\label{remark-covariance-compactly-supported}
Soit $k$ un corps. Soit $f : X \to Y$ un morphisme de schémas séparés
de type fini sur $k$. Si $X$, $Y$ et $f$ vérifient l'une
des conditions suivantes :
\begin{enumerate}
\item $f$ est étale ;
\item $f$ est plat et quasi-fini ;
\item $f$ est quasi-fini et $Y$ est géométriquement unibranche ;
\item $f$ est quasi-fini et il existe une pondération
$w : X \to \mathbf{Z}$ de $f$,
\end{enumerate}
alors la cohomologie à support propre est covariante vis-à-vis de $f$.
Plus précisément, soit $\Lambda$ un anneau. Soit $K$ un objet de
$D^+_{tors}(Y_\etale, \Lambda)$ ou de $D(Y_\etale, \Lambda)$ lorsque
$\Lambda$ est un anneau de torsion. Sous l'une des hypothèses (1) -- (4),
il existe un morphisme canonique
$$
\text{Tr}_{f, w, K} : f_!f^{-1}K \longrightarrow K
$$
Voir la section \ref{section-weightings} pour l'existence du
morphisme trace, et les exemples \ref{example-trace-for-flat-quasi-finite} et
\ref{example-trace-for-quasi-finite-over-normal} pour les cas (2) et (3).
Si $p : X \to \Spec(k)$ et $q : Y \to \Spec(k)$ désignent
les morphismes structuraux, alors $Rq_! \circ f_! = Rp_!$
d'après le lemme \ref{lemma-shriek-composition} et l'égalité
$Rf_! = f_!$ pour le morphisme quasi-fini séparé $f$, donnée par le
lemme \ref{lemma-Rf-shriek-for-quasi-finite}. On peut donc considérer
le morphisme
\begin{align*}
R\Gamma_c(X, f^{-1}K)
& =
R\Gamma(\Spec(k), Rp_!f^{-1}K) \\
& =
R\Gamma(\Spec(k), Rq_!f_!f^{-1}K) \\
& \xrightarrow{Rq_!\text{Tr}_{f, w, K}}
R\Gamma(\Spec(k), Rq_!K) \\
& =
R\Gamma_c(Y, K)
\end{align*}
En particulier, si $\Lambda$ est un anneau de torsion, on obtient une flèche
$$
\text{Tr}_f : R\Gamma_c(X, \Lambda) \longrightarrow R\Gamma_c(Y, \Lambda)
$$
Ce morphisme vérifie de nombreuses propriétés supplémentaires ; il est notamment
compatible avec tout changement de corps de base.
\end{remark}










\section{Un résultat de constructibilité}
\label{section-family-smooth-curves-cohomology}

\noindent
Nous « calculons » la cohomologie d'une famille projective lisse
de courbes à coefficients constants.

\begin{lemma}
\label{lemma-frobenius-linear-on-vb}
Soit $p$ un nombre premier. Soit $S$ un schéma sur $\mathbf{F}_p$.
Soit $\mathcal{E}$ un module localement libre de type fini
sur $\mathcal{O}_S$, vu comme $\mathcal{O}_S$-module sur $S_\etale$.
Soit $F : \mathcal{E} \to \mathcal{E}$ un homomorphisme de faisceaux abéliens
sur $S_\etale$ tel que $F(a e) = a^pF(e)$ pour toutes sections locales $a$, $e$
respectivement de $\mathcal{O}_S$ et $\mathcal{E}$ sur $S_\etale$. Alors
$$
\Coker(F - 1 : \mathcal{E} \to \mathcal{E})
$$
est nul, et
$$
\Ker(F - 1 : \mathcal{E} \to \mathcal{E})
$$
est un faisceau abélien constructible sur $S_\etale$.
\end{lemma}

\noindent
Ce lemme généralise le lemme de Cohomologie étale
\ref{etale-cohomology-lemma-F-1}.

\begin{proof}
Nous pouvons supposer $S = \Spec(A)$, où $A$ est une $\mathbf{F}_p$-algèbre,
et que $\mathcal{E}$ est le module quasi-cohérent associé
au $A$-module libre $Ae_1 \oplus \ldots \oplus Ae_n$.
Écrivons $F(e_i) = \sum a_{ij} e_j$.

\medskip\noindent
Surjectivité de $F - 1$. Il suffit de montrer que tout élément
$\sum a_i e_i$, $a_i \in A$, appartient à l'image de $F - 1$
après remplacement de $A$ par une extension étale fidèlement plate.
Remarquons que
$$
F(\sum x_ie_i) - \sum x_i e_i = \sum x_i^p a_{ij} e_j - \sum x_i e_i
$$
Considérons la $A$-algèbre
$$
A' = A[x_1, \ldots, x_n]/(a_i + x_i - \sum\nolimits_j a_{ji} x_j^p)
$$
Un calcul montre que $\text{d}x_i$ est nul dans $\Omega_{A'/A}$,
donc que $\Omega_{A'/A} = 0$. Comme $A'$ est de type fini sur
$A$, il en résulte que $\Spec(A') \to \Spec(A)$ est non ramifié,
donc quasi-fini. Comme $A'$ est engendré par $n$ éléments
et défini par $n$ équations, on conclut que $A'$
est une intersection complète relative globale sur $A$.
Ainsi, $A'$ est plat sur $A$, et $A \to A'$
est étale (comme morphisme d'anneaux plat et non ramifié). Enfin,
le lecteur peut montrer que $A \to A'$ est fidèlement plat
en vérifiant directement que toutes les fibres géométriques de
$\Spec(A') \to \Spec(A)$ sont non vides ; cela résulte aussi du lemme de
Cohomologie étale
\ref{etale-cohomology-lemma-F-1}.
Enfin, l'élément $\sum x_i e_i \in A'e_1 \oplus \ldots \oplus A'e_n$
a pour image $\sum a_i e_i$ par $F - 1$.

\medskip\noindent
Constructibilité du noyau. Les calculs ci-dessus montrent que
$\Ker(F - 1)$ est représenté par le schéma
$$
\Spec(A[x_1, \ldots, x_n]/(x_i - \sum\nolimits_j a_{ji} x_j^p))
$$
sur $S = \Spec(A)$. Comme il s'agit d'un schéma affine et étale
sur $S$, le résultat découle du
lemme \ref{etale-cohomology-lemma-jshriek-constructible} de Cohomologie étale.
\end{proof}

\begin{lemma}
\label{lemma-proper-smooth-family-curves}
Soit $f : X \to S$ un morphisme propre et lisse de schémas, à
fibres géométriquement connexes de dimension $1$. Soit $\ell$
un nombre premier. Alors $R^qf_*\underline{\mathbf{Z}/\ell\mathbf{Z}}$
est constructible.
\end{lemma}

\begin{proof}
Nous pouvons supposer $S$ affine, disons $S = \Spec(A)$. Si l'on écrit
$A = \bigcup A_i$ comme réunion de ses sous-$\mathbf{Z}$-algèbres de type fini,
on peut trouver un $i$ et un morphisme $f_i : X_i \to S_i = \Spec(A_i)$ de type
fini dont le changement de base à $S$ est $f : X \to S$ ; voir le
lemme \ref{limits-lemma-descend-finite-presentation} de Limites.
Quitte à augmenter $i$, on peut supposer $f_i : X_i \to S_i$
lisse, propre et de dimension relative $1$ ; voir les
lemmes de Limites
\ref{limits-lemma-eventually-proper},
\ref{limits-lemma-descend-smooth}, et
\ref{limits-lemma-descend-dimension-d}.
D'après le lemme \ref{more-morphisms-lemma-proper-flat-geom-red} de Compléments sur les morphismes,
on obtient un sous-schéma ouvert $U_i \subset S_i$ tel que les fibres
de $f_i : X_i \to S_i$ au-dessus de $U_i$ soient géométriquement connexes.
Alors $S \to S_i$ se factorise par $U_i$. On peut remplacer $X \to S$
par $f_i : f_i^{-1}(U_i) \to U_i$ pour se ramener au cas
traité dans le paragraphe suivant.

\medskip\noindent
Supposons $S$ noethérien. On peut écrire $S = U \cup Z$, où $U$ est le
sous-schéma ouvert défini par la non-annulation de $\ell$ et
$Z = V(\ell) \subset S$. Comme la formation de
$R^qf_*\underline{\mathbf{Z}/\ell\mathbf{Z}}$ commute à tout
changement de base (théorème de Cohomologie étale
\ref{etale-cohomology-theorem-proper-base-change}),
il suffit de démontrer le résultat sur $U$ et sur $Z$.
On se ramène ainsi aux deux cas suivants : (a) $\ell$
est inversible sur $S$ ; (b) $\ell$ est nul sur $S$.

\medskip\noindent
Cas (a). Nous affirmons que, dans ce cas, les faisceaux
$R^qf_*\underline{\mathbf{Z}/\ell\mathbf{Z}}$ sont
localement constants à fibres finies sur $S$. D'abord, par changement de base propre
(sous la forme du lemme de Cohomologie étale
\ref{etale-cohomology-lemma-proper-base-change-stalk})
et par finitude
(théorème de Cohomologie étale
\ref{etale-cohomology-theorem-vanishing-affine-curves})
on voit que les fibres de $R^qf_*\underline{\mathbf{Z}/\ell\mathbf{Z}}$
sont finies. D'après le lemme de Cohomologie étale
\ref{etale-cohomology-lemma-sp-isom-proper-loc-cst-torsion}
tous les morphismes de spécialisation sont des isomorphismes.
On conclut grâce au
lemme de Cohomologie étale
\ref{etale-cohomology-lemma-characterize-locally-constant-module}.

\medskip\noindent
Cas (b). Ici, $\ell = p$ est premier et $S$ est un schéma sur
$\Spec(\mathbf{F}_p)$. Les mêmes références que ci-dessus montrent déjà
que les fibres de $R^qf_*\underline{\mathbf{Z}/p\mathbf{Z}}$
sont finies et nulles pour $q \geq 2$. Il résulte du
lemme \ref{etale-cohomology-lemma-sections-upstairs} de Cohomologie étale
que $f_*\underline{\mathbf{Z}/p\mathbf{Z}} =
\underline{\mathbf{Z}/p\mathbf{Z}}$. Il reste à démontrer
que $R^1f_*\underline{\mathbf{Z}/p\mathbf{Z}}$ est constructible.
Considérons la suite d'Artin--Schreier
$$
0 \to \underline{\mathbf{Z}/p\mathbf{Z}} \to \mathcal{O}_X
\xrightarrow{F - 1} \mathcal{O}_X
\to 0
$$
Voir Cohomologie étale, section \ref{etale-cohomology-section-artin-schreier}.
Rappelons que $f_*\mathcal{O}_X = \mathcal{O}_S$ et que
$R^1f_*\mathcal{O}_X$ est un $\mathcal{O}_S$-module localement libre de type fini,
de rang égal au genre des fibres de $X \to S$ ; voir le
lemme \ref{curves-lemma-genus-in-nodal-family-of-curves} de Courbes algébriques.
On en déduit une suite exacte courte
$$
0 \to
\Coker(F - 1 : \mathcal{O}_S \to \mathcal{O}_S)
\to R^1f_*\underline{\mathbf{Z}/p\mathbf{Z}} \to
\Ker(F - 1 : R^1f_*\mathcal{O}_X \to R^1f_*\mathcal{O}_X)
\to 0
$$
Le lemme \ref{lemma-frobenius-linear-on-vb} permet de conclure.
\end{proof}

\begin{lemma}
\label{lemma-proper-smooth-family-curves-modules}
Soit $f : X \to S$ un morphisme propre et lisse de schémas, à
fibres géométriquement connexes de dimension $1$. Soit $\Lambda$
un anneau noethérien. Soit $M$ un $\Lambda$-module de type fini
annulé par un entier $n > 0$. Alors $R^qf_*\underline{M}$
est un faisceau constructible de $\Lambda$-modules sur $S$.
\end{lemma}

\begin{proof}
Si $n = \ell n'$ pour un nombre premier $\ell$, on obtient
une suite exacte courte $0 \to M[\ell] \to M \to M' \to 0$
de $\Lambda$-modules de type fini, et $M'$ est annulé par $n'$.
Elle donne une suite exacte courte correspondante de faisceaux
constants, qui donne à son tour une suite exacte
$$
R^{q - 1}f_*\underline{M'} \to
R^qf_*\underline{M[n]} \to
R^qf_*\underline{M} \to
R^qf_*\underline{M'} \to
R^{q + 1}f_*\underline{M[n]}
$$
Ainsi, si l'on démontre le résultat lorsque $M$ est annulé par
un nombre premier, une récurrence sur $n$ permet de conclure grâce au
lemme \ref{etale-cohomology-lemma-constructible-abelian} de Cohomologie étale.

\medskip\noindent
Soit $\ell$ un nombre premier tel que $\ell$ annule $M$.
On peut alors remplacer $\Lambda$ par la $\mathbf{F}_\ell$-algèbre
$\Lambda/\ell \Lambda$. En effet, le faisceau $R^qf_*\underline{M}$,
où $\underline{M}$ est vu comme faisceau de $\Lambda$-modules,
s'identifie au faisceau $R^qf_*\underline{M}$ calculé en
considérant $\underline{M}$ comme faisceau de $\Lambda/\ell \Lambda$-modules ; voir le
lemme de Cohomologie sur les sites
\ref{sites-cohomology-lemma-modules-abelian-unbounded}.

\medskip\noindent
Supposons que $\ell$ soit un nombre premier tel que $\ell$ annule $M$ et $\Lambda$.
Ramenons-nous au cas où $M$ est un $\Lambda$-module libre de type fini.
Choisissons pour cela une résolution
$$
\ldots \to
\Lambda^{\oplus m_2} \to
\Lambda^{\oplus m_1} \to
\Lambda^{\oplus m_0} \to
M \to 0
$$
Rappelons que $f_*$ est de dimension cohomologique finie sur les faisceaux de
$\Lambda$-modules ; voir le
lemme de Cohomologie étale
\ref{etale-cohomology-lemma-cohomological-dimension-proper}
et le
lemme \ref{derived-lemma-unbounded-right-derived} de Catégories dérivées.
On voit ainsi que $R^qf_*\underline{M}$ est le $q$-ième faisceau
de cohomologie de l'objet
$$
Rf_*(\underline{\Lambda^{\oplus m_a}} \to \ldots \to
\underline{\Lambda^{\oplus m_0}})
$$
de $D(S_\etale, \Lambda)$ pour un entier $a$ assez grand.
En utilisant la première suite spectrale du
lemme \ref{derived-lemma-two-ss-complex-functor} de Catégories dérivées
(ou bien en utilisant un argument par troncations),
on conclut qu'il suffit de démontrer que $R^qf_*\underline{\Lambda})$
est constructible.

\medskip\noindent
Nous pouvons enfin utiliser l'égalité
$$
(Rf_*\underline{\mathbf{Z}/\ell\mathbf{Z}})
\otimes_{\mathbf{Z}/\ell\mathbf{Z}}^\mathbf{L} \underline{\Lambda} =
Rf_*\underline{\Lambda}
$$
du lemme de Cohomologie étale
\ref{etale-cohomology-lemma-projection-formula-proper-mod-n}.
Comme tout module sur le corps $\mathbf{Z}/\ell\mathbf{Z}$
est plat, on obtient
$$
(R^qf_*\underline{\mathbf{Z}/\ell\mathbf{Z}})
\otimes_{\mathbf{Z}/\ell\mathbf{Z}} \underline{\Lambda} =
R^qf_*\underline{\Lambda}
$$
Il suffit donc de démontrer le résultat pour
$R^qf_*\underline{\mathbf{Z}/\ell\mathbf{Z}}$
d'après le lemme \ref{etale-cohomology-lemma-tensor-constructible} de Cohomologie étale.
C'est précisément le lemme \ref{lemma-proper-smooth-family-curves}.
\end{proof}






\section{Complexes à cohomologie constructible}
\label{section-Dc}

\noindent
Nous poursuivons la discussion commencée dans Cohomologie étale, section
\ref{etale-cohomology-section-Dc}. En particulier, pour un schéma
$X$ et un anneau noethérien $\Lambda$, nous notons
$D_c(X_\etale, \Lambda)$ la sous-catégorie triangulée strictement pleine et saturée
de $D(X_\etale, \Lambda)$ formée des objets
dont les faisceaux de cohomologie sont des faisceaux constructibles de $\Lambda$-modules.

\begin{lemma}
\label{lemma-qf-f-shriek-constructible}
Soit $f : X \to Y$ un morphisme de schémas
localement quasi-fini et de présentation finie.
Le foncteur $f_! : D(X_\etale, \Lambda) \to D(Y_\etale, \Lambda)$
du lemme \ref{lemma-lqf-shriek-derived}
envoie $D_c(X_\etale, \Lambda)$ dans $D_c(Y_\etale, \Lambda)$.
\end{lemma}

\begin{proof}
Comme le foncteur $f_!$ est exact, il suffit de montrer que
$f_!\mathcal{F}$ est constructible pour tout faisceau constructible
$\mathcal{F}$ de $\Lambda$-modules sur $X_\etale$.
La question est locale sur $Y$ ; nous pouvons donc supposer
$Y$ affine, ce que nous faisons.
Alors $X$ est quasi-compact et quasi-séparé ; voir la
définition \ref{morphisms-definition-finite-presentation} de Morphismes.
Soit $X = \bigcup_{i = 1, \ldots, n} X_i$ un recouvrement ouvert affine
fini. D'après le lemme \ref{lemma-lqf-colimit-f-shriek},
il suffit de montrer que
$f_{i, !}\mathcal{F}|_{X_i}$ et $f_{ii', !}\mathcal{F}|_{X_i \cap X_{i'}}$
sont constructibles, où $f_i : X_i \to Y$ et
$f_{ii'} : X_i \cap X_{i'} \to Y$ sont les restrictions de $f$.
Comme $X_i$ et $X_i \cap X_{i'}$ sont quasi-compacts et séparés,
on peut donc supposer $f$ séparé.
D'après le théorème principal de Zariski (sous la forme du
lemme de Compléments sur les morphismes
\ref{more-morphisms-lemma-quasi-finite-separated-pass-through-finite-addendum})
on peut choisir une factorisation $f = g \circ j$, où
$j : X \to X'$ est une immersion ouverte et $g : X' \to Y$ est
fini et de présentation finie. Alors $f_! = g_! \circ j_!$ d'après le
lemme \ref{lemma-f-shriek-composition}.
D'après le lemme \ref{etale-cohomology-lemma-jshriek-constructible} de Cohomologie étale,
$j_!\mathcal{F}$ est constructible sur $X'$.
Le morphisme $g$ est fini, donc $g_! = g_*$
d'après le lemme \ref{lemma-proper-f-shriek}.
Ainsi, $f_!\mathcal{F} = g_!j_!\mathcal{F} = g_*j_!\mathcal{F}$
est constructible d'après le
lemme de Cohomologie étale
\ref{etale-cohomology-lemma-finite-pushforward-constructible}.
\end{proof}

\begin{lemma}
\label{lemma-constant-shriek-rel-dim-1}
Soit $S$ un schéma affine noethérien de dimension finie.
Soit $f : X \to S$ un morphisme séparé, affine et lisse
de dimension relative $1$. Soit $\Lambda$ un anneau noethérien
de torsion. Soit $M$ un $\Lambda$-module de type fini.
Alors $Rf_!\underline{M}$ a ses faisceaux de cohomologie constructibles.
\end{lemma}

\begin{proof}
Nous démontrerons le résultat par récurrence sur $d = \dim(S)$.

\medskip\noindent
Cas initial. Si $d = 0$, il suffit de montrer que les fibres
de $R^qf_!\underline{M}$ sont des $\Lambda$-modules de type fini.
Si $\overline{s}$ est un point géométrique de $S$, alors
$(R^qf_!\underline{M})_{\overline{s}} = H^q_c(X_{\overline{s}}, \underline{M})$
d'après le lemme \ref{lemma-stalk-R-f-shriek}.
C'est un $\Lambda$-module de type fini d'après le lemme \ref{lemma-finiteness-curves}.

\medskip\noindent
Étape de récurrence. Il suffit de trouver un ouvert dense $U \subset S$ tel
que $Rf_!\underline{M}|_U$ ait ses faisceaux de cohomologie constructibles.
En effet, la restriction de $Rf_!\underline{M}$ au complémentaire
$S \setminus U$ aura ses faisceaux de cohomologie constructibles par récurrence
et parce que la formation de $Rf_!\underline{M}$ commute
à tout changement de base (lemme \ref{lemma-base-change-shriek}).
Plus précisément, soit $\eta \in S$ un point générique d'une composante
irréductible de $S$. Il suffit alors de trouver un voisinage ouvert $U$
de $\eta$ tel que la restriction de $Rf_!\underline{M}$
à $U$ soit constructible. C'est l'objet du paragraphe suivant.

\medskip\noindent
Étant donné un point générique $\eta \in S$, choisissons un diagramme
$$
\xymatrix{
\overline{Y}_1 \amalg \ldots \amalg \overline{Y}_n \ar[rd] &
Y_1 \amalg \ldots \amalg Y_n \ar[r]_-\nu \ar[d] \ar[l]^j &
X_V \ar[r] \ar[d] &
X_U \ar[r] \ar[d] &
X \ar[d]^f \\
& T_1 \amalg \ldots \amalg T_n \ar[r] &
V \ar[r] &
U \ar[r] &
S
}
$$
comme dans le lemme \ref{more-morphisms-lemma-make-good-curves} de Compléments sur les morphismes.
Nous allons montrer que $Rf_!\underline{M}|_U$ est constructible.
D'abord, comme $V \to U$ est fini et surjectif, il suffit
de montrer que son image inverse sur $V$ est constructible ; voir le
lemme \ref{etale-cohomology-lemma-check-constructible} de Cohomologie étale.
Comme la formation de $Rf_!$ commute au changement de base,
il suffit de montrer que
$R(X_V \to V)_!\underline{M}$ est constructible.
Soit $W \subset X_V$ le sous-schéma ouvert fourni par la partie (4) du
lemme \ref{more-morphisms-lemma-make-good-curves} de Compléments sur les morphismes.
Soit $Z \subset X_V$ le sous-schéma induit réduit sur
le complémentaire de $W$ dans $X_V$. Alors les fibres de $Z \to V$
sont de dimension $0$ (car $W$ est dense dans les fibres), donc
$Z \to V$ est quasi-fini. Du triangle distingué
$$
R(W \to V)_!\underline{M} \to
R(X_V \to V)_!\underline{M} \to
R(Z \to V)_!\underline{M} \to
\ldots
$$
du lemme \ref{lemma-relative-triangle-associated-to-open} et
du lemme \ref{lemma-qf-f-shriek-constructible},
on déduit qu'il suffit de montrer que
$R(W \to V)_!\underline{M}$ a ses faisceaux de cohomologie constructibles.
Ensuite, on a
$$
R(W \to V)_!\underline{M} = R(\nu^{-1}(W) \to V)_!\underline{M}
$$
parce que le morphisme $\nu : \nu^{-1}(W) \to W$ est un épaississement
et que l'on peut appliquer le lemme \ref{lemma-shriek-and-thickening}.
Notons ensuite $Z' \subset \coprod \overline{Y}_i$ le
complémentaire de l'ouvert $j(\nu^{-1}(W))$. Là encore, $Z' \to V$ est quasi-fini.
Utilisons encore le triangle distingué
$$
R(\nu^{-1}(W) \to V)_!\underline{M} \to
R(\coprod \overline{Y}_i \to V)_!\underline{M} \to
R(Z' \to V)_!\underline{M} \to
\ldots
$$
pour conclure qu'il suffit de démontrer que
$$
R(\coprod \overline{Y}_i \to V)_!\underline{M} =
\bigoplus\nolimits_i R(\overline{Y}_i \to V)_!\underline{M} =
\bigoplus\nolimits_i R(T_i \to V)_!R(\overline{Y}_i \to T_i)_!\underline{M}
$$
a ses faisceaux de cohomologie constructibles (la seconde égalité résultant du
lemme \ref{lemma-shriek-composition}).
Le résultat pour $R(\overline{Y}_i \to T_i)_!\underline{M}$
résulte du lemme \ref{lemma-proper-smooth-family-curves-modules},
et l'on conclut parce que $T_i \to V$ est fini étale
et que l'on peut appliquer le lemme \ref{lemma-qf-f-shriek-constructible}.
\end{proof}

\begin{lemma}
\label{lemma-loc-constant-shriek-rel-dim-1}
Soit $Y$ un schéma affine noethérien de dimension finie.
Soit $\Lambda$ un anneau noethérien de torsion.
Soit $\mathcal{F}$ un faisceau de
$\Lambda$-modules localement constant et de type fini sur un sous-schéma ouvert $U \subset \mathbf{A}^1_Y$.
Alors $Rf_!\mathcal{F}$ a ses faisceaux de cohomologie constructibles, où
$f : U \to Y$ est le morphisme structural.
\end{lemma}

\begin{proof}
On peut décomposer $\Lambda$ en un produit
$\Lambda = \Lambda_1 \times \ldots \times \Lambda_r$
où $\Lambda_i$ est $\ell_i$-primaire pour un nombre premier $\ell_i$.
On peut donc supposer qu'il existe un nombre premier $\ell$ et un entier $n > 0$
tels que $\ell^n$ annule $\Lambda$ (et donc $\mathcal{F}$).

\medskip\noindent
Comme $U$ est noethérien, $U$ n'a qu'un nombre fini de composantes
connexes. On peut donc supposer $U$ connexe.
Soit $g : U' \to U$ le revêtement étale fini
construit dans le lemme de Cohomologie étale
\ref{etale-cohomology-lemma-pullback-filtered-modules}.
La discussion de la
section \ref{etale-cohomology-section-trace-method} de Cohomologie étale
donne des morphismes
$$
\mathcal{F} \to g_*g^{-1}\mathcal{F} \to \mathcal{F}
$$
dont la composée est un isomorphisme. Il suffit donc de
démontrer le résultat pour $g_*g^{-1}\mathcal{F}$. D'autre part,
on a $Rf_!g_*g^{-1}\mathcal{F} = R(f \circ g)_!g^{-1}\mathcal{F}$
d'après le lemme \ref{lemma-shriek-composition}.
Comme $g^{-1}\mathcal{F}$ admet une filtration finie dont les quotients sont des
faisceaux constants de $\Lambda$-modules de la forme
$\underline{M}$ pour un $\Lambda$-module de type fini $M$
(par notre choix de $g$), on se ramène au cas démontré au
lemme \ref{lemma-constant-shriek-rel-dim-1}.
\end{proof}

\begin{lemma}
\label{lemma-constructible-shriek-rel-dim-1}
Soit $Y$ un schéma affine. Soit $\Lambda$ un anneau noethérien.
Soit $\mathcal{F}$ un faisceau constructible de $\Lambda$-modules de torsion sur
$\mathbf{A}^1_Y$. Alors $Rf_!\mathcal{F}$ a ses faisceaux de cohomologie
constructibles, où $f : \mathbf{A}^1_Y \to Y$ est le morphisme structural.
\end{lemma}

\begin{proof}
Supposons $\mathcal{F}$ annulé par $n > 0$. On peut alors remplacer
$\Lambda$ par $\Lambda/n\Lambda$ sans changer $Rf_!\mathcal{F}$.
On peut donc supposer, et supposons, que $\Lambda$ est un anneau de torsion.

\medskip\noindent
Écrivons $Y = \Spec(R)$. Si l'on écrit
$R = \bigcup R_i$ comme réunion de ses sous-$\mathbf{Z}$-algèbres de type fini,
on peut trouver un $i$ tel que $\mathcal{F}$ provienne par image inverse
d'un faisceau constructible de $\Lambda$-modules sur $\mathbf{A}^1_{R_i}$ ; voir le
lemme \ref{etale-cohomology-lemma-category-is-colimit} de Cohomologie étale.
On peut donc supposer $Y$ noethérien et de dimension finie.

\medskip\noindent
Supposons $Y$ noethérien de dimension finie $d = \dim(Y)$
et $\Lambda$ un anneau de torsion. Nous démontrerons le résultat par récurrence sur $d$.

\medskip\noindent
Cas initial. Si $d = 0$, il suffit de montrer que les fibres
de $R^qf_!\mathcal{F}$ sont des $\Lambda$-modules de type fini.
Si $\overline{y}$ est un point géométrique de $Y$, alors
$(R^qf_!\mathcal{F})_{\overline{y}} = H^q_c(X_{\overline{y}}, \mathcal{F})$
d'après le lemme \ref{lemma-stalk-R-f-shriek}.
C'est un $\Lambda$-module de type fini d'après le lemme \ref{lemma-finiteness-curves}.

\medskip\noindent
Étape de récurrence. Il suffit de trouver un ouvert dense $V \subset Y$ tel
que $Rf_!\mathcal{F}|_V$ ait ses faisceaux de cohomologie constructibles. En effet, la
restriction de $Rf_!\mathcal{F}$ au complémentaire $Y \setminus V$
aura ses faisceaux de cohomologie constructibles par récurrence
et parce que la formation de $Rf_!\mathcal{F}$ commute
à tout changement de base (lemme \ref{lemma-base-change-shriek}).
Par définition des faisceaux constructibles de $\Lambda$-modules, il existe un
sous-schéma ouvert dense $U \subset \mathbf{A}^1_Y$ tel que $\mathcal{F}|_U$
soit un faisceau de $\Lambda$-modules localement constant et de type fini.
Notons $Z \subset \mathbf{A}^1_Y$ le complémentaire (muni de sa structure de
sous-schéma fermé réduit). Remarquons que $U$ contient tous les points génériques des
fibres de $\mathbf{A}^1_Y \to Y$ au-dessus des points génériques
$\xi_1, \ldots, \xi_n$ des composantes irréductibles de $Y$.
Ainsi, $Z \to Y$ a des fibres finies au-dessus de $\xi_1, \ldots, \xi_n$.
Quitte à remplacer $Y$ par un ouvert dense (ce qui est permis), on peut
supposer $Z \to Y$ fini ; voir le
lemme \ref{morphisms-lemma-generically-finite} de Morphismes.
Le triangle distingué du
lemme \ref{lemma-relative-triangle-associated-to-open}
et le résultat pour $Z \to Y$ (lemme \ref{lemma-qf-f-shriek-constructible})
nous ramènent à montrer que
$R(U \to Y)_!\mathcal{F}$ a ses faisceaux de cohomologie constructibles.
C'est précisément le lemme \ref{lemma-loc-constant-shriek-rel-dim-1}.
\end{proof}

\begin{theorem}
\label{theorem-constructible-shriek}
Soit $f : X \to Y$ un morphisme séparé de présentation finie entre
schémas quasi-compacts et quasi-séparés. Soit $\Lambda$ un anneau
noethérien. Soit $K$ un objet de $D^+_{tors, c}(X_\etale, \Lambda)$
ou de $D_c(X_\etale, \Lambda)$ lorsque $\Lambda$ est un anneau de torsion.
Alors $Rf_!K$ a ses faisceaux de cohomologie constructibles, c'est-à-dire que
$Rf_!K$ appartient à $D^+_{tors, c}(Y_\etale, \Lambda)$
ou à $D_c(Y_\etale, \Lambda)$ lorsque $\Lambda$ est un anneau de torsion.
\end{theorem}

\begin{proof}
La question est locale sur $Y$ ; nous pouvons donc supposer $Y$ affine.
Par le principe de récurrence et le lemme \ref{lemma-relative-mayer-vietoris},
on se ramène au cas où $X$ est aussi affine.

\medskip\noindent
Supposons $X$ et $Y$ affines. Comme $X$ est de présentation finie,
on peut choisir une immersion fermée $i : X \to \mathbf{A}^n_Y$
de présentation finie. Si $p : \mathbf{A}^n_Y \to Y$ désigne
le morphisme structural, on a $Rf_! = Rp_! \circ Ri_!$
d'après le lemme \ref{lemma-shriek-composition}. D'après le
lemme \ref{lemma-qf-f-shriek-constructible},
le résultat vaut pour $Ri_! = i_!$. On peut donc supposer que $f$
est le morphisme de projection $\mathbf{A}^n_Y \to Y$.
Comme on peut écrire $f$ comme la composée
$$
X = \mathbf{A}^n_Y \to \mathbf{A}^{n - 1}_Y \to
\mathbf{A}^{n - 2}_S \to \ldots \to
\mathbf{A}^1_Y \to Y
$$
on peut supposer $n = 1$.

\medskip\noindent
Supposons $Y$ affine et $X = \mathbf{A}^1_Y$. Comme $Rf_!$ est de
dimension cohomologique finie
(lemme \ref{lemma-derived-lower-shriek-bounded}),
on peut supposer $K$ borné inférieurement. En utilisant la première suite spectrale du
lemme \ref{derived-lemma-two-ss-complex-functor} de Catégories dérivées
(ou bien en utilisant un argument par troncations),
on se ramène au résultat du
lemme \ref{lemma-constructible-shriek-rel-dim-1}.
\end{proof}








\section{Applications}
\label{section-applications}

\noindent
Dans cette section, nous donnons quelques applications du
théorème \ref{theorem-constructible-shriek}.

\begin{lemma}
\label{lemma-finiteness-compactly-supported}
Soit $k$ un corps algébriquement clos.
Soit $X$ un schéma séparé de type fini sur $k$.
Soit $\Lambda$ un anneau noethérien. Soit $K$ un objet de
$D^+_{tors, c}(X_\etale, \Lambda)$ ou de $D_c(X_\etale, \Lambda)$
lorsque $\Lambda$ est un anneau de torsion. Alors
$H^i_c(X, K)$ est un $\Lambda$-module de type fini pour tout
$i \in \mathbf{Z}$.
\end{lemma}

\begin{proof}
C'est une conséquence immédiate du théorème \ref{theorem-constructible-shriek}
et de la définition de la cohomologie à support propre donnée à la
section \ref{section-compactly-supported-cohomology}.
\end{proof}

\begin{proposition}
\label{proposition-loc-cst-torsion}
Soit $f : X \to S$ un morphisme propre et lisse entre schémas.
Soit $\Lambda$ un anneau noethérien. Soit $\mathcal{F}$ un
faisceau de $\Lambda$-modules localement constant et de type fini
sur $X_\etale$ tel que, pour tout point géométrique $\overline{x}$ de $X$,
la fibre $\mathcal{F}_{\overline{x}}$ soit annulée par un entier
$n > 0$ premier à la caractéristique résiduelle de $\overline{x}$.
Alors $R^if_*\mathcal{F}$ est un faisceau de
$\Lambda$-modules localement constant et de type fini sur $S_\etale$ pour tout $i \in \mathbf{Z}$.
\end{proposition}

\begin{proof}
La question est locale sur $S$ ; on peut donc supposer $S$ affine.
Pour un point $x$ de $X$, notons $n_x \geq 1$
le plus petit entier annulant $\mathcal{F}_{\overline{x}}$
pour un point géométrique $\overline{x}$ de $X$
situé au-dessus de $x$ (et donc pour tout tel point). Comme $X$ est quasi-compact (étant propre sur un affine),
il existe un recouvrement étale fini $\{U_j \to X\}_{j = 1, \ldots, m}$
tel que $\mathcal{F}|_{U_j}$ soit constant. Comme $U_j \to X$
est ouvert, on en déduit que la fonction $x \mapsto n_x$ est localement
constante et ne prend qu'un nombre fini de valeurs. On obtient donc
une décomposition finie $X = X_1 \amalg \ldots \amalg X_N$
en sous-schémas ouverts et fermés telle que $n_x = n$ si et seulement si $x \in X_n$.
Il suffit alors de démontrer la proposition pour les morphismes induits
$X_n \to S$ et la restriction de $\mathcal{F}$ à $X_n$.
On peut donc supposer qu'il existe un entier $n > 0$ tel
que $\mathcal{F}$ soit annulé par $n$ et que $n$
soit premier aux caractéristiques résiduelles de tous les corps résiduels de $X$.

\medskip\noindent
Comme $f$ est lisse et propre, l'image $f(X) \subset S$ est ouverte et fermée.
On peut donc remplacer $S$ par $f(X)$ et supposer $f(X) = S$. En particulier,
on peut supposer $n$ inversible dans l'anneau de définition du
schéma affine $S$.

\medskip\noindent
Dans ce paragraphe, nous nous ramenons au cas où $S$ est noethérien.
Écrivons $S = \Spec(A)$ pour une $\mathbf{Z}[1/n]$-algèbre $A$.
Écrivons $A = \bigcup A_i$ comme réunion de ses sous-algèbres de type fini
sur $\mathbf{Z}[1/n]$. On peut trouver un $i$ et un morphisme
$f_i : X_i \to S_i = \Spec(A_i)$ de type
fini dont le changement de base à $S$ est $f : X \to S$ ; voir le
lemme \ref{limits-lemma-descend-finite-presentation} de Limites.
Quitte à augmenter $i$, on peut supposer $f_i : X_i \to S_i$
lisse et propre ; voir les lemmes de Limites
\ref{limits-lemma-eventually-proper},
\ref{limits-lemma-descend-smooth} et
\ref{limits-lemma-descend-dimension-d}.
D'après le lemme de Cohomologie étale
\ref{etale-cohomology-lemma-category-loc-cst-is-colimit}
il existe un $i$ et un faisceau
de $\Lambda$-modules localement constant et de type fini
$\mathcal{F}_i$ dont l'image inverse sur $X$ est isomorphe à $\mathcal{F}$.
Comme $\mathcal{F}$ est annulé par $n$, on peut remplacer
$\mathcal{F}_i$ par $\Ker(n : \mathcal{F}_i \to \mathcal{F}_i)$
et supposer qu'il en va de même pour $\mathcal{F}_i$.
On se ramène ainsi au cas traité dans le paragraphe suivant.

\medskip\noindent
Supposons donnés un entier $n \geq 1$, un schéma de base noethérien $S$
sur $\mathbf{Z}[1/n]$, et un faisceau $\mathcal{F}$ de $n$-torsion.
D'après le théorème \ref{theorem-constructible-shriek},
les faisceaux $R^if_*\mathcal{F}$ sont des faisceaux constructibles de $\Lambda$-modules.
D'après le lemme de Cohomologie étale
\ref{etale-cohomology-lemma-sp-isom-proper-torsion-loc-ac}
les morphismes de spécialisation de $R^if_*\mathcal{F}$ sont toujours des isomorphismes.
On conclut par le lemme de Cohomologie étale
\ref{etale-cohomology-lemma-characterize-locally-constant-module}.
\end{proof}







\section{Compléments sur l'image inverse extraordinaire dérivée}
\label{section-more-derived-upper-shriek}

\noindent
Soit $\Lambda$ un anneau de torsion. Considérons un diagramme commutatif
$$
\xymatrix{
U \ar[rr]_j \ar[rd]_g & & U' \ar[ld]^{g'} \\
& Y
}
$$
de schémas quasi-compacts et quasi-séparés, où $g$ et $g'$
sont séparés et de type fini, et où $j$ est étale. Il induit
un morphisme canonique
$$
Rg_!\Lambda \longrightarrow Rg'_!\Lambda
$$
dans $D(Y_\etale, \Lambda)$. En effet, d'après les lemmes \ref{lemma-shriek-composition}
et \ref{lemma-Rf-shriek-for-quasi-finite}, on a $Rg_! = Rg'_! \circ j_!$.
D'autre part, comme $j_!$ est adjoint à gauche de $j^{-1}$, on dispose du
morphisme d'adjonction $\text{Tr}_j : j_!\Lambda = j_!j^{-1}\Lambda \to \Lambda$ ;
nous l'appelons aussi le morphisme trace de $j$ ; voir la
remarque \ref{remark-trace-etale-counit}.
Le morphisme ci-dessus est la composée
$$
Rg_!\Lambda = Rg'_!j_!\Lambda
\xrightarrow{Rg'_! \text{Tr}_j}
Rg'_!\Lambda
$$
Étant donné un second morphisme étale $j' : U' \to U''$
et un morphisme séparé de type fini $g'' : U'' \to Y$,
la composée
$$
Rg_!\Lambda \longrightarrow Rg'_!\Lambda \longrightarrow Rg''_!\Lambda
$$
des morphismes associés à $j$ et $j'$ est égale au morphisme
$Rg_!\Lambda \longrightarrow Rg''_!\Lambda$ construit pour $j' \circ j$.
Cela résulte de l'énoncé correspondant pour les morphismes trace ;
voir le lemme \ref{lemma-trace-composition} pour un cas plus général.

\medskip\noindent
Soit $f : X \to Y$ un morphisme séparé de type fini entre schémas quasi-compacts
et quasi-séparés. On obtient alors un foncteur
$$
X_{affine, \etale}
\longrightarrow
\left\{
\begin{matrix}
\text{schémas séparés de type fini sur }Y\\
\text{et morphismes étales entre eux}
\end{matrix}
\right\}
$$
La construction ci-dessus détermine donc un foncteur
$X_{affine, \etale}^{opp} \to D(Y_\etale, \Lambda)$
qui envoie $U$ sur $R(U \to Y)_!\Lambda$.

\begin{lemma}
\label{lemma-describe-Rf-upper-shriek}
Soit $f : X \to Y$ un morphisme séparé de type fini entre schémas quasi-compacts
et quasi-séparés. Soit $\Lambda$ un anneau de torsion.
Soit $K \in D(Y_\etale, \Lambda)$. Pour $n \in \mathbf{Z}$, le
faisceau de cohomologie $H^n(Rf^!K)$ restreint à $X_{affine, \etale}$
est le faisceau associé au préfaisceau
$$
U \longmapsto \Hom_Y(R(U \to Y)_!\Lambda, K[n])
$$
Voir la discussion ci-dessus pour la fonctorialité de $R(U \to Y)_!\Lambda$.
\end{lemma}

\begin{proof}
Soit $j : U \to X$ un objet de $X_{affine, \etale}$, et posons $g = f \circ j$.
Rappelons que $\Hom_X(j_!\Lambda, M[n]) = H^n(U, M)$ pour tout objet $M$ de
$D(X_\etale, \Lambda)$. Alors $H^n(Rf^!K)$ est le faisceau associé au
préfaisceau
$$
U \mapsto H^n(U, Rf^!K) = \Hom_X(j_!\Lambda, Rf^!K[n]) =
\Hom_Y(Rf_!j_!\Lambda, K[n] = \Hom_Y(Rg_!\Lambda, K[n])
$$
Nous omettons de vérifier que les morphismes de transition sont induits par les morphismes
de transition entre les objets $Rg_!\Lambda = R(U \to Y)_!\Lambda$ construits
ci-dessus.
\end{proof}
















\begin{multicols}{2}[\section{Autres chapitres}]
\noindent
Préliminaires
\begin{enumerate}
\item \hyperref[introduction-section-phantom]{Introduction}
\item \hyperref[conventions-section-phantom]{Conventions}
\item \hyperref[sets-section-phantom]{Théorie des ensembles}
\item \hyperref[categories-section-phantom]{Catégories}
\item \hyperref[topology-section-phantom]{Topologie}
\item \hyperref[sheaves-section-phantom]{Faisceaux sur les espaces}
\item \hyperref[sites-section-phantom]{Sites et faisceaux}
\item \hyperref[stacks-section-phantom]{Champs}
\item \hyperref[fields-section-phantom]{Corps}
\item \hyperref[algebra-section-phantom]{Algèbre commutative}
\item \hyperref[brauer-section-phantom]{Groupes de Brauer}
\item \hyperref[homology-section-phantom]{Algèbre homologique}
\item \hyperref[derived-section-phantom]{Catégories dérivées}
\item \hyperref[simplicial-section-phantom]{Méthodes simpliciales}
\item \hyperref[more-algebra-section-phantom]{Compléments d'algèbre}
\item \hyperref[smoothing-section-phantom]{Lissage des morphismes d'anneaux}
\item \hyperref[modules-section-phantom]{Faisceaux de modules}
\item \hyperref[sites-modules-section-phantom]{Modules sur les sites}
\item \hyperref[injectives-section-phantom]{Objets injectifs}
\item \hyperref[cohomology-section-phantom]{Cohomologie des faisceaux}
\item \hyperref[sites-cohomology-section-phantom]{Cohomologie sur les sites}
\item \hyperref[dga-section-phantom]{Algèbre différentielle graduée}
\item \hyperref[dpa-section-phantom]{Algèbre à puissances divisées}
\item \hyperref[sdga-section-phantom]{Faisceaux différentiels gradués}
\item \hyperref[hypercovering-section-phantom]{Hyperrecouvrements}
\end{enumerate}
Schémas
\begin{enumerate}
\setcounter{enumi}{25}
\item \hyperref[schemes-section-phantom]{Schémas}
\item \hyperref[constructions-section-phantom]{Constructions de schémas}
\item \hyperref[properties-section-phantom]{Propriétés des schémas}
\item \hyperref[morphisms-section-phantom]{Morphismes de schémas}
\item \hyperref[coherent-section-phantom]{Cohomologie des schémas}
\item \hyperref[divisors-section-phantom]{Diviseurs}
\item \hyperref[limits-section-phantom]{Limites de schémas}
\item \hyperref[varieties-section-phantom]{Variétés}
\item \hyperref[topologies-section-phantom]{Topologies sur les schémas}
\item \hyperref[descent-section-phantom]{Descente}
\item \hyperref[perfect-section-phantom]{Catégories dérivées de schémas}
\item \hyperref[more-morphisms-section-phantom]{Compléments sur les morphismes}
\item \hyperref[flat-section-phantom]{Compléments sur la platitude}
\item \hyperref[groupoids-section-phantom]{Schémas en groupoïdes}
\item \hyperref[more-groupoids-section-phantom]{Compléments sur les schémas en groupoïdes}
\item \hyperref[etale-section-phantom]{Morphismes étales de schémas}
\end{enumerate}
Thèmes de théorie des schémas
\begin{enumerate}
\setcounter{enumi}{41}
\item \hyperref[chow-section-phantom]{Homologie de Chow}
\item \hyperref[intersection-section-phantom]{Théorie de l'intersection}
\item \hyperref[pic-section-phantom]{Schémas de Picard des courbes}
\item \hyperref[weil-section-phantom]{Théories cohomologiques de Weil}
\item \hyperref[adequate-section-phantom]{Modules adéquats}
\item \hyperref[dualizing-section-phantom]{Complexes dualisants}
\item \hyperref[duality-section-phantom]{Dualité pour les schémas}
\item \hyperref[discriminant-section-phantom]{Discriminants et différentes}
\item \hyperref[derham-section-phantom]{Cohomologie de de Rham}
\item \hyperref[local-cohomology-section-phantom]{Cohomologie locale}
\item \hyperref[algebraization-section-phantom]{Géométrie algébrique et formelle}
\item \hyperref[curves-section-phantom]{Courbes algébriques}
\item \hyperref[resolve-section-phantom]{Résolution des surfaces}
\item \hyperref[models-section-phantom]{Réduction semi-stable}
\item \hyperref[functors-section-phantom]{Foncteurs et morphismes}
\item \hyperref[equiv-section-phantom]{Catégories dérivées de variétés}
\item \hyperref[pione-section-phantom]{Groupes fondamentaux de schémas}
\item \hyperref[etale-cohomology-section-phantom]{Cohomologie étale}
\item \hyperref[crystalline-section-phantom]{Cohomologie cristalline}
\item \hyperref[proetale-section-phantom]{Cohomologie pro-étale}
\item \hyperref[relative-cycles-section-phantom]{Cycles relatifs}
\item \hyperref[more-etale-section-phantom]{Compléments de cohomologie étale}
\item \hyperref[trace-section-phantom]{La formule des traces}
\end{enumerate}
Espaces algébriques
\begin{enumerate}
\setcounter{enumi}{64}
\item \hyperref[spaces-section-phantom]{Espaces algébriques}
\item \hyperref[spaces-properties-section-phantom]{Propriétés des espaces algébriques}
\item \hyperref[spaces-morphisms-section-phantom]{Morphismes d'espaces algébriques}
\item \hyperref[decent-spaces-section-phantom]{Espaces algébriques décents}
\item \hyperref[spaces-cohomology-section-phantom]{Cohomologie des espaces algébriques}
\item \hyperref[spaces-limits-section-phantom]{Limites d'espaces algébriques}
\item \hyperref[spaces-divisors-section-phantom]{Diviseurs sur les espaces algébriques}
\item \hyperref[spaces-over-fields-section-phantom]{Espaces algébriques sur des corps}
\item \hyperref[spaces-topologies-section-phantom]{Topologies sur les espaces algébriques}
\item \hyperref[spaces-descent-section-phantom]{Descente et espaces algébriques}
\item \hyperref[spaces-perfect-section-phantom]{Catégories dérivées d'espaces}
\item \hyperref[spaces-more-morphisms-section-phantom]{Compléments sur les morphismes d'espaces}
\item \hyperref[spaces-flat-section-phantom]{Platitude sur les espaces algébriques}
\item \hyperref[spaces-groupoids-section-phantom]{Groupoïdes dans les espaces algébriques}
\item \hyperref[spaces-more-groupoids-section-phantom]{Compléments sur les groupoïdes dans les espaces}
\item \hyperref[bootstrap-section-phantom]{Amorçage}
\item \hyperref[spaces-pushouts-section-phantom]{Sommes amalgamées d'espaces algébriques}
\end{enumerate}
Thèmes de géométrie
\begin{enumerate}
\setcounter{enumi}{81}
\item \hyperref[spaces-chow-section-phantom]{Groupes de Chow des espaces}
\item \hyperref[groupoids-quotients-section-phantom]{Quotients de groupoïdes}
\item \hyperref[spaces-more-cohomology-section-phantom]{Compléments sur la cohomologie des espaces}
\item \hyperref[spaces-simplicial-section-phantom]{Espaces simpliciaux}
\item \hyperref[spaces-duality-section-phantom]{Dualité pour les espaces}
\item \hyperref[formal-spaces-section-phantom]{Espaces algébriques formels}
\item \hyperref[restricted-section-phantom]{Algébrisation des espaces formels}
\item \hyperref[spaces-resolve-section-phantom]{Résolution des surfaces revisitée}
\end{enumerate}
Théorie des déformations
\begin{enumerate}
\setcounter{enumi}{89}
\item \hyperref[formal-defos-section-phantom]{Théorie formelle des déformations}
\item \hyperref[defos-section-phantom]{Théorie des déformations}
\item \hyperref[cotangent-section-phantom]{Le complexe cotangent}
\item \hyperref[examples-defos-section-phantom]{Problèmes de déformation}
\end{enumerate}
Champs algébriques
\begin{enumerate}
\setcounter{enumi}{93}
\item \hyperref[algebraic-section-phantom]{Champs algébriques}
\item \hyperref[examples-stacks-section-phantom]{Exemples de champs}
\item \hyperref[stacks-sheaves-section-phantom]{Faisceaux sur les champs algébriques}
\item \hyperref[criteria-section-phantom]{Critères de représentabilité}
\item \hyperref[artin-section-phantom]{Axiomes d'Artin}
\item \hyperref[quot-section-phantom]{Espaces de Quot et de Hilbert}
\item \hyperref[stacks-properties-section-phantom]{Propriétés des champs algébriques}
\item \hyperref[stacks-morphisms-section-phantom]{Morphismes de champs algébriques}
\item \hyperref[stacks-limits-section-phantom]{Limites de champs algébriques}
\item \hyperref[stacks-cohomology-section-phantom]{Cohomologie des champs algébriques}
\item \hyperref[stacks-perfect-section-phantom]{Catégories dérivées de champs}
\item \hyperref[stacks-introduction-section-phantom]{Introduction aux champs algébriques}
\item \hyperref[stacks-more-morphisms-section-phantom]{Compléments sur les morphismes de champs}
\item \hyperref[stacks-geometry-section-phantom]{La géométrie des champs}
\end{enumerate}
Thèmes de théorie des espaces de modules
\begin{enumerate}
\setcounter{enumi}{107}
\item \hyperref[moduli-section-phantom]{Champs de modules}
\item \hyperref[moduli-curves-section-phantom]{Espaces de modules de courbes}
\end{enumerate}
Divers
\begin{enumerate}
\setcounter{enumi}{109}
\item \hyperref[examples-section-phantom]{Exemples}
\item \hyperref[exercises-section-phantom]{Exercices}
\item \hyperref[guide-section-phantom]{Guide bibliographique}
\item \hyperref[desirables-section-phantom]{Desiderata}
\item \hyperref[coding-section-phantom]{Style de programmation}
\item \hyperref[obsolete-section-phantom]{Obsolète}
\item \hyperref[fdl-section-phantom]{Licence de documentation libre GNU}
\item \hyperref[index-section-phantom]{Index généré automatiquement}
\end{enumerate}
\end{multicols}


\bibliography{my}
\bibliographystyle{amsalpha}

\end{document}
